\documentclass[11pt]{article}

\usepackage[margin=1in]{geometry}
\usepackage{amsmath,amssymb,amsthm,mathtools,booktabs,enumitem,microtype,array,tabularx,longtable,float,needspace}
\usepackage[hidelinks]{hyperref}

\newtheorem{theorem}{Theorem}
\newtheorem{proposition}[theorem]{Proposition}
\newtheorem{lemma}[theorem]{Lemma}
\newtheorem{corollary}[theorem]{Corollary}
\newtheorem{definition}[theorem]{Definition}
\newtheorem{remark}[theorem]{Remark}
\newtheorem{assumption}[theorem]{Assumption}

\newcommand{\SO}{\mathrm{SO}(3)}
\newcommand{\so}{\mathfrak{so}(3)}
\newcommand{\OO}{\mathrm{O}}
\newcommand{\R}{\mathbb{R}}
\newcommand{\I}{I_3}
\newcommand{\T}{\mathsf{T}}
\newcommand{\tr}{\operatorname{tr}}
\newcommand{\Exp}{\operatorname{Exp}}
\newcommand{\Log}{\operatorname{Log}}
\newcommand{\Ad}{\operatorname{Ad}}
\newcommand{\argmin}{\operatorname*{arg\,min}}
\newcommand{\argmax}{\operatorname*{arg\,max}}
\newcommand{\lexargmin}{\operatorname*{lex\,arg\,min}}
\newcommand{\clip}{\operatorname{clip}}
\newcommand{\diag}{\operatorname{diag}}
\newcommand{\rank}{\operatorname{rank}}
\newcommand{\median}{\operatorname{median}}
\newcommand{\dist}{\operatorname{dist}}
\newcommand{\wrap}{\operatorname{wrap}}
\newcommand{\TB}{\mathrm{TB}}
\newcommand{\PE}{\mathrm{PE}}
\newcommand{\sgn}{\operatorname{sgn}}
\newcommand{\norm}[1]{\left\lVert #1 \right\rVert}
\newcommand{\abs}[1]{\left\lvert #1 \right\rvert}
\newcommand{\set}[1]{\left\{ #1 \right\}}
\newcommand{\given}{\,\middle|\,}
\newcolumntype{Y}{>{\raggedright\arraybackslash}X}
\newenvironment{sectionrecap}
{\par\medskip\begingroup\small\noindent\rule{\linewidth}{0.4pt}\par
\smallskip\noindent\textbf{Section summary, key equations, and notation.}\par\smallskip}
{\par\smallskip\noindent\rule{\linewidth}{0.4pt}\par\endgroup\medskip}

\title{Recovering Equivalence Classes of Attitude and Angular-Rate Conventions\\
A Five-Module Hypothesis Engine with Hierarchical Cohort Evidence Fusion}
\author{}
\date{Cohort-evidence productization draft --- August 2026}

\begin{document}
    \setlength{\emergencystretch}{2em}


    \maketitle

    \begin{abstract}
        Underdocumented attitude and angular-rate streams create a finite but highly coupled
        convention-recovery problem: quaternion order, Euler sequence, active/passive direction,
        axis wiring, units, body/world rate expression, time scale, and latency may all be
        unknown. The engineering baseline is exhaustive and simple---enumerate plausible
        decoders, normalize them to one semantic model, integrate every admissible rate model,
        and compare the predicted rotation increments with the decoded attitude increments.
        The difficulty is deciding which hypotheses are genuinely different, which can be
        screened before integration, when enough evidence has accumulated to eliminate a
        candidate, and what a ``winner'' means when exact observational symmetries remain.

        This paper organizes the solution as a five-module hypothesis engine. Module~1
        generates and canonically normalizes attitude, rate, timing, and nuisance hypotheses.
        Module~2 quotients exact observational symmetries and reports an excitation/observability
        diagnostic rather than mistaking rank deficiency for model failure. Module~3 supplies
        explicitly heuristic unit, wrapping, Euler-range, and branch cues. Module~4 supplies a
        hierarchical provenance prior derived from primary documentation, with exact API/schema
        information stronger than ecosystem, artifact-role, or broad-domain labels. Module~5
        performs Lie-group integration, robust multiscale scoring, evidence accumulation, and
        safe candidate elimination. The inferential output is an evidence-preserving set of
        \emph{quotient classes}; a single software decoder is chosen only afterward by documented
        metadata, a declared domain prior, or a deterministic canonical policy.

        The framework also distinguishes known time, unknown clock/rate scale, attitude-only,
        and rate-only regimes; proves exact direction--frame--sign and time--rate scale gauges;
        and gives conservative finite-sample and exact-computational elimination rules. The
        result is a workflow in which every reduction has a label: \emph{hard constraint},
        \emph{exact quotient}, \emph{heuristic priority}, \emph{trajectory evidence}, or
        \emph{policy choice}.


        The single-trajectory engine is also given a product boundary. It emits a versioned
        \emph{evidence packet} containing quotient-class-relative evidence, the local prior that
        was used, calibration and observability descriptors, hard metadata, provenance, and
        lineage. A separate cohort-fusion service can then borrow information from strong
        reference packets to classify a weak target packet without double-counting priors or
        silently multiplying dependent evidence. Three transfer models are distinguished: one
        fixed convention shared by all members, an exchangeable partially pooled cohort, and a
        robust sharedness mixture that discounts borrowing when the target conflicts with the
        reference cohort. Numerical posterior probabilities are permitted only when the local
        evidence has been calibrated and the packet-dependence assumptions have been audited.

    \end{abstract}

    \tableofcontents


    \section{Introduction}\label{sec:introduction}

    The practical recovery problem is easy to state. Given one or more numeric orientation
    channels, angular-rate channels, and possibly timestamps, determine which interpretation
    makes the streams kinematically self-consistent. The natural baseline is therefore
    \begin{equation}
        \boxed{
            \text{enumerate decoders}
            \;\longrightarrow\;
            \text{normalize}
            \;\longrightarrow\;
            \text{integrate rates}
            \;\longrightarrow\;
            \text{compare rotations}.}
        \label{eq:baseline-idea}
    \end{equation}
    The implementation may use quaternions, matrices, or both, but the semantic comparison
    should live on the rotation group $\SO$. This prevents storage choices such as
    $q$ versus $-q$ from becoming false geometric differences.

    What makes the problem subtle is not the exhaustive loop itself. It is the existence of
    five logically different operations around that loop:
    \begin{enumerate}[nosep]
    \item generating and normalizing the candidate decoders;
    \item removing hypotheses that are exactly the same observable model and diagnosing
    whether the maneuver excites enough independent rotation directions;
    \item using representation-level clues such as wrapping, units, Euler ranges, and
    branch continuity without promoting those clues to proofs;
    \item using source, ecosystem, artifact-role, and domain knowledge as a prior rather
    than as trajectory evidence; and
    \item integrating, scoring, accumulating evidence, and eliminating candidates only
    when a deterministic or statistical rule justifies doing so.
    \end{enumerate}
    The previous organization of these ideas by mathematical topic makes the method look more
    complicated than it is. This draft instead makes those five operations the spine of the
    paper; detailed proofs, survey tables, timing cases, and benchmark machinery are moved to
    supporting appendices.

    A second organizing principle is that the algorithm never needs to pretend that every
    run has one uniquely identifiable convention. A concrete decoder $h$ is only one member
    of a larger structural class when an exact symmetry produces identical observations.
    The scientifically meaningful object is therefore often the quotient class $[h]$, not
    one arbitrary enum value. Better excitation can resolve many empirical near-ties, but it
    cannot break an exact gauge.

    Finally, this paper distinguishes \emph{elimination} from \emph{priority}. Hard schema
    contradictions and exact quotient identities can remove candidates without risk.
    Heuristic and domain scores normally change evaluation order only. Trajectory evidence
    may eliminate distinct quotient classes through an exact partial-loss bound or a stated
    finite-sample confidence rule. A prior may select among the surviving unresolved
    classes, but that selection is labeled as prior-driven rather than data-identified.


    The same separation creates a clean product boundary beyond one trajectory. A local run
    should not export only a winning enum or even only a posterior. It should export a
    versioned likelihood-like evidence packet over quotient classes, together with the local
    prior, calibration domain, excitation diagnostics, and source lineage. A second-stage
    cohort service can then ask whether several packets share one convention, are merely
    exchangeable members of one family, or are similar enough to justify partial borrowing.
    That follow-on service is developed in Section~\ref{sec:cohort-main} and
    Appendix~\ref{app:cohort}; it never modifies the local geometric score $J_C$.

    \begin{sectionrecap}
        The method is an exhaustive integrate-and-match engine surrounded by five explicit
        modules: normalize, quotient/diagnose, classify heuristically, apply provenance priors,
        and accumulate geometric evidence. The output is allowed to remain set-valued.

        \noindent\textbf{Key equation.}
        Equation~\eqref{eq:baseline-idea} is the engineering baseline; all later machinery
        specifies when candidates may be removed, merely reordered, or reported as equivalent.
    \end{sectionrecap}


    \section{Problem statement, data inventory, and canonical normalization}
    \label{sec:problem}

    Let the recorded data be
    \begin{equation}
        \mathcal D=\{(x_k,y_k,t_k)\}_{k=0}^{N},
        \label{eq:data-record}
    \end{equation}
    where $x_k$ denotes an attitude record when present, $y_k\in\mathbb R^3$ denotes three
    angular-rate channels when present, and $t_k$ is an optional timestamp. A complete
    hypothesis is factored as
    \begin{equation}
        h=(a,g,\xi,\nu)\in\mathcal H,
        \label{eq:focused-hypothesis}
    \end{equation}
    with attitude decoder $a$, rate decoder $g$, timing model $\xi$, and optional nuisance
    parameters $\nu$ such as bias, latency, or a permitted scale correction.

    \subsection{Normalize semantics before comparing candidates}

    For every attitude hypothesis, decode the raw representation into the same semantic
    object: an active body-to-world rotation,
    \begin{equation}
        R_{a,k}=D_a(x_k)\in\SO.
        \label{eq:canonical-attitude}
    \end{equation}
    Euler angle units, sequence names, quaternion scalar position, and active/passive language
    are therefore parser/decoder concerns; after decoding they should not leak into the
    geometric score. Quaternions are an efficient implementation coordinate, but the
    semantic state is $R_{a,k}$.

    Likewise, decode each rate hypothesis as
    \begin{equation}
        \omega_{g,k}=\kappa_g C_g y_k\in\mathbb R^3,
        \label{eq:canonical-rate}
    \end{equation}
    where $C_g$ contains the permitted channel permutation/sign map and $\kappa_g$ converts
    the recorded rate unit to radians per second whenever the physical time scale is known.
    The timing model provides positive interval durations $\Delta t_{\xi,k}$. Thus the
    normalized record under $h$ is
    \begin{equation}
        \mathsf N_h(\mathcal D)
        =\bigl(R_{a,0:N},\;\omega_{g,0:N},\;\Delta t_{\xi,0:N-1}\bigr).
        \label{eq:normalized-record}
    \end{equation}
    If time scale and rate scale are both unknown, the engine must not manufacture separate
    SI quantities. It instead works in the identifiable combined coordinate such as
    $\lambda=s_t\kappa$ or $\lambda=\Delta\kappa$ until an external dimensional anchor is
    available.

    \subsection{Inventory the available streams before claiming recoverability}

    A large part of the apparent complexity disappears once the input regime is named. The
    following table is a front-end contract for the solver: it says what kind of result may
    be returned before any numerical search begins.

    \begingroup
    \scriptsize
    \setlength{\tabcolsep}{3pt}
    \begin{longtable}{@{}p{0.055\textwidth}p{0.055\textwidth}p{0.13\textwidth}p{0.19\textwidth}p{0.26\textwidth}p{0.23\textwidth}@{}}
        \caption{Data-availability modes and required result labels.}\label{tab:availability-modes}\\
        \toprule
        Att. & Rates & Timing information & Required result label & What the engine can legitimately test & What remains unresolved without an external anchor \\
        \midrule
        \endfirsthead
        \toprule
        Att. & Rates & Timing information & Required result label & What the engine can legitimately test & What remains unresolved without an external anchor \\
        \midrule
        \endhead
        \bottomrule
        \endlastfoot
        yes & yes & physical $\Delta t_k$ known
        & \texttt{trajectory-quotient}
        & Relative attitude/rate compatibility, signed channel maps, rate scale, and---with informative motion---latency and selected calibration parameters.
        & Exact gauge partners, semantic axis names without a reference, and unobserved intersample loops. \\
        yes & yes & affine clock or one unknown constant period
        & \texttt{scale-gauge}
        & Directional and noncommutative compatibility plus the combined dimensionless time--rate scale.
        & Clock scale and angular-rate unit scale separately unless one is anchored. \\
        yes & yes & order only; durations constrained
        & \texttt{order-partial}
        & Directional/path compatibility under the declared clock prior or low-dimensional time warp.
        & Absolute angular speed, physical time unit, and any freedom absorbed by the time-warp model. \\
        yes & yes & order only; every $\Delta_k$ free
        & \texttt{overfit/non-ID}
        & Little beyond directional compatibility after nuisance fitting.
        & Rate magnitude and most timing semantics; the clock can absorb the mismatch. \\
        yes & no & physical time known
        & \texttt{attitude-heuristic}
        & Decoder validity, intrinsic smoothness, path length, bounded-speed/acceleration priors, Euler branch/range behavior.
        & Independent kinematic confirmation and every attitude-space isometry preserved by the chosen heuristic. \\
        yes & no & order only
        & \texttt{attitude-heuristic}
        & Ordered path geometry, branch continuity, singularity exposure, wrap/range/source-profile cues.
        & Physical angular speed and derivative-based claims. \\
        no & yes & physical time known
        & \texttt{rate-external}
        & Output matching only if endpoint attitude, vector observations, commanded dynamics, or another external target is added.
        & Convention recovery from self-integration alone, including initial-attitude and body/world gauges. \\
        no & yes & no physical time
        & \texttt{rate-non-ID}
        & Source-profile and raw-distribution plausibility only, unless additional dimensional/kinematic anchors exist.
        & Rate unit, clock scale, attitude, frame, and absolute orientation. \\
        no & no & any
        & \texttt{no-data}
        & Metadata/schema inspection only.
        & All trajectory-based convention claims. \\
    \end{longtable}
    \endgroup

    \paragraph{How to read the label column.}
    The label is part of the result, not decoration. A solver should not return a bare
    ``best convention.'' For example, \texttt{trajectory-quotient} means that the
    attitude/rate cross-check can rank structural quotient classes; the final evidence set
    may contain one class or several. Even when only one class survives, this does \emph{not}
    mean that a unique member of that class was observed. By contrast,
    \texttt{attitude-heuristic} says that no independent rate witness existed and that
    any ordering among decoders came from validity, smoothness, wrap, range, or source-profile
    assumptions. These labels prevent a downstream caller from treating qualitatively
    different evidence regimes as equivalent.


    \begin{sectionrecap}
        Every candidate is normalized to the same semantic objects before comparison. The
        availability label determines whether the run can produce trajectory evidence, only a
        scale-gauge result, heuristic attitude ranking, external-target rate matching, or no
        kinematic result at all.

        \noindent\textbf{Key equations.}\par
        $h=(a,g,\xi,\nu)$,
        $R_{a,k}=D_a(x_k)\in\SO$,
        $\omega_{g,k}=\kappa_gC_gy_k$, and
        $\mathsf N_h(\mathcal D)=(R_{a,0:N},\omega_{g,0:N},\Delta t_{\xi,0:N-1})$.

        \noindent\textbf{Symbols introduced.}
        \begin{center}
            \small
            \begin{tabularx}{0.96\linewidth}{@{}p{0.24\linewidth}Y@{}}
\toprule
Symbol & Meaning \\
\midrule
$\mathcal D$ & raw recorded data \\
$a,g,\xi,\nu$ & attitude, rate, timing, and nuisance parts of one hypothesis \\
$R_{a,k}$ & canonical active body-to-world rotation \\
$\omega_{g,k}$ & decoded angular velocity in the candidate expression frame \\
$\mathsf N_h$ & canonical normalization map for hypothesis $h$ \\
\bottomrule
            \end{tabularx}
        \end{center}
    \end{sectionrecap}


    \section{The five-module recovery engine}\label{sec:workflow}

    The engine is easiest to reason about as a sequence of candidate sets. Let
    $\mathcal H_{\mathrm{raw}}$ be every syntactically plausible decoder combination. Hard
    metadata and numerical validity produce $\mathcal H_0$. Exact observational equivalence
    then produces the quotient set $\mathcal Q_0=\mathcal H_0/\!\sim_s$. Heuristic and
    provenance information may order that set, but normally does not remove elements. Group
    integration and validation produce an evidence-preserving contender set
    $\mathcal T_{\mathrm{ev}}\subseteq\mathcal Q_0$. Only after that set is formed may an
    external prior or canonical convention choose one operational class and representative.
    The complete flow is
    \begin{equation}
        \boxed{
            \mathcal D
            \longrightarrow \mathcal H_0
            \longrightarrow \mathcal Q_0
            \dashrightarrow \mathcal Q_0^{\mathrm{priority}}
            \longrightarrow \mathcal T_{\mathrm{ev}}
            \longrightarrow (C_{\mathrm{op}},h_{\mathrm{op}}),}
        \label{eq:master-pipeline}
    \end{equation}
    where the dashed arrow denotes \emph{priority only}: it may change evaluation order but
    not the evidence semantics.

    \begin{center}
        \small
        \begin{tabularx}{\linewidth}{@{}p{0.13\linewidth}p{0.24\linewidth}p{0.23\linewidth}Y@{}}
\toprule
Module & Question & Main output & May remove candidates? \\
\midrule
1. Generate and normalize & What interpretations are syntactically and dimensionally admissible? & $\mathcal H_0$ and cached canonical streams & Yes, for hard metadata or invalid numerics only. \\
2. Quotient and diagnose & Which hypotheses are exactly the same observable model, and is the maneuver sufficiently exciting? & $\mathcal Q_0$ plus rank/conditioning report & Exact quotienting yes; weak excitation no. \\
3. Representation heuristics & Do wrapping, units, Euler ranges, or branch behavior favor some decoders? & heuristic key $K_{\mathrm{heur}}$ & Normally no; it orders/screens with an abstention mode. \\
4. Domain prior & What conventions are typical for the suspected source context? & $\pi_{\mathrm{dom}}(C\mid z)$ and representative preference & Hard only for exact documented schema; otherwise no before evidence. \\
5. Integrate and accumulate evidence & Which quotient classes actually predict the observed relative rotations, and when can losers be dropped? & $J_C$, active contender sets $\mathcal T_n$, final $\mathcal T_{\mathrm{ev}}$ & Yes, by exact lower bounds or stated statistical rules. \\
\bottomrule
        \end{tabularx}
    \end{center}

    \subsection{Module 1: generate and normalize hypotheses}

    The raw Cartesian product can be written
    \begin{equation}
        \mathcal H_{\mathrm{raw}}
        =\mathcal A\times\mathcal G\times\mathcal T_{\rm clk}\times\mathcal N,
        \label{eq:raw-cartesian}
    \end{equation}
    where $\mathcal A$ contains attitude decoders, $\mathcal G$ rate decoders,
    $\mathcal T_{\rm clk}$ timing models, and $\mathcal N$ allowed nuisance models. This is where the
    engine enumerates quaternion storage conventions, Euler sequences and angle units,
    axis/sign maps, body/world rate expression, candidate time scales, and permitted latency
    or bias models. Decoder outputs are cached in the canonical form of
    Section~\ref{sec:problem}. Trusted schema facts are applied here as hard constraints;
    ``probably NASA'' or ``looks like a game engine'' is not a hard constraint.

    \subsection{Module 2: quotient exact duplicates and diagnose observability}

    Define structural observational equivalence by equality of the forward observable map:
    \begin{equation}
        h\sim_s h'
        \quad\Longleftrightarrow\quad
        F_h(\mathcal D)=F_{h'}(\mathcal D)
        \ \text{for every admissible data record }\mathcal D.
        \label{eq:structural-equivalence-focused}
    \end{equation}
    The search should score one representative of each class
    \begin{equation}
        C=[h],\qquad \mathcal Q_0=\mathcal H_0/\!\sim_s,
        \label{eq:focused-quotient}
    \end{equation}
    while retaining the full member list for reporting. Important exact examples include
    quaternion sign, Euler intrinsic/extrinsic reversal with reversed channels, the
    attitude-direction/rate-frame/global-sign dual
    \begin{equation}
        (R,\omega,B/W)\longleftrightarrow(R^{\mathsf T},-\omega,W/B),
        \label{eq:focused-global-dual}
    \end{equation}
    and the unknown time/rate-scale gauge. These pairs remain tied even with perfect,
    full-rank trajectories; Appendix~\ref{app:identifiability} gives the proofs and full
    catalog.

    Excitation is a different issue. When a three-channel rate stream exists, a useful
    front-end diagnostic can be computed before choosing any signed-permutation decoder.
    For positive weights $w_k^{\rm exc}$ (physical durations when known, otherwise unit
    weights), form the raw-channel Gramian
    \begin{equation}
        G_y=\sum_{k\in\mathcal V}w_k^{\rm exc}y_k y_k^{\mathsf T},
        \qquad
        \lambda_1\ge\lambda_2\ge\lambda_3\ge0,
        \label{eq:focused-excitation-gramian}
    \end{equation}
    and, for $\lambda_1>0$, report
    \begin{equation}
        \operatorname{rank}_{\varepsilon_{\rm rank}}(G_y)
        =\sum_{i=1}^{3}\mathbf 1\{\lambda_i\ge\varepsilon_{\rm rank}\lambda_1\},
        \qquad
        \eta_{\rm exc}(G_y)=\frac{\lambda_3}{\lambda_1}.
        \label{eq:focused-excitation}
    \end{equation}
    For the base rate decoder $\omega=\kappa Cy$ with signed-permutation $C$,
    $G_\omega=\kappa^2 C G_y C^{\mathsf T}$, so $\operatorname{rank}_{\varepsilon_{\rm rank}}$ and $\eta_{\rm exc}$ are invariant
    to the unknown common rate scale, channel order, and channel signs. Thus
    $\operatorname{rank}_{\varepsilon_{\rm rank}}<3$ is a direct warning that the record does not provide well-conditioned
    three-axis excitation. Conversely, $\operatorname{rank}_{\varepsilon_{\rm rank}}=3$ is \emph{not} a uniqueness
    theorem: exact quotient symmetries and symmetric maneuvers can still tie. With attitude
    but no rates, an analogous decoder-dependent increment Gramian can be reported instead;
    Appendix~\ref{app:excitation} gives both forms. Excitation diagnostics de-rate confidence
    but do not reject a candidate merely for being weakly excited.

    \subsection{Module 3: classify units and representation behavior heuristically}

    Representation clues can be highly useful before expensive integration, especially for
    Euler-valued logs. Candidate one-turn periods are
    \begin{equation}
        T_u\in\{2\pi,360,1,400\}
        \quad\text{for}\quad
        u\in\{\mathrm{rad},\mathrm{deg},\mathrm{turn},\mathrm{grad}\}.
        \label{eq:focused-unit-periods}
    \end{equation}
    A candidate unit is favored when subtracting integer multiples of $T_u$ explains a small
    number of otherwise extreme boundary jumps. If no candidate boundary crossing occurs,
    this feature must abstain: remaining inside roughly $[-\pi,\pi]$ is compatible with
    radians but does not prove radians.

    Euler principal ranges provide a second, conditional cue. Under the common principal
    branch, a Tait--Bryan middle angle lies in $[-\pi/2,\pi/2]$ and a proper-Euler middle
    angle in $[0,\pi]$. Exceeding those intervals contradicts a producer that promises
    principal output, but it does not invalidate the physical rotation because equivalent
    nonprincipal branches exist. Appendix~\ref{app:heuristics} defines the wrap-gain,
    branch-continuity, singularity, and range scores in detail.

    The output of this module is therefore a labeled key
    \begin{equation}
        K_{\mathrm{heur}}(h)
        =(H_{\mathrm{wrap}},H_{\mathrm{range}},H_{\mathrm{branch}},f_{\mathrm{sing}},\ldots),
        \label{eq:focused-heur-key}
    \end{equation}
    not an unlabeled probability unless a classifier has been trained and calibrated on a
    representative labeled corpus.

    \subsection{Module 4: construct a domain-conditioned prior}

    Let the provenance context be
    \begin{equation}
        z=(d,r,p,\iota),
        \label{eq:focused-provenance}
    \end{equation}
    for broad domain $d$, artifact role $r$, ecosystem/protocol $p$, and exact API/schema
    $\iota$. Prior strength follows
    \begin{equation}
        \boxed{
            \text{exact schema/API}\succ
            \text{ecosystem/protocol}\succ
            \text{artifact role}\succ
            \text{broad domain}\succ
            \text{uninformative}.}
        \label{eq:focused-prior-order}
    \end{equation}
    An exact message or file specification should normally be promoted to hard metadata.
    Broader context supplies a class prior $\pi_{\mathrm{dom}}(C\mid z)$ and a separate
    within-class representative preference. It may order the search early, but it is not
    allowed to overwrite a later geometric loss. If provenance itself is uncertain, an
    upstream filename/header/schema classifier may return a distribution
    $q_{\rm ctx}(z\mid\mathfrak m)$ from metadata features $\mathfrak m$; the engine should then
    marginalize rather than pretend the guessed domain is exact:
    \begin{equation}
        \bar\pi_{\mathrm{dom}}(C\mid\mathfrak m)
        =\sum_z q_{\rm ctx}(z\mid\mathfrak m)\,\pi_{\mathrm{dom}}(C\mid z).
        \label{eq:focused-provenance-mixture}
    \end{equation}
    A bag-of-words, fuzzy-name, or embedding match therefore changes prior weight only; a
    verified schema remains a hard constraint. After trajectory evidence has produced
    $\mathcal T_{\mathrm{ev}}$, the prior may rank only those surviving classes. The survey
    and its NASA/JPL, ROS, MAVLink/PX4, scientific-library, and game-engine scorecards are
    reported in Appendix~\ref{app:survey}.

    \subsection{Module 5: integrate, score, accumulate evidence, and eliminate}

    For a body-expressed rate, the kinematics are
    \begin{equation}
        \dot R=R[\omega_B]_\times,
    \end{equation}
    while world-expressed rate uses $\dot R=[\omega_W]_\times R$. The corresponding
    observed increments are
    \begin{equation}
        \Delta R_{B,k}=R_k^{\mathsf T}R_{k+1},
        \qquad
        \Delta R_{W,k}=R_{k+1}R_k^{\mathsf T}.
        \label{eq:focused-increments}
    \end{equation}
    Each quotient-class representative is integrated using its own timing model and compared
    with the compatible observed increments using the intrinsic distance $d_g$.
    For a window of $\ell$ intervals define
    \begin{equation}
        \rho_{C,k,\ell}
        =d_g(P_{C,k}^{(\ell)},\Delta R_{C,k}^{(\ell)}),
    \end{equation}
    and rank classes by the robust multiscale risk
    \begin{equation}
        J_C
        =\sum_{\ell\in\mathcal L}w_\ell
        T_{\alpha_{\rm trim}}\!\left(
                                   \left\{\frac{\rho_{C,k,\ell}^2}{\ell}\right\}_k
        \right),
        \label{eq:focused-risk}
    \end{equation}
    where $\alpha_{\rm trim}\in[0,1)$ is the configured trimming fraction and $T_{\alpha_{\rm trim}}$ retains the smallest fraction $1-\alpha_{\rm trim}$ of the per-window losses. Section~\ref{sec:core-score}
    provides the full group integration and scoring definitions.

    There are two different ways a candidate can stop consuming compute, and the distinction
    is important.

    \paragraph{Exact computational abandonment.}
    If a partial lower bound $\mathrm{LB}_r(C)$ on the final robust risk already exceeds the
    verified incumbent score $J_{\mathrm{inc}}$, then
    \begin{equation}
        \mathrm{LB}_r(C)>J_{\mathrm{inc}}
        \quad\Longrightarrow\quad
        \text{stop evaluating class }C.
        \label{eq:focused-exact-abandon}
    \end{equation}
    This is a computational shortcut; it does not change the exhaustive answer.

    \paragraph{Statistical evidence elimination.}
    For independent bounded validation blocks, let
    \begin{equation}
        X_{C,j}=\frac{\rho_{C,j}^2}{\pi^2}\in[0,1],
        \qquad
        \widehat L_n(C)=\frac1n\sum_{j=1}^nX_{C,j}.
    \end{equation}
    Using the simultaneous radius
    \begin{equation}
        r_n(\delta)
        =\sqrt{\frac{1}{2n}
        \log\!\left(\frac{\pi^2|\mathcal Q_0|n^2}{3\delta}\right)},
        \label{eq:focused-anytime-radius}
    \end{equation}
    define the active evidence set
    \begin{equation}
        \mathcal T_n(\delta)
        =\left\{C\in\mathcal Q_0:
             \widehat L_n(C)-\min_{C'\in\mathcal Q_0}\widehat L_n(C')
             \le2r_n(\delta)\right\}.
        \label{eq:focused-active-set}
    \end{equation}
    A class leaves the active set only when its lower confidence bound exceeds the best
    competing upper confidence bound. Under the stated bounded-IID assumptions, this rule
    preserves every population-risk minimizer on the simultaneous confidence event.
    For a post-quotient population gap $\gamma>0$, let
    $q_{\rm bad}=|\mathcal Q_0|-|\mathcal Q_\star|\ge1$ be the number of nonoptimal
    classes. The fixed-sample guarantee implies the useful planning relation
    \begin{equation}
        n\ge\frac{2}{\gamma^2}
        \log\!\left(\frac{q_{\rm bad}}{\delta}\right)
        \quad\Longrightarrow\quad
        \Pr(\widehat C_n\notin\mathcal Q_\star)\le\delta,
        \label{eq:focused-sample-complexity}
    \end{equation}
    when the theorem's independence and model assumptions hold. Serially correlated or
    overlapping residuals require disjoint blocks, a justified dependence bound, or a block
    bootstrap rather than this formula.

    \paragraph{Emit a reusable evidence packet.}
    The final act of Module~5 is serialization rather than another tie breaker. For packet
    $i$, define the class-relative log-evidence vector
    \begin{equation}
        \lambda_i(C)
        =\log L_i(C)-\max_{G\in\mathcal Q_0}\log L_i(G),
        \qquad \max_C\lambda_i(C)=0,
        \label{eq:focused-packet-log-evidence}
    \end{equation}
    where $L_i$ is either a stated likelihood or a calibrated generalized-evidence map.
    The engine stores $\lambda_i$, the local prior $\pi_i^{(0)}$, the quotient-space version,
    $\mathcal T_{\mathrm{ev},i}$, observability diagnostics, hard metadata, provenance,
    calibration identity, and atomic source lineage. Storing the local prior is essential:
    a downstream service must be able to remove it before applying a different cohort prior.
    An uncalibrated risk vector may still be exported, but the packet must label its semantics
    as \texttt{weight-only} rather than \texttt{probability-ready}. Section~\ref{sec:cohort-main}
    uses this packet as the interface between local inference and cohort borrowing.

    \begin{sectionrecap}
        The candidate life cycle is now explicit. Hard constraints and exact quotient identities
        may remove hypotheses without approximation. Excitation/conditioning, wrap/Euler cues, and broad
        domain knowledge normally change confidence or priority, not truth. Integration supplies
        the geometric evidence, and candidates are eliminated either by an exact partial-risk
        bound or by a stated evidence rule.

        \noindent\textbf{Key equations.}
        The master flow is \eqref{eq:master-pipeline}; exact structural reduction is
        $\mathcal Q_0=\mathcal H_0/\!\sim_s$; excitation uses
        \eqref{eq:focused-excitation}; the final geometric risk is \eqref{eq:focused-risk}; and
        sequential evidence is represented by the active set \eqref{eq:focused-active-set}.

        \noindent\textbf{Symbols introduced.}
        \begin{center}
            \small
            \begin{tabularx}{0.97\linewidth}{@{}p{0.25\linewidth}Y@{}}
\toprule
Symbol & Meaning \\
\midrule
$\mathcal H_0$ & hypotheses surviving hard metadata and numerical validity \\
$\sim_s,\mathcal Q_0$ & exact structural equivalence and quotient-class set \\
$G_y,\operatorname{rank}_{\varepsilon_{\rm rank}},\eta_{\rm exc}$ & raw-rate excitation Gramian, thresholded rank, and weakest-axis ratio \\
$K_{\mathrm{heur}},T_u$ & nonprobabilistic representation/unit heuristic key and candidate one-turn period \\
$z,\mathfrak m,q_{\rm ctx},\pi_{\mathrm{dom}}$ & provenance context, soft metadata features/classifier, and domain-conditioned class prior \\
$J_C,\mathrm{LB}_r,J_{\mathrm{inc}}$ & class trajectory risk, exact partial lower bound, and verified incumbent score \\
$r_n(\delta),\gamma$ & simultaneous statistical radius and post-quotient population-risk gap \\
$\mathcal T_n,\mathcal T_{\mathrm{ev}}$ & running and final evidence-preserving contender sets \\
\bottomrule
            \end{tabularx}
        \end{center}
    \end{sectionrecap}


    \section{Core geometric prediction and score}\label{sec:core-score}

    The workflow above can be implemented with quaternions, but the mathematical score is
    best stated intrinsically on $\SO$. This section collects only the equations required by
    the engine; the identifiability proofs and special cases are deferred to the appendices.

    \subsection{Rotation-group geometry and body/world kinematics}

    For $u\in\R^3$, let $[u]_\times\in\so$ satisfy $[u]_\times v=u\times v$.
    Define the vector exponential and principal logarithm by
    \begin{equation}
        \Exp(u)=\exp([u]_\times),
        \qquad
        \Log(Q)=\bigl(\log Q\bigr)^\vee,
    \end{equation}
    where $\Log(Q)$ is single-valued for rotations with angle strictly less than $\pi$.
    At angle $\pi$, the rotation axis has a sign ambiguity, although the angle remains
    well-defined. The canonical geodesic distance is
    \begin{align}
        d_g(A,B)
        &=\norm{\Log(A^{\T} B)} \nonumber\\
        &=\arccos\!\left(\clip\!\left(\frac{\tr(A^{\T} B)-1}{2},-1,1\right)\right)
        \in[0,\pi].
    \end{align}
    The trace formula defines the distance even at the cut locus. This metric is left- and
    right-invariant \cite{huynh2009}:
    \begin{equation}
        d_g(QA,QB)=d_g(AQ,BQ)=d_g(A,B),
    \end{equation}
    and inversion is an isometry:
    \begin{equation}
        d_g(A^{-1},B^{-1})=d_g(A,B).
    \end{equation}

    Let $R(t)\in\SO$ map body coordinates into world coordinates. Body- and
    world-expressed angular velocities obey
    \begin{align}
        \dot R(t)&=R(t)[\omega_B(t)]_\times, \\
        \dot R(t)&=[\omega_W(t)]_\times R(t),
    \end{align}
    with
    \begin{equation}
        \omega_W(t)=R(t)\omega_B(t).
    \end{equation}
    These standard kinematic identities are reviewed in \cite{barfoot2017,markley2014}.
    The corresponding observed increments are
    \begin{align}
        \Delta R_{B,k}&=R_k^{\T} R_{k+1}, \\
        \Delta R_{W,k}&=R_{k+1}R_k^{\T}.
    \end{align}
    They are conjugate and therefore have the same rotation angle, but their axes are
    expressed in different coordinates.

    \begin{remark}[Angular velocity is axial]
If a coordinate-basis hypothesis includes an improper orthogonal transformation
$Q\in\OO(3)$, angular velocity cannot be transformed as an ordinary polar vector.
The identity
\begin{equation}
    Q^{\T}[\omega]_\times Q
    =[\det(Q)Q^{\T}\omega]_\times
\end{equation}
shows the required axial-vector sign. A search that reflects attitude axes but transforms
rates as polar vectors is physically inconsistent. A mere sensor-channel sign decoder
$C_h y_k$, however, need not be interpreted as a basis transformation.
    \end{remark}

    \begin{sectionrecap}
        The comparison is performed intrinsically on $\SO$, so it is independent of a chosen
        matrix or quaternion coordinate chart. Body and world angular velocities satisfy
        different left/right kinematics, and improper basis changes require the axial-vector
        determinant sign.

        \noindent\textbf{Key equations.}
        \(d_g(A,B)=\arccos((\operatorname{tr}(A^{\mathsf T}B)-1)/2)\) after clipping;
        \(\dot R=R[\omega_B]_\times=[\omega_W]_\times R\);
        \(\omega_W=R\omega_B\); and
        \(Q^{\mathsf T}[\omega]_\times Q=[\det(Q)Q^{\mathsf T}\omega]_\times\).

        \noindent\textbf{Symbols introduced.}
        \par\noindent\begin{tabularx}{\linewidth}{@{}p{0.22\linewidth}Y@{}}
\toprule
Symbol & Meaning \\
\midrule
$\Exp,\Log$ & vector exponential and principal rotation logarithm \\
$d_g$ & bi-invariant geodesic rotation distance in $[0,\pi]$ \\
$R$ & active body-to-world rotation \\
$\omega_B,\omega_W$ & angular velocity expressed in body and world coordinates \\
$\Delta R_{B,k},\Delta R_{W,k}$ & observed relative rotations in the two frames \\
$Q$ & orthogonal coordinate transformation; improper cases have $\det Q=-1$ \\
\bottomrule
        \end{tabularx}
    \end{sectionrecap}

    \subsection{Discrete Lie-group integration}

    Fix one interval $[t_k,t_{k+1}]$, let $M\ge1$ be the number of integration substeps,
    and set
    \begin{equation}
        \Delta t_k^{(M)}=\frac{\Delta t_k}{M},
        \qquad
        \alpha_j=\frac{j+1/2}{M},
        \qquad
        \omega_{k,j}=(1-\alpha_j)\omega_k+\alpha_j\omega_{k+1}.
    \end{equation}
    Define
    \begin{equation}
        A_{k,j}=\Exp\!\left(\Delta t_k^{(M)}\omega_{k,j}\right).
    \end{equation}
    Product notation is easy to misread on a noncommutative group, so the algorithm should
    be defined recursively. For body-expressed rates,
    \begin{equation}
        P_{B,k}^{(0)}=\I,
        \qquad
        P_{B,k}^{(j+1)}=P_{B,k}^{(j)}A_{k,j}.
    \end{equation}
    For world-expressed rates,
    \begin{equation}
        P_{W,k}^{(0)}=\I,
        \qquad
        P_{W,k}^{(j+1)}=A_{k,j}P_{W,k}^{(j)}.
    \end{equation}
    The predicted interval increments are $P_{B,k}=P_{B,k}^{(M)}$ and
    $P_{W,k}=P_{W,k}^{(M)}$. Thus
    \begin{align}
        P_{B,k}&=A_{k,0}A_{k,1}\cdots A_{k,M-1}, \\
        P_{W,k}&=A_{k,M-1}\cdots A_{k,1}A_{k,0}.
    \end{align}
    Each approximation lies exactly in $\SO$.

    \begin{proposition}[Accuracy of exponential midpoint substepping]\label{prop:midpoint-accuracy}
        Suppose the true angular velocity is twice continuously differentiable on an interval of
        length $T$, and its norm and first two derivatives are bounded. If exact midpoint values
        are used, the substepped exponential-midpoint approximation has interval error
        \begin{equation}
            d_g(P_{\mathrm{exact}},P_M)\le K_{\mathrm{int}}\frac{T^3}{M^2}
        \end{equation}
        for sufficiently small $T/M$, with $K_{\mathrm{int}}$ depending on the stated bounds. If midpoint
        values are obtained from endpoint-linear interpolation, an additional $O(T^3)$ model
        error remains even as $M\to\infty$. Consequently, increasing $M$ controls numerical
        quadrature error but cannot remove error caused by the assumed interpolation model.
    \end{proposition}

    \begin{proof}[Proof sketch]
On each subinterval, exponential midpoint is a second-order Lie-group method, so its
local defect is $O((T/M)^3)$. Bi-invariance of $d_g$ allows the $M$ local defects to be
accumulated without leaving the group, giving $M O((T/M)^3)=O(T^3/M^2)$ over the
interval. Linear interpolation approximates a $C^2$ signal with pointwise error
$O(T^2)$; its integrated effect over an interval is $O(T^3)$. Standard Lie-group and
Magnus-expansion analyses make these statements precise \cite{iserles2000,blanes2009}.
    \end{proof}

    A practical numerical diagnostic is therefore to rescore leading hypotheses with $M$
    and $2M$. If the score gap or winner changes materially, integration error is not yet
    subordinate to measurement/model error.

    \begin{sectionrecap}
        Exponential midpoint substepping keeps every prediction exactly on $\SO$ and makes the
        noncommutative product order explicit. Increasing $M$ removes numerical integration
        error at second order but cannot remove error in the assumed intersample rate model.

        \noindent\textbf{Key equations.}
        \(A_{k,j}=\Exp(\Delta t_k^{(M)}\omega_{k,j})\),
        \(P_{B}^{(j+1)}=P_B^{(j)}A_j\), and
        \(P_W^{(j+1)}=A_jP_W^{(j)}\). With exact midpoint values,
        \(d_g(P_{\mathrm{exact}},P_M)\le K_{\mathrm{int}}T^3/M^2\); endpoint-linear interpolation contributes
        an additional $O(T^3)$ model term.

        \noindent\textbf{Symbols introduced.}
        \par\noindent\begin{tabularx}{\linewidth}{@{}p{0.22\linewidth}Y@{}}
\toprule
Symbol & Meaning \\
\midrule
$M$ & integration substeps per recorded interval \\
$T,K_{\rm int}$ & interval length and integration-error constant in Proposition~\ref{prop:midpoint-accuracy} \\
$\Delta t_k^{(M)}$ & substep duration $\Delta t_k/M$ \\
$\alpha_j$ & normalized substep midpoint \\
$\omega_{k,j}$ & interpolated angular velocity at that midpoint \\
$A_{k,j}$ & one exponential substep \\
$P_{B,k},P_{W,k}$ & predicted body- and world-frame interval increments \\
\bottomrule
        \end{tabularx}
    \end{sectionrecap}

    \subsection{Intrinsic residuals and robust multiscale risk}

    Exact quotienting precedes numerical scoring. Let $C\in\mathcal Q_0$ be one structural
    quotient class and let $h_C\in C$ be the representative stored by the implementation.
    Write the predicted and observed increments produced by that representative as
    $P_{C,k}$ and $\Delta R_{C,k}^{\mathrm{obs}}$. Structural equivalence requires every
    member of $C$ to produce the same risk; this equality is an implementation invariant,
    not an additional modeling assumption.

    Define the interval error rotation and geodesic residual by
    \begin{equation}
        E_{C,k}=P_{C,k}^{-1}\Delta R_{C,k}^{\mathrm{obs}},
        \qquad
        \rho_{C,k}=d_g(P_{C,k},\Delta R_{C,k}^{\mathrm{obs}}).
    \end{equation}
    When $\rho_{C,k}<\pi$, the tangent residual
    \begin{equation}
        e_{C,k}=\Log(E_{C,k})
    \end{equation}
    is well-defined and satisfies $\norm{e_{C,k}}=\rho_{C,k}$. Statistical procedures that
    use the residual vector, rather than only its norm, must flag samples near $\pi$ because
    the axis of the principal logarithm is ill-conditioned there.

    \begin{proposition}[Noise-free consistency]
If a quotient class contains a candidate that decodes the true canonical attitude and
angular velocity, and if its forward model is integrated exactly, then
$E_{C,k}=\I$ and $\rho_{C,k}=0$ for every interval.
    \end{proposition}

    \begin{proof}
        For body rates, $R_{k+1}=R_kP_{B,k}$ and hence $R_k^{\T} R_{k+1}=P_{B,k}$. For world
        rates, $R_{k+1}=P_{W,k}R_k$ and hence $R_{k+1}R_k^{\T}=P_{W,k}$.
    \end{proof}

    This is a consistency statement, not an identifiability statement. Multiple raw
    hypotheses may belong to the same exact class, and multiple distinct classes may remain
    empirically competitive.

    For a window of $\ell$ intervals, define
    \begin{align}
        \Delta R_{B,k}^{(\ell)}&=R_k^{\T} R_{k+\ell},
        &P_{B,k}^{(\ell)}&=P_{B,k}P_{B,k+1}\cdots P_{B,k+\ell-1},\\
        \Delta R_{W,k}^{(\ell)}&=R_{k+\ell}R_k^{\T},
        &P_{W,k}^{(\ell)}&=P_{W,k+\ell-1}\cdots P_{W,k+1}P_{W,k}.
    \end{align}
    For the representative of class $C$, let
    \begin{equation}
        \rho_{C,k,\ell}
        =d_g\!\left(P_{C,k}^{(\ell)},\Delta R_{C,k}^{(\ell),\mathrm{obs}}\right).
    \end{equation}

    For nonnegative numbers $z_1,\ldots,z_n$, let
    $z_{(1)}\le\cdots\le z_{(n)}$ be their order statistics and define the upper-trimmed
    mean
    \begin{equation}
        T_{\alpha_{\rm trim}}(z_1,\ldots,z_n)
        =\frac{1}{m}\sum_{i=1}^{m}z_{(i)},
        \qquad
        m=\max\set{1,\left\lfloor(1-\alpha_{\rm trim})n\right\rfloor}.
    \end{equation}
    For a configured window set $\mathcal L$ and weights $w_\ell>0$ with
    $\sum_{\ell\in\mathcal L}w_\ell=1$, define the class-level robust geodesic risk
    \begin{equation}
        J_C
        =\sum_{\ell\in\mathcal L}w_\ell
        T_{\alpha_{\rm trim}}\!\left(
                                   \left\{\frac{\rho_{C,k,\ell}^2}{\ell}\right\}_{k=0}^{N-\ell}
        \right),
        \qquad C\in\mathcal Q_0.
        \label{eq:robust-risk}
    \end{equation}
    The geometric data ranking is therefore set-valued at the class level:
    \begin{equation}
        \widehat{\mathcal C}
        =\argmin_{C\in\mathcal Q_0}J_C.
    \end{equation}
    Using one trimmed mean per scale prevents a scale with more available windows from
    receiving an accidental weight. The factor $1/\ell$ is the squared form of dividing
    window angles by $\sqrt{\ell}$.

    \begin{lemma}[Small-noise justification for window scaling]
Suppose interval error rotations are $\Exp(e_j)$ with independent,
zero-mean, isotropic small errors satisfying
$\operatorname{Cov}(e_j)=\sigma^2\I$. Then the $\ell$-interval error satisfies
\begin{equation}
    \operatorname{Cov}\!\left(\Log E^{(\ell)}\right)
    =\ell\sigma^2\I+O(\!\sigma^3),
\end{equation}
so $\rho_{k,\ell}/\sqrt{\ell}$ has approximately scale-independent RMS.
    \end{lemma}

    \begin{proof}[Proof sketch]
The Baker--Campbell--Hausdorff expansion gives
\begin{equation}
    \Log E^{(\ell)}
    =\sum_{j=1}^{\ell}Q_j e_j
    +O\!\left(\sum_j\norm{e_j}^2\right),
\end{equation}
where each $Q_j\in\SO$ is an adjoint rotation induced by intervening predicted
increments. Isotropy gives
$\operatorname{Cov}(Q_je_j)=Q_j(\sigma^2\I)Q_j^{\T}=\sigma^2\I$, and independence
makes the first-order covariances add.
    \end{proof}

    The lemma does not make overlapping windows independent, and it does not justify the
    same scaling for strongly anisotropic or biased errors. In those cases, a propagated
    covariance or per-scale empirical calibration is preferable.

    A class-dependent normalization such as dividing $J_C^{1/2}$ by a motion magnitude
    decoded under the same representative should not determine the ranking: a wrong decoder
    can inflate its own denominator. Any normalized display score must use a denominator
    fixed independently of $C$ and must be labeled as a presentation heuristic, not a
    probability.

    Upper trimming is also not free. A wrong class can tie if the intervals that separate it
    from the truth occupy no more than the trimmed fraction. The report should therefore
    include untrimmed and trimmed per-scale risks and show whether the winner changes with
    $\alpha_{\rm trim}$.

    \begin{sectionrecap}
        The primary data score is an intrinsic geometric risk defined on structural quotient
        classes, with a normalization fixed independently of the candidate. Window scales are
        normalized according to the small independent-error approximation and are combined with
        explicit weights. Trimming is a robustness choice that must be audited against the
        untrimmed score because it can discard the only discriminating intervals.

        \noindent\textbf{Key equations.}
        \(E_{C,k}=P_{C,k}^{-1}\Delta R_{C,k}^{\mathrm{obs}}\),
        \(\rho_{C,k}=d_g(P_{C,k},\Delta R_{C,k}^{\mathrm{obs}})\), and
        \[
            J_C=\sum_{\ell\in\mathcal L}w_\ell
            T_{\alpha_{\rm trim}}\!\left(\left\{\rho_{C,k,\ell}^2/\ell\right\}_k\right).
        \]
        For small isotropic independent interval errors,
        $\operatorname{Cov}(\Log E^{(\ell)})\approx\ell\sigma^2I$.

        \noindent\textbf{Symbols introduced.}
        \par\noindent\begin{tabularx}{\linewidth}{@{}p{0.22\linewidth}Y@{}}
\toprule
Symbol & Meaning \\
\midrule
$h_C$ & stored representative used to evaluate quotient class $C$ \\
$E_{C,k},e_{C,k}$ & interval error rotation and its principal tangent residual \\
$\rho_{C,k,\ell}$ & geodesic residual angle for an $\ell$-interval window \\
$T_{\alpha_{\rm trim}}$ & upper-trimmed mean retaining the smallest fraction $1-\alpha_{\rm trim}$ \\
$\mathcal L,w_\ell$ & configured window lengths and their explicit weights \\
$J_C$ & robust multiscale geometric risk used for class-level data ranking \\
\bottomrule
        \end{tabularx}
    \end{sectionrecap}


    \section{Interpreting the result in plain language}\label{sec:result}

    The engine produces four different objects that should not be collapsed into a single
    ``best convention'' field.
    \begin{center}
        \small
        \begin{tabularx}{\linewidth}{@{}p{0.23\linewidth}Y@{}}
\toprule
Object & Plain-language meaning \\
\midrule
Hypothesis $h$ & One fully specified decoder: storage/order, units, axes, frame direction, timing, and nuisance choices. \\
Quotient class $C=[h]$ & Several hypotheses that are mathematically indistinguishable under the available forward observations because an exact symmetry maps one to another. \\
Evidence contender set $\mathcal T_{\mathrm{ev}}$ & One or more \emph{different} quotient classes that the measured trajectory has not separated at the selected numerical/statistical standard. \\
Operational representative $h_{\mathrm{op}}$ & One software decoder selected after inference for interoperability and reproducibility; its selection may be metadata-, prior-, or policy-driven. \\
\bottomrule
        \end{tabularx}
    \end{center}

    The recommended returned object is
    \begin{equation}
        \boxed{
            \mathcal R=
            (\mathcal T_{\mathrm{ev}},\,C_{\mathrm{op}},\,h_{\mathrm{op}},\,
            \ell_{\mathrm{reason}}),}
        \label{eq:focused-result-object}
    \end{equation}
    where $C_{\mathrm{op}}\in\mathcal T_{\mathrm{ev}}$ is the operational class when a
    single decoder is requested and $h_{\mathrm{op}}\in C_{\mathrm{op}}$ is its chosen
    representative; an inference-only caller may leave both operational fields unset. The reason label is
    mandatory. Examples include \texttt{trajectory-unique}, \texttt{exact-gauge-class},
    \texttt{hard-metadata}, \texttt{domain-prior}, \texttt{heuristic-only}, and
    \texttt{canonical-policy}.

    If $|\mathcal T_{\mathrm{ev}}|=1$, the data and stated assumptions selected one quotient
    class. This still does not mean a unique member of that class was observed. For
    example,
    \begin{align}
        h_1&=(\text{body-to-world attitude},\text{body rates},+\omega),\\
        h_2&=(\text{world-to-body attitude},\text{world rates},-\omega)
    \end{align}
    are exact global duals when all three semantic dimensions are left free. A run may
    therefore return one class $C=\{h_1,h_2\}$ even with zero residual and rich three-axis
    motion. Choosing $h_1$ because the suspected API normally uses that description is a
    representative policy, not additional trajectory evidence.

    If $|\mathcal T_{\mathrm{ev}}|>1$, the ambiguity is empirical rather than structural.
    More informative motion, longer validation, an external frame/vector anchor, or trusted
    metadata may shrink the set. If it cannot be shrunk, a domain prior may select an
    operational class only from within the retained set:
    \begin{equation}
        C_{\mathrm{op}}
        \in\argmax_{C\in\mathcal T_{\mathrm{ev}}}
        \pi_{\mathrm{dom}}(C\mid z).
        \label{eq:focused-prior-tiebreak}
    \end{equation}
    If the prior has no support or negligible margin inside the contender set, the prior
    must abstain and the inferential report remains the full set. If a caller nevertheless
    requires one operational decoder, a deterministic canonical class key supplies the
    fallback and the reason label becomes \texttt{canonical-policy}. Within the selected
    $C_{\mathrm{op}}$, choose one stable representative by the exact schema,
    source-profile preference, or the same declared minimum-intervention policy. Appendix~\ref{app:selection} gives the detailed hierarchy and
    explains why AIC/BIC, MDL, Bayesian priors, and canonicalization solve different parts of
    the decision problem.

    \subsection{Persistent preference profiles}
    \label{sec:preference-profile}

    Applications often need a stable decoder even when the scientific result is a class or
    a set of classes. We therefore make the operational policy explicit as a versioned
    preference profile
    \begin{equation}
        P=(\mathrm{id},\mathrm{source},\mathrm{context},
        w_{\mathrm{class}},w_{\mathrm{rep}},\beta_P),
        \label{eq:preference-profile}
    \end{equation}
    where $w_{\mathrm{class}}$ assigns nonnegative relative weights to quotient classes
    and $w_{\mathrm{rep}}$ assigns relative weights to convention dimensions of a
    representative. The class weights are evaluated only on
    $\mathcal T_{\mathrm{ev}}$. After $C_{\mathrm{op}}$ is chosen, the representative is
    selected by
    \begin{equation}
        h_{\mathrm{op}}
        \in \argmax_{h\in C_{\mathrm{op}}}
        \sum_{d\in\mathcal D(h)}\log w_{\mathrm{rep},d}(h_d),
        \label{eq:representative-preference}
    \end{equation}
    with the stable canonical key used only for an exact tie. Thus a profile can select
    the body-rate representative of a body/world global-dual class, but it cannot make the
    dual class cease to be an exact quotient or create trajectory evidence.

    The profile is stored in the ordinary JSON configuration and carried into the report
    as its identifier, source, supported class IDs, abstention state, and representative
    selection reason. This makes persistent user policy reproducible without confusing it
    with $J_C$, the data-only geometric risk. A profile with no positive support among the
    evidence contenders abstains; the scientific output remains set-valued and a required
    decoder falls back to \texttt{canonical-policy}.

    Research-backed profiles should be assembled from exact, versioned API and schema
    documentation before ecosystem surveys or naming heuristics. For each claim, record
    the artifact role, field names, source version, access date, whether it selects a class
    or merely a representative, and any conflicting documentation. The repository
    provides a worksheet in \texttt{docs/PREFERENCE\_RESEARCH.md}; relative weights in a
    profile are rankings until a held-out corpus validates their calibration.

    \begin{sectionrecap}
        The final scientific result is set-valued whenever the observations demand it. A quotient
        class represents an exact equivalence; a contender set represents unresolved evidence
        between distinct classes; the operational decoder is a subsequent implementation choice.

        \noindent\textbf{Key equations.}
        The result object is \eqref{eq:focused-result-object}. A domain prior may rank only the
        surviving evidence classes as in \eqref{eq:focused-prior-tiebreak}; it does not change
        $J_C$. A persistent profile extends this with the class/representative separation in
        \eqref{eq:preference-profile}--\eqref{eq:representative-preference}.
    \end{sectionrecap}


    \section{From one trajectory to a cohort evidence service}
    \label{sec:cohort-main}

    The local engine answers, ``Which quotient classes are supported by this trajectory?''
    A cohort service answers a different question: ``Given what trusted packets from related
    artifacts have established, how should the uncertainty for a weaker target packet be
    updated?'' The separation is important. Cohort information is a prior or a shared
    latent-variable model; it is not an additional residual secretly added to $J_C$.

    \subsection{The evidence packet is the product boundary}

    Let the common, versioned quotient universe be
    $\mathcal Q_v=\{C_1,\ldots,C_K\}$. A reusable packet for trajectory $i$ is
    \begin{equation}
        \boxed{
            \mathcal E_i=
            (v_{\mathcal Q},\lambda_i,\pi_i^{(0)},\mathcal T_i,
            \mathcal O_i,\mathcal M_i,z_i,\mathsf{lin}_i,\mathsf{cal}_i),}
        \label{eq:cohort-packet}
    \end{equation}
    where $\lambda_i$ is the local relative log evidence from
    \eqref{eq:focused-packet-log-evidence}; $\pi_i^{(0)}$ is the prior used in the local run;
    $\mathcal T_i$ is its evidence-preserving class set; $\mathcal O_i$ contains excitation
    and conditioning diagnostics; $\mathcal M_i$ contains hard anchors and representative
    metadata; $z_i$ is provenance context; $\mathsf{lin}_i$ records atomic measurement and
    derivation lineage; and $\mathsf{cal}_i$ identifies the score-to-probability calibration.
    Packets can be combined only when their quotient-space definitions are compatible.

    A packet that stores only a posterior $q_i$ can still expose relative local evidence if
    it also stores the strictly positive local prior:
    \begin{equation}
        \begin{split}
            \lambda_i(C)-\lambda_i(G)
            ={}&\log\frac{q_i(C)}{q_i(G)}
            -\log\frac{\pi_i^{(0)}(C)}{\pi_i^{(0)}(G)}.
            \label{eq:posterior-deprior}
        \end{split}
    \end{equation}
    Without $\pi_i^{(0)}$, multiplying packet posteriors generally counts the same domain
    prior more than once.

    \subsection{Three relationship models}

    The relationship among packets must be declared rather than inferred from the desired
    answer.
    \begin{center}
        \small
        \begin{tabularx}{\linewidth}{@{}>{\raggedright\arraybackslash}p{0.19\linewidth}>{\raggedright\arraybackslash}p{0.29\linewidth}Y@{}}
\toprule
Model & Scientific claim & Appropriate use \\
\midrule
M1: shared fixed class
& Every packet was generated by one common quotient class $G$.
& Repeated trajectories from the same unchanged interface, schema, or logger configuration. \\
M2: exchangeable cohort
& Member classes are draws from an unknown cohort distribution $\theta$.
& A vendor, program, or artifact family that usually but not always reuses conventions. \\
M3: robust sharedness mixture
& The target shares the reference prototype with prior probability $\rho_0$ and otherwise follows a baseline prior.
& Product default for strong A/B references and a hazy C whose relatedness is plausible but not guaranteed. \\
\bottomrule
        \end{tabularx}
    \end{center}

    Under M1, conditionally independent reference packets $A$ and $B$ give
    \begin{equation}
        q_{AB}(c)
        =\frac{\pi_0(c)\exp\{\lambda_A(c)+\lambda_B(c)\}}
        {\sum_g\pi_0(g)\exp\{\lambda_A(g)+\lambda_B(g)\}},
        \label{eq:cohort-prototype}
    \end{equation}
    where $q_{AB}$ is a posterior over the \emph{shared quotient class}, not a vote over raw
    representatives. M2 replaces that single class by a categorical cohort distribution and
    is developed in Appendix~\ref{app:cohort}.

    For the product-first M3 model, let $\rho_{0,C}$ be the provenance-derived probability,
    set before looking at C's local residuals, that target C follows the A/B prototype. Its
    borrowed prior is
    \begin{equation}
        \pi_C^{\mathrm{borrow}}(c)
        =\rho_{0,C}q_{AB}(c)+(1-\rho_{0,C})\pi_0(c),
        \label{eq:cohort-mixture-prior}
    \end{equation}
    and C's fused class posterior is
    \begin{equation}
        q_C^{\mathrm{fused}}(c)
        =\frac{\exp\{\lambda_C(c)\}\pi_C^{\mathrm{borrow}}(c)}
        {\sum_g\exp\{\lambda_C(g)\}\pi_C^{\mathrm{borrow}}(g)}.
        \label{eq:cohort-fused-posterior}
    \end{equation}
    The model automatically exposes prior--target conflict. Define
    \begin{equation}
        \begin{aligned}
            m_{AB}(D_C)&=\sum_c q_{AB}(c)\exp\{\lambda_C(c)\},\\
            m_0(D_C)&=\sum_c \pi_0(c)\exp\{\lambda_C(c)\}.
        \end{aligned}
    \end{equation}
    Then the posterior probability that C belongs to the shared component is
    \begin{equation}
        \rho_C^{\mathrm{post}}
        =\frac{\rho_{0,C}m_{AB}(D_C)}
        {\rho_{0,C}m_{AB}(D_C)+(1-\rho_{0,C})m_0(D_C)}.
        \label{eq:cohort-posterior-sharedness}
    \end{equation}
    If C agrees with A/B, $\rho_C^{\mathrm{post}}$ rises; if C strongly conflicts, borrowing
    falls rather than forcing the cohort answer. This robust-mixture logic is closely related
    to dynamic borrowing, commensurate-prior, and robust meta-analytic-predictive methods for
    historical information \cite{ibrahim2000,hobbs2011,schmidli2014,yang2023}.

    The explanation is additive on a pairwise log-odds scale:
    \begin{equation}
        \log\frac{q_C^{\mathrm{fused}}(c)}{q_C^{\mathrm{fused}}(g)}
        =\underbrace{\lambda_C(c)-\lambda_C(g)}_{\text{C's local trajectory}}
        +\underbrace{\log\frac{\pi_C^{\mathrm{borrow}}(c)}
        {\pi_C^{\mathrm{borrow}}(g)}}_{\text{cohort/provenance contribution}}.
        \label{eq:cohort-log-odds-explanation}
    \end{equation}
    A product report can therefore say exactly how much of a conclusion came from C and how
    much came from the cohort model.

    \subsection{What can and cannot transfer}

    Trajectory evidence transfers only at the quotient-class level. If two decoders are
    members of the same exact class, additional unanchored trajectories cannot distinguish
    them. A/B may help select a particular representative for C only when their packets
    carry an external anchor---for example, a verified API, documented attitude direction,
    or known rate frame. The fusion service must therefore keep two channels separate:
    \begin{equation}
        \boxed{\begin{aligned}
                   &\text{class evidence from trajectories}\\[-1mm]
                   &\quad\text{is kept separate from representative evidence from metadata.}
        \end{aligned}}
        \label{eq:cohort-two-channels}
    \end{equation}

    Likewise, packet likelihoods may be multiplied only after a lineage audit. Overlapping
    time ranges, reused sensor samples, derived copies, or packets produced from one another
    violate conditional independence. The safe default is to reject the fusion, retain one
    packet per dependency cluster, or construct a joint likelihood from the underlying data.
    An optional fractional power is a declared discounting model, not a free confidence knob
    \cite{ibrahim2000}. When predicting C, the cohort prior must be built from A/B without C;
    leave-one-out construction prevents a packet from reinforcing itself through the cohort.

    \subsection{When the output may be called a probability}

    A normalized number is a posterior probability only when all of the following hold:
    \begin{enumerate}[nosep]
    \item every packet uses the same quotient-space definition and compatible class support;
    \item $\lambda_i$ is a calibrated likelihood or declared generalized-Bayes update;
    \item overlapping evidence and derivation lineage have been removed or modeled jointly;
    \item the relationship model and $\rho_{0,C}$ were declared from provenance rather than
    tuned to obtain a preferred target answer; and
    \item calibration is validated on held-out trajectories from the stated operating domain.
    \end{enumerate}
    Otherwise the service returns \emph{cohort-adjusted evidence weights} and abstains from a
    probability claim. Appendix~\ref{app:cohort} gives a temperature-calibration construction
    and recommends held-out logarithmic and Brier scores, which are proper scoring rules
    \cite{gneiting2007}.

    \paragraph{Illustrative A/B/C result.}
    Let the baseline prior over three quotient classes be uniform. Suppose A and B have
    calibrated likelihood ratios $(9,1,0.5)$ and $(6,1,0.5)$, so
    $q_{AB}=(0.977,0.018,0.005)$. If $\rho_{0,C}=0.8$ and hazy C contributes only
    $(1.4,1,0.8)$, then C alone would have posterior
    $(0.438,0.313,0.250)$, while the robust fused posterior is
    \begin{equation}
        q_C^{\mathrm{fused}}=(0.896,0.061,0.042),
        \qquad \rho_C^{\mathrm{post}}=0.839.
        \label{eq:cohort-worked-result}
    \end{equation}
    The report should say that C's own evidence is weak and most certainty comes from the
    A/B cohort plus the stated $0.8$ sharedness prior. If C instead strongly favors another
    class, the same model lowers $\rho_C^{\mathrm{post}}$ and can abstain rather than conceal
    the conflict.

    The cohort result object is
    \begin{equation}
        \boxed{\begin{aligned}
                   \mathcal F_C=\bigl(&q_C^{\mathrm{local}},q_{AB},\rho_{0,C},
                   \rho_C^{\mathrm{post}},q_C^{\mathrm{fused}},\\
                   &\mathcal T_C^{\mathrm{fused}},\ell_{\mathrm{status}},
                   \mathsf{explain}_C\bigr).
        \end{aligned}}
        \label{eq:cohort-result-object}
    \end{equation}
    with one of four software statuses: \texttt{posterior}, \texttt{weight-only},
    \texttt{conflict}, or \texttt{abstain}, each namespaced by the cohort service. The
    original local result $\mathcal R_C$ remains unchanged and is always reported alongside
    $\mathcal F_C$.

    \begin{sectionrecap}
        Cohort fusion is a second-stage evidence model, not a sixth local decoder module. The
        local engine exports prior-separated quotient-class evidence and lineage. The cohort
        service declares whether members share one class, are exchangeable, or borrow through a
        robust sharedness mixture. It reports local and borrowed contributions separately and
        abstains when calibration, compatibility, or independence is not justified.

        \noindent\textbf{Key equations.}
        The packet contract is \eqref{eq:cohort-packet}; local priors are removed through
        \eqref{eq:posterior-deprior}; the robust borrowed prior and target posterior are
        \eqref{eq:cohort-mixture-prior}--\eqref{eq:cohort-fused-posterior}; and
        \eqref{eq:cohort-posterior-sharedness} reports target--cohort compatibility.

        \noindent\textbf{Symbols introduced.}
        \begin{center}
            \small
            \begin{tabularx}{0.97\linewidth}{@{}p{0.25\linewidth}Y@{}}
\toprule
Symbol & Meaning \\
\midrule
$\mathcal E_i,v_{\mathcal Q}$ & versioned evidence packet and quotient-space identifier \\
$\lambda_i(C),\pi_i^{(0)}(C)$ & prior-separated relative log evidence and local prior \\
$\mathsf{lin}_i,\mathsf{cal}_i$ & atomic lineage/dependency and calibration descriptors \\
$q_{AB}$ & prototype class posterior inferred from reference packets A and B \\
$\rho_{0,C},\rho_C^{\mathrm{post}}$ & prior and posterior probability that C follows the reference prototype \\
$\pi_C^{\mathrm{borrow}},q_C^{\mathrm{fused}}$ & robust cohort-derived prior and target posterior \\
$\mathcal F_C$ & auditable cohort-fusion result reported beside the unchanged local result \\
\bottomrule
            \end{tabularx}
        \end{center}
    \end{sectionrecap}


    \section{Conclusion}\label{sec:conclusion}

    Orientation-convention recovery does not require a mysterious optimizer. Its correct
    baseline is still the intuitive one: enumerate plausible decoders, convert them to one
    semantic representation, integrate the decoded angular rates, and compare the resulting
    relative rotations. The paper's contribution is to make the surrounding decision logic
    explicit enough that an implementation can be both aggressive and honest.

    The five-module organization separates the questions that had previously been mixed
    together. Normalization defines what is actually being compared. Quotient analysis
    removes exact duplicates and states which ambiguities no trajectory can resolve.
    Excitation diagnostics explain when a particular maneuver fails to exercise enough axes.
    Wrapping, Euler ranges, and unit cues prioritize plausible representations but include an
    abstention state. Domain knowledge becomes a hierarchical prior whose strength depends
    on whether the evidence is an exact API, an ecosystem, an artifact role, or only a broad
    realm. Finally, Lie-group integration and multiscale residuals provide the actual
    trajectory evidence, with explicit exact and statistical conditions for stopping work on
    a losing class.

    This organization also clarifies what the caller receives. The engine should return the
    surviving quotient classes and the reason they survived. A single decoder may be chosen
    for software use, but its reason label must distinguish trajectory identification from
    metadata, domain prior, heuristic preference, or canonical policy. That distinction is
    not cosmetic: it is the line between what the measurements established and what the
    system assumed in order to proceed.


    The product boundary extends the same discipline across multiple trajectories. Each
    local run exports prior-separated quotient-class evidence, calibration state, and source
    lineage. A cohort service can then use strong packets to form a prototype or partially
    pooled prior for a weak target, while a robust sharedness component discounts borrowing
    under target--cohort conflict. The local result is never overwritten: the service reports
    C's own evidence, the borrowed contribution, posterior sharedness, and the fused result
    side by side. Exact within-class ambiguities remain exact unless a packet contributes an
    external representative anchor.


    \appendix


    \section{Detailed data-availability regimes}\label{app:availability}
    The main text uses Table~\ref{tab:availability-modes} as the front-end decision contract.
    This appendix gives the detailed examples and identifiability arguments for time-origin,
    time-scale, latency, attitude-only, and rate-only cases.

    \subsection{Worked availability examples}

    The following examples show why ``time present'' and ``rates present'' are not binary
    labels.
    \begin{enumerate}[nosep]
    \item \textbf{Stamped attitude and gyro data in physical units.}  Quaternion or Euler
    attitudes, gyroscope values in a known unit, and trustworthy seconds give the
    standard cross-sensor problem. The residual tests both direction and magnitude,
    subject to the exact gauges cataloged later.
    \item \textbf{Hardware ticks and uncalibrated rate counts.}  If
    $t=s_t\tau+t_0$ and $\omega=\kappa Cy$, the record determines only
    $\lambda=s_t\kappa$. It cannot decide whether a numerical mismatch belongs to the
    tick duration or to the rate conversion.
    \item \textbf{No stamps, but a fixed-rate logger.}  A single unknown period $\Delta$
    is a defensible nuisance parameter. If the rate unit is anchored, $\Delta$ can
    often be estimated from attitude/rate consistency; if both are unknown, only
    $\Delta\kappa$ is observable.
    \item \textbf{Irregular event order with no durations.}  Giving every interval its
    own free $\Delta_k$ absorbs rate magnitude and leaves mainly directional
    compatibility. Bounds, a smooth clock model, or a low-dimensional time warp are
    needed to avoid fitting the clock to every residual.
    \item \textbf{Attitude and timestamps, but no rate stream.}  One may compare decoder
    validity, intrinsic path length, or acceleration regularity. These are
    source/motion priors; they are not an independent kinematic witness and cannot
    break every attitude-space isometry.
    \item \textbf{Rates and time, but no attitude stream.}  A known initial attitude alone
    only initializes each candidate curve. A known terminal attitude creates an
    endpoint residual; one observed inertial vector generally leaves a rotation about
    that vector unresolved at an instant, while two noncollinear vector observations
    determine an attitude. These added observations convert the problem into a
    genuine output-matching problem.
    \end{enumerate}

    \subsection{Time origin, time scale, and order-only data}

    A common shift of every timestamp has no effect on any increment.
    \begin{proposition}[Common time-origin gauge]
If $t_k^\dagger=t_k+b$ for fixed $b\in\R$, every interval, interpolated rate, integrated
increment, and residual in this paper is unchanged. A trajectory-consistency objective
cannot identify an absolute time origin. Only a \emph{relative} offset between two
streams can enter the residual.
    \end{proposition}

    When raw timestamps have an unknown scale, the sample-period gauge becomes an
    affine-clock gauge.
    \begin{theorem}[Affine clock--rate-scale gauge]
Let physical time and decoded rate be
\begin{equation}
    t=s_t\tau+t_0,
    \qquad
    \omega(t)=\kappa C y(\tau),
    \qquad s_t,\kappa>0.
\end{equation}
Every group integral depends on $s_t$ and $\kappa$ only through
\begin{equation}
    \lambda=s_t\kappa,
    \qquad
    \int \omega(t)\,dt
    =s_t\kappa\int C y(\tau)\,d\tau.
    \label{eq:affine-time-rate-product}
\end{equation}
Consequently,
\begin{equation}
    (s_t,\kappa)\longmapsto(cs_t,\kappa/c),
    \qquad c>0,
\end{equation}
leaves all predictions and residuals unchanged.
    \end{theorem}

    Thus an implementation should optimize the identifiable product $\lambda$, not two
    perfectly collinear parameters. A known timestamp unit, known gyro unit, clock-rate
    calibration, or other dimensional anchor is needed to split the product.

    If every interval is assigned an independent unknown duration, rate magnitude can be
    absorbed almost completely. In the small-step approximation
    $\phi_k\approx\Delta_k\bar\omega_k$, and for $\bar\omega_k\ne0$, the best positive
    duration is
    \begin{equation}
        \widehat\Delta_k
        =\frac{\max(0,\phi_k^{\T}\bar\omega_k)}{\norm{\bar\omega_k}^2}.
    \end{equation}
    When $\bar\omega_k=0$, the duration is unidentifiable to first order and the residual is
    $\norm{\phi_k}^2$. For $\bar\omega_k\ne0$, the minimized first-order residual is
    \begin{equation}
        \min_{\Delta_k\ge0}\norm{\phi_k-\Delta_k\bar\omega_k}^2
        =\norm{\phi_k}^2
        -\frac{\max(0,\phi_k^{\T}\bar\omega_k)^2}
        {\norm{\bar\omega_k}^2}.
        \label{eq:order-only-direction-residual}
    \end{equation}
    Equation~\eqref{eq:order-only-direction-residual} retains directional information but
    removes most magnitude information. In the full noncommutative model, direction changes
    within and across intervals still constrain the candidate, but one free $\Delta_k$ per
    interval is usually an over-flexible nuisance model. A defensible order-only model must
    therefore impose a clock prior, bounds, smoothness, or a low-dimensional time warp.

    \subsection{Unknown relative latency}

    Let the rate stream be evaluated at $t+\delta_t$ relative to attitude time. For a fixed
    quotient class $C$, the natural profile objective is
    \begin{equation}
        \widehat\delta_{t,C}
        \in\argmin_{\delta_t\in[\delta_{t,\min},\delta_{t,\max}]}J_C(\delta_t).
        \label{eq:latency-profile}
    \end{equation}
    Its local sensitivity is controlled by
    \begin{equation}
        \frac{\partial}{\partial\delta_t}\omega(t+\delta_t)
        =\dot\omega(t+\delta_t).
    \end{equation}
    Constant or nearly constant angular velocity therefore provides little latency
    information, while periodic motion produces repeated or near-repeated modes separated by
    the motion period. Temporal-calibration literature reaches the same practical
    conclusion: time offset is estimable under suitable motion and becomes weak or
    unobservable in degenerate regimes \cite{liMourikis2014}. A coarse lag search should
    report every competitive mode rather than only the local optimizer reached from one
    initial value.

    \subsection{Attitude without rates}

    With no independent rate stream, every attitude decoder produces a candidate curve
    $R_h(t)$. Smoothness can still be used as a \emph{model prior}, for example through
    \begin{equation}
        S_h=\Phi\!\left(
                      \set{d_g(R_{h,i},R_{h,j})}_{i,j},
                      \set{t_k}
        \right),
        \label{eq:attitude-only-functional}
    \end{equation}
    but $S_h$ is not a kinematic cross-check against another sensor.

    \begin{proposition}[Attitude-only isometry obstruction]
Let $F:\SO\rightarrow\SO$ be an isometry of $d_g$. Every score of the form
\eqref{eq:attitude-only-functional} satisfies
\begin{equation}
    S(F\circ R)=S(R).
\end{equation}
The same holds for intrinsic path length and norms of intrinsic velocity or covariant
acceleration because Riemannian isometries preserve those quantities.
    \end{proposition}

    This proposition explains why timestamps do not rescue scalar-first versus scalar-last
    quaternion recovery when the only criterion is intrinsic attitude regularity: the fixed
    orthogonal quaternion-coordinate maps discussed later induce isometries of the
    orientation manifold. Euler sequence candidates, by contrast, need not be related by an
    isometry when the same raw columns are fed to different decoders, so a smoothness prior
    may rank them. That ranking is conditional on the prior and should be labeled
    \texttt{attitude-heuristic}, not \texttt{trajectory-confirmed}.

    \subsection{Rates without attitude}

    \begin{theorem}[Rate-only self-integration is non-identifying]
Let $y(t)$ be a piecewise continuous raw rate record. For every admissible rate decoder
$g=(C,\kappa,f)$ and every initial attitude $R_0\in\SO$, the corresponding body- or
world-rate differential equation has a unique attitude solution $R_g(t)$. If no attitude,
endpoint orientation, known vector observation, commanded-motion model, or other output
is observed, a score based only on whether the decoded rate integrates to a valid
attitude curve is identical for all $g$.
    \end{theorem}

    For body rates, $Q R_g(t)$ is also a solution with the same rate for every constant world
    rotation $Q$; for world rates, $R_g(t)Q$ is also a solution for every constant body-basis
    rotation. These are initial-orientation gauges. Rate-only convention recovery becomes
    possible only after adding an external target, such as a known initial and final
    orientation,
    \begin{equation}
        d_g\!\left(R_T^{\mathrm{known}},R_g(T;R_0^{\mathrm{known}})\right),
    \end{equation}
    a gravity or magnetic vector observation, a known maneuver axis, actuator commands and a
    dynamics model, or calibrated magnitude bounds. Those additions are new measurements or
    priors; they are not contained in the rate record itself.

    \begin{sectionrecap}
        The availability regime determines the identifiable object before any optimizer is run.
        Known time plus attitude and rates gives the strongest problem. Unknown time scale leaves
        only a rate--clock product; order-only data preserve some directional compatibility but
        not physical speed. Attitude-only scoring is prior-based. Rate-only self-integration has
        no observed target and therefore cannot rank decoders without added measurements.

        \noindent\textbf{Key equations.}
        The affine gauge is \(\lambda=s_t\kappa\). With independently fitted interval duration,
        \(\widehat\Delta_k=[\phi_k^{\mathsf T}\bar\omega_k]_+/\|\bar\omega_k\|^2\), so the
        remaining first-order loss is the component of \(\phi_k\) orthogonal to the positive ray
        spanned by \(\bar\omega_k\). Rate-only recovery becomes a test only after an external
        term such as \(d_g(R_T^{\mathrm{known}},R_g(T;R_0^{\mathrm{known}}))\) is supplied.

        \noindent\textbf{Symbols introduced.}
        \par\noindent\begin{tabularx}{\linewidth}{@{}p{0.22\linewidth}Y@{}}
\toprule
Symbol & Meaning \\
\midrule
$\mathsf A,\mathsf G$ & availability of attitude and angular-rate streams \\
$\tau,s_t,t_0$ & raw time coordinate, positive clock scale, and common time origin \\
$\lambda=s_t\kappa$ & identifiable clock-scale/rate-scale product \\
$\Delta_k$ & physical duration of interval $k$ when known or modeled \\
$\delta_t$ & relative rate-to-attitude latency \\
$S_h$ & attitude-only regularity functional, not a cross-sensor residual \\
\bottomrule
        \end{tabularx}
    \end{sectionrecap}


    \section{Structural equivalences and identifiability proofs}\label{app:identifiability}

    \begin{definition}[Structural and data-dependent equivalence]
Two hypotheses are structurally equivalent for a stated scoring procedure if they
produce the same residual angles for every admissible data set. They are equivalent on
$\mathcal D$ if they produce the same residual angles on the observed data only. The
first is a property of the model; the second can arise from insufficient excitation or
an accidental trajectory symmetry.
    \end{definition}

    Convention recovery is well-posed only on the quotient of the hypothesis set by these
    equivalences. Let
    \begin{equation}
        \mathcal Q=\mathcal H/{\sim_s},
        \qquad
        q:\mathcal H\rightarrow\mathcal Q
    \end{equation}
    denote the structural quotient and its projection. Exact members of one class must not
    be ordered as though the trajectory distinguished them.

    \begin{definition}[Canonical section]
For metadata $z$ and a declared source profile $p$, a canonicalizer is a section
\begin{equation}
    s_{p,z}:\mathcal Q\rightarrow\mathcal H,
    \qquad
    q\!\left(s_{p,z}(C)\right)=C
    \quad\text{for every }C\in\mathcal Q.
    \label{eq:canonical-section}
\end{equation}
It selects one representation of each class without changing the class itself.
    \end{definition}

    This is the same logical operation as fixing a datum in an inverse problem with gauge
    freedom: it removes representational nonuniqueness for computation but does not add
    measurement information \cite{triggs2000}. The section may depend on documented metadata,
    interoperability profiles, or a published fallback policy. It must not depend on
    hypothesis enumeration order, hash order, thread scheduling, or floating-point noise.
    The full quotient class remains the inferential result.

    \subsection{Attitude-direction, rate-frame, and sign duality}

    \begin{theorem}[Exact direction--frame--sign duality]\label{thm:direction-frame-sign}
        Consider any decoded attitude sequence $R_k$, decoded rate sequence $\omega_k$, and the
        Lie-group integrator above. Define the dual sequence by
        \begin{equation}
            R_k^\dagger=R_k^{\T},
            \qquad
            \omega_k^\dagger=-\omega_k,
            \qquad
            B^\dagger=W,
            \quad W^\dagger=B.
        \end{equation}
        For every interval and every window length,
        \begin{equation}
            \rho_{h^\dagger,k,\ell}=\rho_{h,k,\ell}.
        \end{equation}
        Therefore any score constructed only from these geodesic residual angles assigns
        exactly the same value to $h$ and $h^\dagger$.
    \end{theorem}

    \begin{proof}
        For body kinematics, differentiating $R^{-1}$ gives
        \begin{equation}
            \frac{d}{dt}R^{-1}=-[\omega_B]_\times R^{-1},
        \end{equation}
        which is world-form kinematics for $R^{-1}$ with rate $-\omega_B$. The world-to-body
        case is analogous.

        The discrete statement is exact as well. If
        $P_B=A_0A_1\cdots A_{M-1}$, then switching to world composition and negating the rate
        gives
        \begin{equation}
            P_W^\dagger=A_{M-1}^{-1}\cdots A_1^{-1}A_0^{-1}=P_B^{-1}.
        \end{equation}
        The dual observed increment is also the inverse:
        \begin{equation}
            \Delta_W^\dagger
            =R_{k+1}^{\T} R_k
            =(R_k^{\T} R_{k+1})^{-1}=\Delta_B^{-1}.
        \end{equation}
        The same identities hold for composed windows. Since inversion is an isometry of the
        bi-invariant metric,
        $d_g(P^{-1},\Delta^{-1})=d_g(P,\Delta)$.
    \end{proof}

    \begin{corollary}
        If the search simultaneously includes both attitude directions, both rate frames, and
        both global rate signs, and if it contains the dual of every candidate, trajectory
        consistency alone cannot return a unique member of each dual pair. At least one of those
        three dimensions must be anchored externally, or the output must report the pair as one
        equivalence class.
    \end{corollary}

    This is stronger than a low-excitation warning: it is an exact symmetry for every
    trajectory, including noiseless, richly excited motion.

    \subsection{Rate-scale and sample-period gauge}

    \begin{theorem}[Unknown-time scale gauge]\label{thm:time-scale-gauge}
        Assume timestamps are absent and one constant sample period $\Delta$ is fitted. Let the
        rate decoder include a positive scale $\kappa$. For any $c>0$, the transformation
        \begin{equation}
            \kappa^\dagger=c\kappa,
            \qquad
            \Delta^\dagger=\frac{\Delta}{c}
        \end{equation}
        leaves every predicted substep matrix and therefore every residual unchanged.
    \end{theorem}

    \begin{proof}
        Every substep exponent contains the product
        \begin{equation}
            \frac{\Delta}{M}\,\kappa C y.
        \end{equation}
        The transformed product is
        $(\Delta/c)(c\kappa)Cy/M$, which is identical.
    \end{proof}

    \begin{corollary}
        Without timestamps or another time-scale anchor, radians-per-second versus
        degrees-per-second cannot be identified separately from the sample period. Bounds or
        priors can select one member of the class, but that selection is external information.
    \end{corollary}

    \subsection{Quaternion storage order from orientation alone}

    Let $S^3$ be the unit quaternions and identify antipodal points $q\sim-q$. The quotient
    $S^3/\{\pm1\}$ is $\mathbb{RP}^3$, which is the orientation manifold $\SO$ with its
    standard metric up to a constant scale.

    \begin{theorem}[Orthogonal quaternion-coordinate invariance]\label{thm:quaternion-o4}
        Let $q_i,q_j\in S^3$ and let $Q\in\OO(4)$ be any fixed orthogonal coordinate
        transformation. Then
        \begin{equation}
            d_g(R(q_i),R(q_j))
            =d_g(R(Qq_i),R(Qq_j)).
        \end{equation}
        Consequently, every orientation-only objective built exclusively from intrinsic
        Riemannian distances is invariant under a fixed quaternion-coordinate permutation,
        including scalar-first versus scalar-last storage.
    \end{theorem}

    \begin{proof}
        For unit quaternions,
        \begin{equation}
            d_g(R(q_i),R(q_j))
            =2\arccos\!\left(\abs{q_i^{\T} q_j}\right).
        \end{equation}
        Orthogonality gives
        $(Qq_i)^{\T}(Qq_j)=q_i^{\T}q_j$. Moreover, $Q(-q)=-Qq$, so $Q$ descends to an
        isometry of the antipodal quotient.
    \end{proof}

    The theorem applies to orientation-only intrinsic path length and intrinsic smoothness,
    not only to one-step angle differences. A rate signal, known reference orientation,
    wire-format specification, or another physical observable can break the ambiguity
    because a generic $Q\in\OO(4)$ is not a quaternion group automorphism and does not
    preserve the attitude-rate kinematic relation.

    \subsection{When body and world rates coincide numerically}

    \begin{proposition}[Coaxial body/world ambiguity]
Let $R(t)$ be differentiable and let the same numerical vector $u(t)$ be proposed as
both body- and world-expressed angular velocity. Both kinematic equations hold if and
only if
\begin{equation}
    R(t)u(t)=u(t)
\end{equation}
for every $t$ at which $u(t)\ne0$.
    \end{proposition}

    \begin{proof}
        Both equations hold exactly when
        $R[u]_\times=[u]_\times R$. Multiplying on the right by $R^{\T}$ gives
        $R[u]_\times R^{\T}=[u]_\times$. Since
        $R[u]_\times R^{\T}=[Ru]_\times$ and the hat map is injective, this is equivalent to
        $Ru=u$.
    \end{proof}

    Pure or nearly pure rotation about one persistent axis therefore creates an exact or
    near body/world ambiguity when the numerical axis is fixed by the attitude.

    \subsection{Coordinate relabeling gauge}

    For fixed proper basis changes $Q_W,Q_B\in\SO$, the transformed quantities
    \begin{equation}
        R'(t)=Q_W^{\T} R(t)Q_B,
        \qquad
        \omega_W'(t)=Q_W^{\T}\omega_W(t),
        \qquad
        \omega_B'(t)=Q_B^{\T}\omega_B(t)
    \end{equation}
    preserve the kinematic equations. Body increments are conjugated by $Q_B$ and world
    increments by $Q_W$, so all bi-invariant residual angles are unchanged. If a hypothesis
    space allows coordinated relabeling of both attitude axes and rate axes without an
    external axis reference, those labels form another gauge class. Improper basis changes
    must use the axial-vector rule stated earlier.

    \subsection{Euler reversal and branch equivalence}

    Equation~\eqref{eq:intrinsic-extrinsic-reversal} is a structural duplicate, not a
    low-excitation coincidence. If the hypothesis space independently enumerates angle
    channel permutations, the pair
    \begin{equation}
        (\mathrm{intrinsic}\ abc,\;\alpha,\beta,\gamma)
    \quad\longleftrightarrow\quad
    (\mathrm{extrinsic}\ cba,\;\gamma,\beta,\alpha)
    \end{equation}
    returns exactly the same rotation at every sample and therefore the same trajectory
    residual under every rate model. Likewise, decoders that differ only by the exact branch
    maps in \eqref{eq:tait-bryan-branch} or \eqref{eq:proper-euler-branch} are the same
    attitude decoder after canonicalization. Enumerating them as separate statistical models
    artificially multiplies prior mass and computational cost.

    \subsection{Exact equivalences that survive perfect data and full-rank motion}

    Call a trajectory full-rank when an appropriate observed-increment or rate Gramian is
    positive definite. Full rank helps with empirical axis excitation, but it cannot break a
    structural symmetry. Table~\ref{tab:exact-equivalences} lists exact pairs and gauges that
    survive noiseless, arbitrarily long, full-rank data whenever the stated information is
    absent.

    \begingroup
    \small
    \setlength{\tabcolsep}{3.5pt}
    \begin{longtable}{@{}p{0.25\textwidth}p{0.30\textwidth}p{0.35\textwidth}@{}}
        \caption{Exact observational equivalences.}\label{tab:exact-equivalences}\\
        \toprule
        Transformation or gauge & Why the score is identical & Information that can break it \\
        \midrule
        \endfirsthead
        \toprule
        Transformation or gauge & Why the score is identical & Information that can break it \\
        \midrule
        \endhead
        \bottomrule
        \endlastfoot
        Quaternion sign $q_k\mapsto-q_k$, independently by sample
        & Both quaternions represent the same rotation exactly.
        & Nothing is needed for attitude; sign continuity is only a storage canonicalization. \\
        Direction--frame--sign dual of Theorem~\ref{thm:direction-frame-sign}
        & Observed and predicted increments are both inverted, and $d_g$ is inversion-invariant.
        & One documented member of attitude direction, rate frame, or rate sign. \\
        Clock/rate scale $(s_t,\kappa)\mapsto(cs_t,\kappa/c)$ or
        $(\Delta,\kappa)\mapsto(\Delta/c,c\kappa)$
        & Every exponential depends only on the product of time scale and rate scale.
        & Known timestamp unit, known rate unit, or another dimensional calibration. \\
        Common timestamp origin $t\mapsto t+b$
        & Only time differences and relative alignment enter.
        & An external epoch, which is semantically useful but irrelevant to the kinematic score. \\
        Intrinsic $abc$ with $(\alpha,\beta,\gamma)$ versus extrinsic $cba$ with
        $(\gamma,\beta,\alpha)$
        & The elemental rotation product is identical sample by sample.
        & Metadata that distinguishes moving-axis from fixed-axis language or fixes channel order. \\
        Euler branch transformations in \eqref{eq:tait-bryan-branch} and
        \eqref{eq:proper-euler-branch}
        & Multiple triples are coordinates of the same rotation.
        & A documented output branch; physical orientation data cannot distinguish coordinate representatives. \\
        Coordinated basis relabeling $(Q_W,Q_B)$
        & Attitudes and rates transform covariantly, while bi-invariant residual angles are unchanged.
        & Gravity, magnetic north, a surveyed axis, vehicle geometry, or another external vector reference. \\
        Fixed $Q\in\OO(4)$ quaternion-coordinate map in attitude-only mode
        & It is an isometry of $S^3/\{\pm1\}$ and preserves all intrinsic curve regularity.
        & A rate stream, known reference orientation, vector observation, or wire-format metadata. \\
        Rate-only decoder and initial-attitude gauges
        & Every decoded rate integrates to a valid curve, and no attitude output is observed.
        & Endpoint attitude, vector observations, commanded dynamics, or calibrated magnitude/axis references. \\
    \end{longtable}
    \endgroup

    The following phenomena are different: they are ties for a particular trajectory or
    objective, not universal symmetries.
    \begin{center}
        \small
        \par\noindent\begin{tabularx}{\linewidth}{@{}p{0.31\linewidth}YY@{}}
\toprule
Data-dependent degeneracy & Typical cause & Useful remedy \\
\midrule
Body/world tie without the global dual & Coaxial or nearly coaxial rotation & Noncommuting multi-axis excitation \\
Signed-permutation near ties & One or more weakly exercised axes & Maneuvers spanning all axes with unequal temporal structure \\
Latency nonuniqueness & Constant or periodic angular velocity & Aperiodic changes in rate direction and magnitude \\
Sample-period alias modes & Large rotations or closed loops between samples & Higher-rate sampling or credible intersample bounds \\
Euler ambiguity & Motion near a chart singularity & Switch representation or add motion away from the singular set \\
Aggregate-score tie under full-rank samples & Symmetric experiment design or trimming away discriminating intervals & Pairwise residual inspection and asymmetric validation maneuvers \\
\bottomrule
        \end{tabularx}
    \end{center}

    A report should explicitly classify every retained pair as \emph{structural},
    \emph{availability-induced}, or \emph{trajectory-dependent}. The phrase ``insufficient
    excitation'' is incorrect for the first two categories.


    \begin{sectionrecap}
        Identifiability belongs to the quotient $\mathcal H/\!\sim_s$, not to the raw enum list.
        Structural gauges survive perfect, arbitrarily long, full-rank data; data-dependent
        near-ties may be reduced by better maneuvers. A canonical section chooses one software
        representative but does not turn a gauge choice into measurement evidence.

        \noindent\textbf{Key equations.}
        \(\mathcal Q=\mathcal H/\!\sim_s\) and
        \(q(s_{p,z}(C))=C\). The exact global dual is
        \((R,\omega,B/W)\mapsto(R^{\mathsf T},-\omega,W/B)\). Missing time scale yields
        \((\Delta,\kappa)\mapsto(\Delta/c,c\kappa)\), and attitude-only quaternion geometry is
        invariant under every fixed $Q\in\mathrm O(4)$.

        \noindent\textbf{Symbols introduced.}
        \par\noindent\begin{tabularx}{\linewidth}{@{}p{0.24\linewidth}Y@{}}
\toprule
Symbol & Meaning \\
\midrule
$\sim_s$ & structural observational equivalence relation \\
$\mathcal Q,q$ & quotient-class set and projection from hypotheses \\
$s_{p,z}$ & metadata/profile-dependent canonical section \\
$h^\dagger$ & exact direction--frame--sign dual of hypothesis $h$ \\
$Q_W,Q_B$ & coordinated world- and body-basis relabelings \\
$C\in\mathcal Q$ & one structural equivalence class \\
\bottomrule
        \end{tabularx}
    \end{sectionrecap}


    \section{Aliasing, wrapping, and the cut locus}\label{app:aliasing}

    The exponential map is one-to-one on the open ball
    $\set{\phi\in\R^3:\norm{\phi}<\pi}$. Endpoint attitude samples alone cannot reveal
    how many closed rotational loops occurred between samples. A sufficient, model-based
    no-wrap diagnostic for interval $k$ is
    \begin{equation}
        L_{h,k}
        =\int_{t_k}^{t_{k+1}}\norm{\omega_h(t)}\,dt
        \le
        \frac{\Delta t_k}{2}
        \left(\norm{\omega_{h,k}}+\norm{\omega_{h,k+1}}\right)
        <\pi-\varepsilon_{\mathrm{alias}},
    \end{equation}
    where the inequality follows from convexity of the norm under linear interpolation.
    If the bound exceeds $\pi$, aliasing is not proved, but it cannot be excluded by the
    sampled endpoints. Intervals whose observed relative angle is close to $\pi$ should also
    be flagged because the logarithm axis is unstable there.

    No algorithm can recover an unsampled closed loop from endpoint attitudes without a
    credible intersample model or additional measurements. High-rate tests should therefore
    include trajectories that differ by full loops but share the same sampled attitudes.

    \begin{sectionrecap}
        Sampled endpoint attitudes determine only a principal relative rotation, not the number
        of unobserved full turns between samples. A credible intersample model can certify a
        no-wrap regime; exceeding its bound means aliasing is possible, not that it definitely
        occurred.

        \noindent\textbf{Key equation.}
        For linearly interpolated candidate rates, a sufficient diagnostic is
        \(L_{h,k}\le(\Delta t_k/2)(\|\omega_{h,k}\|+\|\omega_{h,k+1}\|)
        <\pi-\varepsilon_{\mathrm{alias}}\).

        \noindent\textbf{Symbols introduced.}
        \par\noindent\begin{tabularx}{\linewidth}{@{}p{0.23\linewidth}Y@{}}
\toprule
Symbol & Meaning \\
\midrule
$L_{h,k}$ & upper bound on integrated angular path length in interval $k$ \\
$\varepsilon_{\mathrm{alias}}$ & configured margin from the $\pi$ no-wrap boundary \\
$\pi$ cut locus & relative rotations whose principal logarithm axis is nonunique \\
\bottomrule
        \end{tabularx}
    \end{sectionrecap}


    \section{Excitation and empirical distinguishability diagnostics}\label{app:excitation}

    Let $\mathcal V\subseteq\{0,\ldots,N-1\}$ denote intervals retained after numerical
    validity and cut-locus screening. Excitation diagnostics answer a conditioning question:
    \emph{did the observed maneuver exercise enough independent directions for the finite
    search to separate convention choices reliably?}  They do not replace structural
    quotient analysis.

    \subsection{Decoder-invariant rate-channel excitation}

    When a three-channel rate stream is present and the base decoder family consists of a
    common scalar unit conversion and signed permutations,
    \begin{equation}
        \omega_{h,k}=\kappa_h C_h y_k,
        \qquad C_h\in\mathcal B_3,
    \end{equation}
    a particularly useful diagnostic can be computed directly from the raw channels. For
    positive weights $w_k^{\rm exc}$, define
    \begin{equation}
        G_y=\sum_{k\in\mathcal V}w_k^{\rm exc}y_k y_k^{\mathsf T}.
        \label{eq:raw-rate-gramian}
    \end{equation}
    Use $w_k^{\rm exc}=\Delta t_k$ when physical durations are known and unit weights when
    only sample order is available. The decoded Gramian is
    \begin{equation}
        G_{\omega,h}
        =\sum_{k\in\mathcal V}w_k^{\rm exc}\omega_{h,k}\omega_{h,k}^{\mathsf T}
        =\kappa_h^2 C_hG_yC_h^{\mathsf T}.
        \label{eq:decoded-rate-gramian}
    \end{equation}
    Because signed-permutation matrices are orthogonal, $G_{\omega,h}$ and $G_y$ have the
    same eigenvalue ratios. Thus basic excitation quality is invariant to the unknown
    common rate scale, channel order, and channel signs in this decoder family.

    Let
    \begin{equation}
        \lambda_1\ge\lambda_2\ge\lambda_3\ge0
    \end{equation}
    be the eigenvalues of $G_y$. For a declared numerical threshold
    $\varepsilon_{\rm rank}\in(0,1)$, report the thresholded rank and weakest-axis ratio
    \begin{align}
        \operatorname{rank}_{\varepsilon_{\rm rank}}(G_y)
        &=\sum_{i=1}^{3}\mathbf 1\{\lambda_i\ge
        \varepsilon_{\rm rank}\lambda_1\},\\
        \eta_{\rm exc}(G_y)&=\frac{\lambda_3}{\lambda_1},
        \qquad (\lambda_1>0).
        \label{eq:excitation-rank-ratio}
    \end{align}
    A threshold-free companion is the participation-ratio effective rank
    \begin{equation}
        r_{\rm eff}(G_y)
        =\frac{\operatorname{tr}(G_y)^2}{\operatorname{tr}(G_y^2)}\in[1,3],
        \qquad G_y\ne0.
        \label{eq:excitation-effective-rank}
    \end{equation}
    Roughly, $\operatorname{rank}_{\varepsilon_{\rm rank}}=1$ signals nearly coaxial motion, $\operatorname{rank}_{\varepsilon_{\rm rank}}=2$ signals
    one weakly exercised channel direction, and $\operatorname{rank}_{\varepsilon_{\rm rank}}=3$ indicates that all three
    raw channel directions exceed the configured threshold. The threshold must be reported;
    ``rank three'' without a scale-relative tolerance is not reproducible.

    Full numerical rank is still not sufficient for convention identifiability. Structural
    gauges survive every excitation pattern, and symmetric or nearly isotropic maneuvers can
    make distinct signed permutations or timing hypotheses difficult to separate even when
    $\lambda_3>0$. The Gramian is therefore a conditioning diagnostic, not a model-selection
    criterion.

    \subsection{Attitude-only fallback}

    When rates are unavailable, an analogous motion summary can be constructed after an
    attitude decoder $a$ has produced observed increment vectors
    \begin{equation}
        \phi_{a,k}=\Log(\Delta R_{a,k}),
        \qquad
        G^{R}_a=\sum_{k\in\mathcal V}w_k^{\rm exc}
        \phi_{a,k}\phi_{a,k}^{\mathsf T}.
        \label{eq:attitude-increment-gramian}
    \end{equation}
    The same eigenvalue diagnostics may be reported, but unlike $G_y$ this quantity is
    candidate dependent: different attitude decoders can produce different curves. It is
    therefore useful for describing the motion under a candidate, not as independent evidence
    that the candidate is correct.

    \subsection{Empirical class separation}

    After exact symmetries have been quotiented, the most direct distinguishability quantity
    is the class-risk gap beyond the empirical minimizer set. Define
    \begin{equation}
        J_{\min}=\min_{C\in\mathcal Q_0}J_C,
        \qquad
        \widehat{\mathcal C}_{\min}
        =\set{C\in\mathcal Q_0:J_C=J_{\min}}.
    \end{equation}
    When at least one class lies outside that set, report
    \begin{equation}
        \widehat\gamma
        =\min_{C\in\mathcal Q_0\setminus\widehat{\mathcal C}_{\min}}
        \left(J_C-J_{\min}\right),
        \label{eq:empirical-class-gap}
    \end{equation}
    and set $\widehat\gamma=+\infty$ if every class is an empirical minimizer. The size of
    $\widehat{\mathcal C}_{\min}$ must be reported separately: a positive next-class gap does
    not erase a tie among several minimizers. Also report the full near-tie set and the
    per-scale contributions. Pairwise score-gap
    matrices under representative synthetic maneuvers are often more informative than one
    scalar excitation statistic because they reveal exactly which convention dimensions are
    confused by a maneuver.

    \begin{sectionrecap}
        Rate-channel excitation can be diagnosed before convention decoding. For signed
        permutations with one common scale, $G_{\omega,h}=\kappa_h^2C_hG_yC_h^{\mathsf T}$, so its
        eigenvalue ratios are decoder invariant. A thresholded rank below three warns that the
        record lacks well-conditioned three-axis excitation; rank three is useful but not
        sufficient for identifiability. With no rates, an attitude-increment Gramian provides a
        candidate-dependent fallback. Final empirical distinguishability is measured by gaps
        between quotient-class risks, not by excitation rank alone.

        \noindent\textbf{Key equations.}
        \(G_y=\sum_k w_k^{\rm exc}y_ky_k^{\mathsf T}\),
        \(G_{\omega,h}=\kappa_h^2C_hG_yC_h^{\mathsf T}\),
        \(\operatorname{rank}_{\varepsilon_{\rm rank}}=\sum_i\mathbf1\{\lambda_i\ge\varepsilon_{\rm rank}\lambda_1\}\),
        \(\eta_{\rm exc}=\lambda_3/\lambda_1\), and
        \(\widehat\gamma=\min_{C\notin\widehat{\mathcal C}_{\min}}(J_C-J_{\min})\).

        \noindent\textbf{Symbols introduced.}
        \par\noindent\begin{tabularx}{\linewidth}{@{}p{0.25\linewidth}Y@{}}
\toprule
Symbol & Meaning \\
\midrule
$G_y$ & raw angular-rate channel excitation Gramian \\
$G_{\omega,h}$ & decoded rate Gramian under hypothesis representative $h$ \\
$w_k^{\rm exc}$ & positive excitation weight for interval/sample $k$ \\
$\lambda_1,\lambda_2,\lambda_3$ & ordered eigenvalues of $G_y$ \\
$\operatorname{rank}_{\varepsilon_{\rm rank}}$ & scale-relative thresholded numerical rank \\
$\eta_{\rm exc}$ & weakest-axis eigenvalue ratio $\lambda_3/\lambda_1$ \\
$r_{\rm eff}$ & participation-ratio effective rank \\
$G^R_a$ & candidate-dependent attitude-increment Gramian used when rates are absent \\
$\widehat{\mathcal C}_{\min},\widehat\gamma$ & empirical minimizer set and gap from it to the next class \\
$\mathcal V$ & intervals admitted to excitation and timing diagnostics \\
\bottomrule
        \end{tabularx}
    \end{sectionrecap}


    \section{Timing, time scale, latency, and time warp}\label{app:timing}

    For compactness, this appendix writes $h$ for the stored representative $h_C$ of a
    quotient class $C$ while profiling continuous timing parameters. After profiling, the
    resulting geometric risk is attached to the class as $J_C$.

    For a candidate frame, let
    \begin{equation}
        \phi_{h,k}=\Log(\Delta R_{h,k}^{\mathrm{obs}}),
        \qquad
        \bar\omega_{h,k}=\frac{\omega_{h,k}+\omega_{h,k+1}}{2},
    \end{equation}
    with both vectors expressed in the coordinates associated with that frame. Under a
    small-step, no-alias approximation,
    \begin{equation}
        \phi_{h,k}\approx \Delta\,\bar\omega_{h,k}.
    \end{equation}
    For linearly interpolated rates, the first noncommutative correction is exposed by the
    Magnus expansion:
    \begin{align}
        \phi_{B,k}
        &=\Delta\,\bar\omega_k
        +\frac{\Delta^2}{12}(\omega_k\times\omega_{k+1})
        +O(\Delta^3\Omega_k^3),\\
        \phi_{W,k}
        &=\Delta\,\bar\omega_k
        -\frac{\Delta^2}{12}(\omega_k\times\omega_{k+1})
        +O(\Delta^3\Omega_k^3),
    \end{align}
    where $\Omega_k=\max(\norm{\omega_k},\norm{\omega_{k+1}})$ and the remainder notation
    assumes $\Delta\Omega_k$ is small. The sign difference follows from right- versus
    left-ordered evolution.

    A weighted first-order least-squares initializer is
    \begin{equation}
        \widetilde\Delta_h
        =\frac{\sum_{k\in\mathcal V}w_k^{\rm ls}\bar\omega_{h,k}^{\T}\phi_{h,k}}
        {\sum_{k\in\mathcal V}w_k^{\rm ls}\norm{\bar\omega_{h,k}}^2},
        \label{eq:dt-initializer}
    \end{equation}
    where $w_k^{\rm ls}\ge0$ are optional reliability weights and $\mathcal V$ excludes intervals
    with a nearly zero denominator contribution or a logarithm near the cut locus. Equation
    \eqref{eq:dt-initializer} is a useful initializer, not a final estimator when rotations
    are large or noncommuting.

    The mathematically consistent estimate is the bounded one-dimensional group fit
    \begin{equation}
        \widehat\Delta_h
        \in\argmin_{\Delta\in[\Delta_{\min},\Delta_{\max}]}J_h(\Delta),
    \end{equation}
    initialized by \eqref{eq:dt-initializer}. Because aliasing can make $J_h(\Delta)$
    multimodal, a coarse global grid followed by local refinement is safer than one local
    optimizer. All near-optimal modes and active bounds should be reported. The unknown-time
    scale theorem implies that rate-unit hypotheses must be quotient-grouped unless the
    period bounds or another external time reference break the scale gauge.

    If one candidate fits a continuous nuisance parameter while another does not, raw
    in-sample residuals are not directly comparable without a complexity correction or
    held-out evaluation. In the common case where every candidate fits the same one scalar
    period, parameter counts are equal, but validation is still preferable because the
    period is data-selected.

    \subsection{Known irregular timestamps}

    When physical timestamps are known, no constant-period parameter is needed. Each
    interval uses its own $\Delta t_k$, and nonuniform sampling is not itself a problem. The
    interpolation model must, however, be defined on physical time rather than sample index.
    A candidate that silently treats irregular data as uniformly sampled changes the forward
    model and can bias both convention and latency estimates.

    If timestamps are known only in an unknown unit, set $t=s_t\tau+t_0$ and optimize the single
    identifiable product $\lambda=s_t\kappa$ from
    \eqref{eq:affine-time-rate-product}. Do not report separate confidence intervals for $s_t$
    and $\kappa$ unless an external anchor is introduced.

    \subsection{Regularized time warp when only order is known}

    A low-dimensional or regularized time warp can preserve some convention information
    without granting one arbitrary scale per interval. One example is
    \begin{equation}
        \min_{h,\Delta_0,\ldots,\Delta_{N-1}>0}
        J_h(\Delta_{0:N-1})
        +\lambda_t\sum_{k=0}^{N-2}
        \left(\log\Delta_{k+1}-\log\Delta_k\right)^2,
        \label{eq:regularized-time-warp}
    \end{equation}
    subject to average-period or endpoint-duration bounds. The penalty is a clock-smoothness
    prior. Its strength must be selected outside the final validation blocks, and the report
    must show how the winning class changes with $\lambda_t$. A spline clock with a small
    number of coefficients is another option. An unconstrained monotone warp is generally
    too flexible for convention recovery.

    \subsection{Joint scale and latency search}

    For asynchronous streams with an unknown affine rate clock, the identifiable timing
    parameters are a relative offset $\delta_t$ and a scale-rate product $\lambda$. The profile
    problem is
    \begin{equation}
        (\widehat\lambda_h,\widehat\delta_{t,h})
    \in\argmin_{\lambda\in[\lambda_{\min},\lambda_{\max}],\,
    \delta_t\in[\delta_{t,\min},\delta_{t,\max}]}
    J_h(\lambda,\delta_t).
    \label{eq:joint-scale-latency}
    \end{equation}
    A practical solver first uses a coarse two-dimensional grid or a cross-correlation proxy
    to locate modes, then refines each competitive basin with the full group objective. The
    profile curvature near a mode should be reported; a sharp numerical optimum can still be
    semantically gauge-coupled to an unmodeled unit choice.

    Constant rate makes $\delta_t$ unobservable, and periodic rate can generate repeated modes.
    A linearly changing rate may identify latency in principle but can be strongly correlated
    with gyro bias or interpolation assumptions. Temporal calibration should therefore be
    benchmarked jointly with bias, scale, and trajectory family rather than in isolation.

    \begin{sectionrecap}
        Known irregular intervals should be used directly. With unknown clock units, optimize
        the single identifiable product $\lambda=s_t\kappa$. Order-only data require a constrained
        time model; one free duration per interval removes too much magnitude information.
        Relative latency is a profile parameter whose observability comes from time variation and
        whose objective may be multimodal under periodic motion.

        \noindent\textbf{Key equations.}
        The first-order initializer is
        \(\widetilde\Delta_h=(\sum w_k^{\rm ls}\bar\omega_{h,k}^{\mathsf T}\phi_{h,k})/
        (\sum w_k^{\rm ls}\|\bar\omega_{h,k}\|^2)\). Final fitting uses
        \(\widehat\Delta_h\in\arg\min J_h(\Delta)\) or the joint profile
        \((\widehat\lambda_h,\widehat\delta_{t,h})\in\arg\min J_h(\lambda,\delta_t)\).

        \noindent\textbf{Symbols introduced.}
        \par\noindent\begin{tabularx}{\linewidth}{@{}p{0.24\linewidth}Y@{}}
\toprule
Symbol & Meaning \\
\midrule
$\phi_{h,k}$ & principal observed-increment rotation vector \\
$\bar\omega_{h,k}$ & interval-average decoded rate used by the initializer \\
$\widetilde\Delta_h$ & first-order constant-period initializer \\
$\mathcal V,w_k^{\rm ls}$ & valid interval set and optional least-squares reliability weights \\
$\lambda_t$ & regularization strength for a smooth positive time warp \\
$J_h(\lambda,\delta_t)$ & full group risk profiled over time/rate scale and latency \\
\bottomrule
        \end{tabularx}
    \end{sectionrecap}


    \section{Statistical interpretation and finite-sample guarantees}\label{app:statistics}

    \subsection{A loss is not automatically a likelihood}

    The robust multi-window risk in \eqref{eq:robust-risk} is a model-selection loss. Its
    trimming, overlapping windows, and multiple scales make it deliberately unlike a simple
    independent-sample likelihood. A transformed score such as
    $100\exp(-J_C/(2s_J^2))$ is only a monotone display mapping when the display scale
    $s_J$ is fixed independently of $C$; it is not a posterior probability.

    A globally defined probabilistic residual model is the isotropic matrix-Fisher family
    \cite{lee2018}
    \begin{equation}
        p(E\mid\kappa_F)
        =c(\kappa_F)^{-1}\exp\!\left(\kappa_F\tr(E)\right),
        \qquad E\in\SO.
    \end{equation}
    With a common concentration $\kappa_F$, maximum likelihood ranks quotient classes by
    $\sum_k\tr(E_{C,k})$, equivalently by the chordal loss
    \begin{equation}
        3-\tr(E_{C,k})
        =4\sin^2\!\left(\frac{\rho_{C,k}}{2}\right)
        =\rho_{C,k}^2+O(\rho_{C,k}^4).
    \end{equation}
    This model is globally defined and avoids choosing a logarithm axis at $\pi$. It still
    does not make temporally overlapping residuals independent.

    For small errors away from the cut locus, one may instead use the tangent-space
    quasi-model
    \begin{equation}
        e_{C,k}\approx\mathcal N(0,\Sigma).
    \end{equation}
    If $\Sigma=\sigma^2\I$ is fixed independently of the class, the corresponding
    quasi-log-likelihood differs between quotient classes only through
    $\sum_k\norm{e_{C,k}}^2/(2\sigma^2)$. An estimated covariance should be reported with
    its eigenvalues; forcing isotropy when residuals are strongly anisotropic can distort
    evidence units.

    \subsection{Adjacent residuals are structurally correlated}

    Even white attitude-sample noise produces correlated increment residuals because
    adjacent increments share a sample.

    \begin{proposition}[First-order residual correlation from attitude noise]
Suppose the true body-frame increment is $P_k=R_k^{\T} R_{k+1}$ and measured attitudes
have independent right-multiplicative small errors
\begin{equation}
    R_k^m=R_k\Exp(\eta_k),
    \qquad
    \mathbb E\eta_k=0,
    \qquad
    \operatorname{Cov}(\eta_k)=\Sigma_R.
\end{equation}
If the predicted increment is exact, then the right-invariant residual satisfies
\begin{equation}
    e_k
    =-\Ad_{P_k^{-1}}\eta_k+\eta_{k+1}+O(\norm{\eta}^2),
    \qquad
    \Ad_Qv=Qv.
\end{equation}
Hence $e_k$ and $e_{k+1}$ are generally correlated even though the $\eta_k$ are
independent.
    \end{proposition}

    \begin{proof}
        The measured increment is
        \begin{equation}
            (R_k^m)^{\T} R_{k+1}^m
        =\Exp(-\eta_k)P_k\Exp(\eta_{k+1}).
        \end{equation}
        Premultiplying by $P_k^{-1}$, conjugating the first exponential through $P_k$, and using
        the first-order Baker--Campbell--Hausdorff formula gives the result.
    \end{proof}

    Endpoint rate interpolation likewise shares rate samples across adjacent intervals.
    Therefore an independent Gaussian likelihood is, at best, a quasi-likelihood. Options
    for stronger inference include a block covariance model, heteroskedasticity-and-
    autocorrelation-consistent estimates, nonoverlapping validation blocks, or a moving-block
    bootstrap.

    \subsection{Finite-class recovery under explicit assumptions}

    The following theorem states a limited but honest recovery guarantee. It applies to an
    untrimmed, bounded validation loss on independent blocks, not automatically to the full
    robust overlapping-window score.

    \begin{theorem}[Finite-class risk recovery]
Let $\mathcal Q$ be a finite set of structural quotient classes. For independent and
identically distributed validation blocks $j=1,\ldots,n$, let
\begin{equation}
    X_{C,j}=\frac{\rho_{C,j}^2}{\pi^2}\in[0,1],
    \qquad
    L(C)=\mathbb E X_{C,j},
    \qquad
    \widehat L_n(C)=\frac1n\sum_{j=1}^{n}X_{C,j}.
\end{equation}
Let
\begin{equation}
    \mathcal Q_\star=\argmin_{C\in\mathcal Q}L(C)
\end{equation}
be the population-risk minimizer set, and suppose the post-quotient gap
\begin{equation}
    \gamma
    =\min_{C\notin\mathcal Q_\star}
    \left[L(C)-\min_{C'\in\mathcal Q}L(C')\right]
    >0.
\end{equation}
Then any empirical minimizer
$\widehat C_n\in\argmin_{C\in\mathcal Q}\widehat L_n(C)$ satisfies
\begin{equation}
    \Pr(\widehat C_n\notin\mathcal Q_\star)
    \le
    \bigl(\abs{\mathcal Q}-\abs{\mathcal Q_\star}\bigr)
    \exp\!\left(-\frac{n\gamma^2}{2}\right).
    \label{eq:finite-Q-bound}
\end{equation}
    \end{theorem}

    \begin{proof}
        Fix $C_\star\in\mathcal Q_\star$ and $C'\notin\mathcal Q_\star$. The differences
        $D_j=X_{C',j}-X_{C_\star,j}$ lie in $[-1,1]$ and have mean at least $\gamma$. Hoeffding's
        inequality \cite{hoeffding1963} gives
        \begin{equation}
            \Pr\!\left(\widehat L_n(C')\le\widehat L_n(C_\star)\right)
            \le\exp(-n\gamma^2/2).
        \end{equation}
        A union bound over quotient classes outside $\mathcal Q_\star$ proves
        \eqref{eq:finite-Q-bound}.
    \end{proof}

    The theorem makes the central identifiability condition explicit: a positive population
    risk gap after quotienting exact symmetries. More motion is not sufficient when
    $\gamma=0$.

    An anytime-valid, conservative prefix rule follows from the same bounded loss. For a
    target error probability $\delta\in(0,1)$, define
    \begin{equation}
        r_n(\delta)
        =\sqrt{\frac{1}{2n}
        \log\!\left(\frac{\pi^2\abs{\mathcal Q}n^2}{3\delta}\right)}.
    \end{equation}
    By Hoeffding's inequality \cite{hoeffding1963} and a union bound over all quotient classes and all
    $n\ge1$, with probability at least $1-\delta$,
    \begin{equation}
        \abs{\widehat L_n(C)-L(C)}\le r_n(\delta)
    \end{equation}
    for every $C\in\mathcal Q$ and every prefix length simultaneously. Therefore, if the empirical best
    class at prefix $n$ is separated from every other class by more than $2r_n(\delta)$,
    then it is also the unique population-risk class on that event. This is a statistical
    statement about the specified validation loss and assumptions, not a proof that the
    forward model contains the physical truth.

    The same uniform event yields a set-valued finite-sample result instead of forcing a
    winner. Let $\widehat C_n\in\argmin_{C\in\mathcal Q}\widehat L_n(C)$ and define
    \begin{equation}
        \mathcal A_n(\delta)
        =\set{C\in\mathcal Q:
        \widehat L_n(C)-\widehat L_n(\widehat C_n)
            \le 2r_n(\delta)}.
        \label{eq:ambiguity-set}
    \end{equation}
    \begin{proposition}[Coverage and exclusion of the ambiguity set]
On the simultaneous event
$\sup_{C\in\mathcal Q}\abs{\widehat L_n(C)-L(C)}\le r_n(\delta)$,
$\mathcal A_n(\delta)$ contains every population-risk minimizer. Every class whose
population gap exceeds $4r_n(\delta)$ is excluded from $\mathcal A_n(\delta)$.
    \end{proposition}
    \begin{proof}
        If $C_\star$ minimizes $L$, then
        $\widehat L_n(C_\star)\le L_\star+r_n$ while
        $\widehat L_n(\widehat C_n)\ge L_\star-r_n$, so the empirical difference is at most
        $2r_n$. Conversely, if $L(C)-L_\star>4r_n$, then
        $\widehat L_n(C)>L_\star+3r_n$ whereas
        $\widehat L_n(\widehat C_n)\le\widehat L_n(C_\star)\le L_\star+r_n$; hence the
        empirical difference exceeds $2r_n$.
    \end{proof}
    The construction is intentionally conservative. It plays the same operational role as
    a model confidence set---uninformative data retain many candidates, while informative
    data shrink the set---but it follows directly from the bounded-loss uniform event above,
    not from the bootstrap elimination algorithm of Hansen, Lunde, and Nason
    \cite{hansen2011}.

    To make independence plausible, use disjoint blocks that do not share attitude or rate
    samples. Under temporal dependence, replace the formula with a justified mixing bound
    or a block bootstrap. A fixed threshold such as a plug-in log-likelihood ratio of five
    has no universal interpretation after winner selection, variance estimation, trimming,
    and serial correlation.


    \begin{sectionrecap}
        The robust engineering risk is not automatically a likelihood. A probabilistic model
        must state its residual distribution and temporal dependence. Finite-sample recovery
        requires a positive post-quotient population gap; otherwise the mathematically correct
        output is a set of unresolved classes.

        \noindent\textbf{Key equations.}
        The matrix-Fisher chordal loss is
        \(3-\operatorname{tr}(E)=4\sin^2(\rho/2)\). Shared attitude samples give
        \(e_k=-\operatorname{Ad}_{P_k^{-1}}\eta_k+\eta_{k+1}+O(\|\eta\|^2)\). Under independent
        bounded validation blocks,
        \[
            \Pr(\widehat C_n\notin\mathcal Q_\star)
            \le(|\mathcal Q|-|\mathcal Q_\star|)e^{-n\gamma^2/2},
        \]
        and $\mathcal A_n(\delta)$ retains classes within the simultaneous risk radius.

        \noindent\textbf{Symbols introduced.}
        \par\noindent\begin{tabularx}{\linewidth}{@{}p{0.24\linewidth}Y@{}}
\toprule
Symbol & Meaning \\
\midrule
$s_J$ & fixed display-scale parameter for monotone score visualization \\
$\kappa_F$ & matrix-Fisher concentration parameter in the probabilistic subsection \\
$\Sigma,\Sigma_R$ & tangent residual and attitude-sample covariance matrices \\
$X_{C,j},L(C),\widehat L_n(C)$ & bounded block loss, population risk, and empirical risk of a quotient class \\
$\gamma$ & positive population-risk gap after quotienting exact ties \\
$r_n(\delta)$ & simultaneous Hoeffding radius over quotient classes and prefixes \\
$\mathcal A_n(\delta)$ & evidence-preserving finite-sample ambiguity set \\
\bottomrule
        \end{tabularx}
    \end{sectionrecap}


    \section{Hierarchical cohort evidence fusion and probability calibration}
    \label{app:cohort}

    This appendix formalizes the follow-on service introduced in
    Section~\ref{sec:cohort-main}. The local orientation engine and the cohort service solve
    different inverse problems. The local engine estimates evidence over a versioned set of
    quotient classes. The cohort service specifies a probabilistic relationship among those
    class variables and combines compatible packets. Keeping those layers separate prevents
    source priors, repeated evidence, and representative conventions from being counted as
    new trajectory measurements.

    \subsection{The packet is a likelihood interface, not a winner label}

    Let $\mathcal Q_v=\{C_1,\ldots,C_K\}$ be a fixed quotient universe. For packet $i$,
    write its local data as $D_i$ and define
    \begin{equation}
        L_i(c)=p(D_i\mid C_i=c),
        \qquad
        \lambda_i(c)=\log L_i(c)-\max_g\log L_i(g).
        \label{eq:cohort-local-likelihood}
    \end{equation}
    Only likelihood ratios matter for finite-class updating, so the max-zero representation is
    numerically stable. A posterior from the local engine is
    \begin{equation}
        q_i(c)
        =\frac{\pi_i^{(0)}(c)e^{\lambda_i(c)}}
        {\sum_g\pi_i^{(0)}(g)e^{\lambda_i(g)}}.
        \label{eq:cohort-local-posterior}
    \end{equation}
    Equations~\eqref{eq:posterior-deprior} and \eqref{eq:cohort-local-posterior} show why the
    packet must retain $\pi_i^{(0)}$: a downstream updater can reconstruct relative evidence
    from $(q_i,\pi_i^{(0)})$, but not from $q_i$ alone.

    A production packet should contain at least the fields below.
    \begin{center}
        \small
        \begin{tabularx}{\linewidth}{@{}p{0.26\linewidth}Y@{}}
\toprule
Field & Required semantics \\
\midrule
Packet and quotient version & Stable packet ID, schema version, $v_{\mathcal Q}$, ordered class IDs, and representative-map version. \\
Class evidence & $\lambda_i(c)$, score semantics, local prior, local posterior if available, and the evidence contender set. \\
Calibration & Calibration ID, training/validation domain, effective-sample convention, temperature or likelihood model, and validity status. \\
Observability & Excitation rank/eigenvalues, conditioning, window counts, trimming policy, and warnings about weak axes or aliasing. \\
Hard anchors & Documented units, schema, frame direction, rate frame, API identity, and any representative-level facts. \\
Provenance & Domain, artifact role, ecosystem/protocol, exact interface guess, software/version fingerprints, and similarity features. \\
Lineage & Atomic measurement IDs, sensor/time ranges, derivation parents, dependency clusters, and packet-generation code version. \\
Audit & All local warnings, rejected classes with reasons, and the immutable local result $\mathcal R_i$. \\
\bottomrule
        \end{tabularx}
    \end{center}

    The class likelihood is constant across exact representatives $h\in C$. Representative
    anchors therefore belong in a separate conditional channel,
    \begin{equation}
        A_i(h\mid C)=p(M_i\mid h,C),
        \qquad h\in C,
        \label{eq:cohort-representative-channel}
    \end{equation}
    where $M_i$ is external metadata or provenance. No number of unanchored trajectory
    packets can make $A_i$ informative about an exact direction--frame--sign dual.

    \subsection{M1: one fixed class shared by the cohort}

    Suppose all members of a reference set $S$ were generated by one latent class $G$:
    \begin{equation}
        G\sim\operatorname{Categorical}(\pi_0),
        \qquad D_i\perp D_j\mid G,
        \qquad D_i\mid G=c\sim L_i(c).
        \label{eq:cohort-m1-model}
    \end{equation}
    Then
    \begin{equation}
        p(G=c\mid D_S)
        =\frac{\pi_0(c)\exp\{\sum_{i\in S}\lambda_i(c)\}}
        {\sum_g\pi_0(g)\exp\{\sum_{i\in S}\lambda_i(g)\}}.
        \label{eq:cohort-m1-posterior}
    \end{equation}
    For any two classes $c$ and $g$, posterior log odds decompose as
    \begin{equation}
        \log\frac{p(G=c\mid D_S)}{p(G=g\mid D_S)}
        =\log\frac{\pi_0(c)}{\pi_0(g)}
        +\sum_{i\in S}\log\frac{L_i(c)}{L_i(g)}.
        \label{eq:cohort-m1-log-odds}
    \end{equation}
    Thus independent packets contribute additive log Bayes factors. M1 is appropriate only
    when ``same convention'' is part of the source model, such as repeated files from one
    unchanged schema. It is not justified merely because the top local labels happen to
    agree.

    A discounted form
    \begin{equation}
        p_a(G=c\mid D_S)
        \propto\pi_0(c)\exp\left\{\sum_{i\in S}a_i\lambda_i(c)\right\},
        \qquad 0\le a_i\le1,
        \label{eq:cohort-power-posterior}
    \end{equation}
    is a categorical analogue of a power-prior update \cite{ibrahim2000}. The $a_i$ values
    must be predeclared or calibrated from operating characteristics. They are not a license
    to lower the weight of an inconvenient packet after inspecting the desired class.

    \subsection{M2: an exchangeable cohort with partial pooling}

    When members usually but not always reuse conventions, let
    \begin{equation}
        \theta\sim\operatorname{Dirichlet}(\alpha_0\pi_0),
        \qquad
        C_i\mid\theta\sim\operatorname{Categorical}(\theta),
        \qquad
        D_i\mid C_i=c\sim L_i(c).
        \label{eq:cohort-m2-model}
    \end{equation}
    The packet labels remain latent. Integrating each $C_i$ gives
    \begin{equation}
        p(\theta\mid D_S)
        \propto
        \operatorname{Dirichlet}(\theta;\alpha_0\pi_0)
        \prod_{i\in S}\left(\sum_{c=1}^{K}\theta_c L_i(c)\right).
        \label{eq:cohort-m2-posterior}
    \end{equation}
    The posterior-predictive prior for a new target is
    \begin{equation}
        \widetilde\pi_S(c)
        =\mathbb E[\theta_c\mid D_S],
        \qquad
        p(C_T=c\mid D_T,D_S)
        \propto L_T(c)\widetilde\pi_S(c).
        \label{eq:cohort-m2-predictive}
    \end{equation}
    Unlike hard majority voting, this model lets an uncertain A or B contribute soft evidence.
    The posterior is no longer Dirichlet-conjugate because the labels are uncertain, but the
    state space is finite and the expectation can be computed by MCMC, quadrature for small
    $K$, or a checked variational approximation. A fast fractional-count approximation,
    \begin{equation}
        \widetilde\pi_S(c)
        \approx
        \frac{\alpha_0\pi_0(c)+\sum_{i\in S}w_i q_i(c)}
        {\alpha_0+\sum_{i\in S}w_i},
        \label{eq:cohort-fractional-count}
    \end{equation}
    is useful for initialization and dashboards, but should be labeled an approximation
    unless derived from a specified fractional-label model.

    \subsection{M3: robust sharedness borrowing for one hazy target}

    The robust mixture in the main text has a simple latent-variable interpretation. Let
    $S_T\in\{0,1\}$ indicate whether target T follows the reference prototype:
    \begin{equation}
        \begin{aligned}
            S_T\mid z_T,z_S&\sim\operatorname{Bernoulli}(\rho_{0,T}),\\
            C_T\mid S_T=1,D_S&\sim\operatorname{Categorical}(q_S),\\
            C_T\mid S_T=0&\sim\operatorname{Categorical}(\pi_0),\\
            D_T\mid C_T=c&\sim L_T(c).
        \end{aligned}
        \label{eq:cohort-m3-model}
    \end{equation}
    Here $q_S$ is the M1 prototype posterior or the M2 posterior-predictive distribution.
    Marginalizing $S_T$ yields \eqref{eq:cohort-mixture-prior}; Bayes' rule yields
    \eqref{eq:cohort-fused-posterior} and \eqref{eq:cohort-posterior-sharedness}.

    The prior sharedness $\rho_{0,T}$ should be derived from information available before the
    target residual is used, for example
    \begin{equation}
        \operatorname{logit}(\rho_{0,T})
        =b_0+b^{\mathsf T}s(z_T,z_S),
        \label{eq:cohort-sharedness-regression}
    \end{equation}
    where $s$ contains schema equality, vendor/program identity, file-role similarity,
    software-version proximity, and other provenance features. This classifier requires its
    own held-out calibration. Target data then update sharedness through
    \eqref{eq:cohort-posterior-sharedness}; they should not also be reused to hand-tune
    $\rho_{0,T}$. Commensurate priors and robust mixture priors were developed precisely to
    reduce borrowing when current and historical information conflict
    \cite{hobbs2011,schmidli2014}; self-adapting mixtures provide a related data-conflict
    mechanism \cite{yang2023}. These works motivate the architecture, while the finite
    categorical model above is specific to convention quotient classes.

    \subsection{Dependency, lineage, and leave-one-out discipline}

    The product formula in \eqref{eq:cohort-m1-posterior} assumes conditional independence.
    Let $U_i$ be the set of atomic observations or derivation roots represented by packet $i$.
    A sufficient strict-mode check is
    \begin{equation}
        U_i\cap U_j=\varnothing
        \qquad\text{for every pair multiplied as independent evidence.}
        \label{eq:cohort-lineage-disjoint}
    \end{equation}
    A shared sensor file, overlapping time interval, copied orientation estimate, or packet
    constructed from another packet violates this condition. The preferred remedies are,
    in order:
    \begin{enumerate}[nosep]
    \item fuse the underlying data jointly and emit one packet;
    \item keep one representative packet per dependency cluster;
    \item construct a defensible cluster likelihood; or
    \item use a declared power/discount model and label the result as such.
    \end{enumerate}
    Simply multiplying both packets is not conservative; it can create arbitrarily high
    confidence from duplicated evidence.

    When predicting each member of an existing cohort, use leave-one-out borrowing:
    \begin{equation}
        \widetilde\pi_{-i}(c)
        =p(C_i=c\mid\{D_j:j\ne i\}),
        \qquad
        q_i^{\mathrm{fused}}(c)
        \propto L_i(c)\widetilde\pi_{-i}(c).
        \label{eq:cohort-loo}
    \end{equation}
    The target packet must not help construct the prior later used to claim that the target
    agrees with its cohort. After a leave-one-out report is produced, the target may be
    included in the cohort posterior for future packets.

    \subsection{Calibrating geometric risk into class probability}

    The robust risk $J_i(C)$ orders classes but is not automatically a log likelihood.
    One practical calibration starts with the nonnegative class gap
    \begin{equation}
        \Delta_i(C)
        =n_{\mathrm{eff},i}\,[J_i(C)-\min_GJ_i(G)]
        \label{eq:cohort-risk-gap}
    \end{equation}
    and a temperature family
    \begin{equation}
        q_i(C;T)
        =\frac{\pi_i^{(0)}(C)\exp\{-\Delta_i(C)/T\}}
        {\sum_G\pi_i^{(0)}(G)\exp\{-\Delta_i(G)/T\}},
        \qquad T>0.
        \label{eq:cohort-temperature-posterior}
    \end{equation}
    For a labeled calibration corpus $\mathcal C$ with true quotient classes
    $C_i^{\mathrm{true}}$, fit
    \begin{equation}
        T^\star
        \in\argmin_{T>0}
        \left[-\sum_{i\in\mathcal C}
            \log q_i(C_i^{\mathrm{true}};T)\right].
        \label{eq:cohort-temperature-fit}
    \end{equation}
    The effective count $n_{\mathrm{eff},i}$ must reflect disjoint validation blocks or a
    separately justified dependence correction; it is not the raw number of overlapping
    windows. Calibration should be assessed on an untouched test set with reliability
    plots, logarithmic score, Brier score, and stratification by domain, excitation, sample
    length, and missing-data regime. Logarithmic and Brier scores are strictly proper under
    their standard formulations \cite{gneiting2007}.

    The packet's calibration status should be one of:
    \begin{center}
        \small
        \begin{tabularx}{0.95\linewidth}{@{}p{0.26\linewidth}Y@{}}
\toprule
Status & Permitted claim \\
\midrule
\texttt{calibrated-likelihood} & Bayesian posterior under the stated residual model and cohort assumptions. \\
\texttt{calibrated-generalized} & Calibrated predictive class probability under the stated temperature/domain mapping. \\
\texttt{generalized-weight} & Normalized loss-based evidence weight; no literal probability claim. \\
\texttt{uncalibrated-risk} & Ranking and contender-set use only. \\
\bottomrule
        \end{tabularx}
    \end{center}
    Calibration is versioned. A software, sensor, or operating-domain change that moves
    performance outside the calibration envelope should invalidate the probability-ready
    flag until revalidated.

    \subsection{Worked A/B/C calculation and conflict behavior}

    Let $\mathcal Q=\{C_1,C_2,C_3\}$ and $\pi_0=(1/3,1/3,1/3)$. Use the likelihood-ratio
    vectors
    \begin{equation}
        L_A=(9,1,0.5),
        \qquad
        L_B=(6,1,0.5),
        \qquad
        L_C=(1.4,1,0.8).
    \end{equation}
    M1 applied to A and B gives
    \begin{equation}
        q_{AB}
        =\frac{(54,1,0.25)}{55.25}
        =(0.9774,0.0181,0.0045).
    \end{equation}
    With $\rho_{0,C}=0.8$,
    \begin{equation}
        \pi_C^{\mathrm{borrow}}
        =0.8q_{AB}+0.2\pi_0
        =(0.8486,0.0811,0.0703).
    \end{equation}
    C alone has
    \begin{equation}
        q_C^{\mathrm{local}}
        =\frac{(1.4,1,0.8)}{3.2}
        =(0.4375,0.3125,0.2500),
    \end{equation}
    whereas the fused posterior is
    \begin{equation}
        q_C^{\mathrm{fused}}
        =\operatorname{normalize}
        \bigl(L_C\odot\pi_C^{\mathrm{borrow}}\bigr)
        =(0.8963,0.0612,0.0424).
    \end{equation}
    The shared and baseline predictive ordinates are
    $m_{AB}(D_C)=1.3900$ and $m_0(D_C)=1.0667$, producing
    $\rho_C^{\mathrm{post}}=0.8390$.

    Now keep A/B fixed but let target evidence be
    $L_C^{\mathrm{conflict}}=(0.3,3,1)$. The result becomes
    \begin{equation}
        q_C^{\mathrm{fused}}
        =(0.4480,0.4284,0.1237),
        \qquad
        \rho_C^{\mathrm{post}}=0.4956.
    \end{equation}
    The model no longer produces a confident cohort assignment: C's strong disagreement
    nearly equalizes the first two classes and halves the probability of sharedness. A
    configured acceptance rule should abstain.

    \subsection{Decision contract and service decomposition}

    Let $q_{(1)}$ and $q_{(2)}$ be the two largest fused class probabilities. A simple
    probability-ready decision rule is
    \begin{equation}
        \text{accept class }\widehat C
        \quad\Longleftrightarrow\quad
        q_{(1)}\ge p_{\min},
        \quad q_{(1)}-q_{(2)}\ge\Delta_p,
        \quad \rho_T^{\mathrm{post}}\ge\rho_{\min},
        \label{eq:cohort-acceptance-rule}
    \end{equation}
    plus successful compatibility, calibration, and lineage checks. The thresholds are
    operating-policy parameters chosen from validation studies, not universal constants.
    Failure produces an abstention or conflict status and leaves the full fused contender set.

    A product implementation naturally separates six components:
    \begin{enumerate}[nosep]
    \item a local recovery engine that emits immutable evidence packets;
    \item a packet registry that validates schema and quotient-space versions;
    \item a lineage/dependency auditor;
    \item a cohort relationship resolver that selects M1, M2, or M3 from provenance;
    \item a class-level fusion core plus a separate representative-anchor resolver; and
    \item a calibration/monitoring service that tracks reliability and domain drift.
    \end{enumerate}
    The output must preserve $\mathcal R_T$, $\mathcal E_T$, the reference packet IDs, the
    relationship model, all priors and calibration versions, the local/cohort log-odds
    contributions, and the final $\mathcal F_T$. That audit trail is what makes the sentence
    ``C is probably class $C_1$ because A/B are strong'' reproducible rather than rhetorical.

    \begin{sectionrecap}
        Cohort borrowing is a finite latent-class problem over quotient classes. M1 multiplies
        independent evidence only when one common convention is a real source assumption. M2
        learns an exchangeable cohort distribution. M3 mixes a reference predictive
        probability with a baseline prior and reports posterior sharedness, providing a practical
        robust default for a hazy target. Probabilities require calibrated local evidence,
        version compatibility, lineage control, and held-out operating validation.

        \noindent\textbf{Key equations.}
        Local class evidence is \eqref{eq:cohort-local-likelihood}. M1 uses
        \eqref{eq:cohort-m1-posterior}; M2 uses \eqref{eq:cohort-m2-posterior} and its predictive
        \eqref{eq:cohort-m2-predictive}; M3 uses \eqref{eq:cohort-m3-model}. Risk calibration is
        \eqref{eq:cohort-temperature-posterior}--\eqref{eq:cohort-temperature-fit}, and the
        acceptance contract is \eqref{eq:cohort-acceptance-rule}.

        \noindent\textbf{Symbols introduced.}
        \par\noindent\begin{tabularx}{\linewidth}{@{}p{0.25\linewidth}Y@{}}
\toprule
Symbol & Meaning \\
\midrule
$D_i,L_i,\lambda_i$ & local packet data, class likelihood, and max-zero relative log evidence \\
$G,q_S$ & one shared cohort class and its reference posterior \\
$\theta,\alpha_0$ & exchangeable cohort class-frequency vector and Dirichlet concentration \\
$S_T,\rho_{0,T},\rho_T^{\mathrm{post}}$ & target sharedness indicator and its prior/posterior probabilities \\
$U_i$ & atomic lineage set represented by packet $i$ \\
$n_{\mathrm{eff},i},T^\star$ & effective independent evidence count and fitted calibration temperature \\
$p_{\min},\Delta_p,\rho_{\min}$ & validated operational acceptance thresholds \\
$\odot$ & componentwise multiplication \\
\bottomrule
        \end{tabularx}
    \end{sectionrecap}


    \section{Detailed unit, wrapping, and Euler-angle heuristics}\label{app:heuristics}

    Euler data require more than enumerating sequence names. The raw triple must be assigned
    to axes, angle units, intrinsic or extrinsic composition, an attitude direction, and a
    branch policy. The twelve classical three-axis sequences split into two families:
    \begin{itemize}[nosep]
    \item \emph{Tait--Bryan} sequences use three distinct axes; and
    \item \emph{proper Euler} sequences repeat the first axis as the third axis.
    \end{itemize}
    \begin{center}
        \small
        \begin{tabular}{@{}lll@{}}
            \toprule
            Family       & Six axis sequences        & Principal middle interval \\
            \midrule
            Tait--Bryan  & $XYZ,XZY,YXZ,YZX,ZXY,ZYX$ & $[-\pi/2,\pi/2]$          \\
            Proper Euler & $XYX,XZX,YXY,YZY,ZXZ,ZYZ$ & $[0,\pi]$                 \\
            \bottomrule
        \end{tabular}
    \end{center}
    The word ``pitch'' is a source-specific semantic label, not a generic name for the
    middle coordinate of every sequence. For a fixed sequence $s$ and composition
    convention, write the decoder as
    $E_s(\alpha,\beta,\gamma)\in\SO$, where $\beta$ is the middle angle.

    \subsection{A unit-and-wrap hypothesis classifier}

    Suppose one or more raw attitude channels appear angular but the angular unit is unknown.
    Let the candidate unit set be
    \begin{equation}
        \mathcal U_\theta=\{\mathrm{rad},\mathrm{deg},\mathrm{turn},\mathrm{grad}\},
    \end{equation}
    with one-turn periods and radian conversion factors
    \begin{equation}
        \begin{array}{c|cccc}
            u & \mathrm{rad} & \mathrm{deg} & \mathrm{turn} & \mathrm{grad}\\
            \hline
            T_u & 2\pi & 360 & 1 & 400\\
            s_u=2\pi/T_u & 1 & \pi/180 & 2\pi & \pi/200.
        \end{array}
        \label{eq:unit-periods}
    \end{equation}
    The period $T_u$ is expressed in the \emph{raw numerical unit}; $s_u x$ converts a raw
    angle to radians under hypothesis $u$.

    A visible boundary crossing is informative because a wrapped source can show a large raw
    jump even when the underlying motion is smooth. For one raw channel $x_k$, let
    $d_k=x_k-x_{k-1}$. For candidate period $T_u$, use the period-complement correction
    \begin{equation}
        m_{u,k}=\begin{cases}
                    \operatorname{round}(d_k/T_u), & \abs{d_k}>T_u/2,\\
                    0, & \text{otherwise},
        \end{cases}
        \qquad
        r_{u,k}=d_k-m_{u,k}T_u,
        \label{eq:wrap-correction}
    \end{equation}
    optionally clipping $m_{u,k}$ to $\{-1,0,1\}$ when the sampling model excludes more than
    one full revolution per sample. This is the same complement-of-large-deltas idea used
    by standard phase-unwrapping routines, generalized to an arbitrary period
    \cite{numpyUnwrap}.

    Let $\rho_H$ be a robust loss and let $S>0$ be a robust raw-difference scale computed
    once from the file, independent of $u$. Define the wrap gain and inferred-wrap count
    \begin{align}
        G_u &= \sum_k
        \left[
            \rho_H(d_k/S)-\rho_H(r_{u,k}/S)
        \right]_+,\\
        n_u &= \sum_k \mathbf 1\{m_{u,k}\ne0\}.
    \end{align}
    A period that explains a few otherwise extreme jumps receives large $G_u$; a period so
    small that it ``unwraps'' ordinary motion produces an implausibly large $n_u$. One
    transparent heuristic is therefore
    \begin{equation}
        H_{\mathrm{wrap}}(u)
        =-G_u+\lambda_{\mathrm{wrap}}n_u,
        \label{eq:wrap-heuristic}
    \end{equation}
    with the penalty fixed by a source profile or tuned on labeled data. The absolute value
    of $H_{\mathrm{wrap}}$ has no probabilistic meaning by itself.

    \paragraph{The important abstention case.}
    If the record contains no candidate boundary crossing, then $m_{u,k}=0$ for every
    plausible $u$ and the wrap feature supplies essentially no information. A stream staying
    inside approximately $[-\pi,\pi]$ is compatible with radians, but it is also compatible
    with a degree-valued source whose motion simply stayed in a narrow range. In that case
    the classifier should say \texttt{unit-undetermined-by-wrap}, not manufacture a confident
    unit choice from the absence of evidence.

    \paragraph{From hand-built heuristic to statistical classifier.}
    The representation heuristic and the domain prior should remain separate statistical
    objects. Let $v$ denote only a representation attribute that the current raw stream can
    plausibly inform---for example angle unit, Euler family, principal-versus-unwrapped
    behavior, or a branch policy. Construct a representation feature vector
    \begin{equation}
        \phi_{\rm repr}=\bigl(
        H_{\mathrm{wrap}},
        H_{\mathrm{range}},
        H_{\mathrm{branch}},
        f_{\mathrm{sing}},
        H_{\mathrm{physical}}
        \bigr),
        \label{eq:heuristic-features}
    \end{equation}
    using only components available in the current data regime. Exact schema facts have
    already acted as hard constraints in Module~1, source/domain evidence belongs to
    Module~4, and coarse trajectory loss belongs to Module~5. Keeping those quantities out
    of $\phi_{\rm repr}$ prevents the same information from being counted twice.

    With a labeled corpus, a calibrated multiclass model may estimate
    \begin{equation}
        q_{\theta}(v\mid\phi_{\rm repr})
        =\frac{\exp\!\left(b_v+\beta_v^{\T}\phi_{\rm repr}\right)}
        {\sum_{v'\in\mathcal V_{\rm repr}}
            \exp\!\left(b_{v'}+\beta_{v'}^{\T}\phi_{\rm repr}\right)},
        \label{eq:heuristic-softmax}
    \end{equation}
    followed by held-out calibration if the output is to be interpreted probabilistically.
    The classifier predicts the representation attribute $v$, not a full physical
    convention class that also contains rate frame, timing, or domain semantics. Each raw
    hypothesis inherits the score associated with its attribute value, and exact quotient
    members still contribute only one evidential class downstream.

    The engine should support abstention, for example when the largest calibrated attribute
    probability is below $p_{\min}$ or the top-two margin is below $\delta_p$. Without a
    representative labeled corpus, the same features may be used lexicographically or in a
    documented screening score, but normalized numbers must be called \emph{heuristic
    weights}, not posterior probabilities. A separate coarse-trajectory screen may combine
    $\phi_{\rm repr}$ with $\widetilde J_C$ for acceleration, but that combined screen is a
    Module~5 computational device and must be validated by quotient-class recall.

    \paragraph{Angle unit versus rate unit versus time unit.}
    Wrapping can suggest the period of an \emph{angle-valued channel}. It does not by itself
    identify radians/second versus degrees/second for a gyro, nor seconds versus milliseconds
    for a clock. If attitude, rates, and physical timestamps are all present, rate-unit
    hypotheses are best tested directly by trajectory consistency. If the timestamp scale
    and rate scale are both unknown, only their product is identifiable, as established by the time--rate scale gauge.
    Timestamp-unit hypotheses such as seconds, milliseconds, microseconds, or hardware ticks
    therefore belong to the timing hypothesis set and require a dimensional anchor or a
    source-profile prior to be separated.

    \subsection{Principal ranges are output conventions, not physical laws}

    A common principal extraction convention, used throughout this paper, is
    \begin{center}
        \begin{tabular}{@{}lll@{}}
            \toprule
            Sequence family & First and third angles & Middle angle     \\
            \midrule
            Tait--Bryan     & $(-\pi,\pi]$           & $[-\pi/2,\pi/2]$ \\
            Proper Euler    & $(-\pi,\pi]$           & $[0,\pi]$        \\
            \bottomrule
        \end{tabular}
    \end{center}
    SciPy uses the equivalent degree ranges $[-180,180]$ for the outer angles,
    $[-90,90]$ for a Tait--Bryan middle angle, and $[0,180]$ for a proper-Euler middle
    angle \cite{scipyEulerRanges}. General extraction formulas and singularity treatment are
    given by Shuster and Markley and by Bernardes and Viollet
    \cite{shusterMarkley2006,bernardes2022}.

    These ranges describe one representative of a many-to-one chart. They do not imply that
    a physical vehicle cannot pitch beyond $90$ degrees, nor do they imply that a producer
    must wrap angles. A source may emit outer angles in $[0,2\pi)$, continuously unwrap them
    past multiple turns, or deliberately use a nonprincipal Euler branch. Therefore a range
    violation is a hard rejection only when metadata promises principal-value output.

    The exact branch identities explain why. Up to independent additions of $2\pi$ to any
    angle, every Tait--Bryan sequence satisfies
    \begin{equation}
        E_s(\alpha,\beta,\gamma)
        =E_s(\alpha+\pi,\pi-\beta,\gamma+\pi),
        \qquad s\in\TB,
        \label{eq:tait-bryan-branch}
    \end{equation}
    and every proper-Euler sequence satisfies
    \begin{equation}
        E_s(\alpha,\beta,\gamma)
        =E_s(\alpha+\pi,-\beta,\gamma+\pi),
        \qquad s\in\PE.
        \label{eq:proper-euler-branch}
    \end{equation}
    At the singular middle angles the outer angles are not individually identifiable; only a
    sum or difference is fixed, depending on the sequence convention.

    Intrinsic and extrinsic naming also creates exact duplicates when channel reversal is
    allowed:
    \begin{equation}
        E_{abc}^{\mathrm{int}}(\alpha,\beta,\gamma)
        =E_{cba}^{\mathrm{ext}}(\gamma,\beta,\alpha).
        \label{eq:intrinsic-extrinsic-reversal}
    \end{equation}
    An enumerator that includes both composition modes and all angle-channel permutations
    should quotient this identity before numerical scoring.

    \subsection{A principled hierarchy of Euler heuristics}

    Let $\alpha_{h,k},\beta_{h,k},\gamma_{h,k}$ be the raw channels assigned to the
    first, middle, and third angles after a candidate unit conversion. For an interval
    $I\subset\R$, distinguish the ordinary and circular distances
    \begin{align}
        d_{\R}(\theta,I)&=\inf_{u\in I}\abs{\theta-u},\\
        d_{2\pi}(\theta,I)&=\min_{m\in\mathbb Z}d_{\R}(\theta+2\pi m,I).
    \end{align}
    The distinction matters. Circular distance asks whether a value is compatible after a
    whole-turn shift; it cannot verify that the producer literally emitted principal raw
    values.

    Let $I_{\mathrm{out}}=[-\pi,\pi]$ and
    \begin{equation}
        I_{s_h}=\begin{cases}
                    [-\pi/2,\pi/2], & s_h\in\TB,\\
            [0,\pi], & s_h\in\PE.
        \end{cases}
    \end{equation}
    If source profile $p$ promises principal output, use the literal diagnostic
    \begin{equation}
        \begin{split}
            H_{\mathrm{principal}}(h;p)
            =\median_k\bigl[&d_{\R}(\alpha_{h,k},I_{\mathrm{out}})^2
            +d_{\R}(\beta_{h,k},I_{s_h})^2\\
            &+d_{\R}(\gamma_{h,k},I_{\mathrm{out}})^2\bigr].
            \label{eq:euler-principal-score}
        \end{split}
    \end{equation}
    For a profile that permits whole-turn wrapping or unwrapping, the weaker middle-channel
    cue
    \begin{equation}
        H_{\mathrm{middle,mod}}(h)
        =\median_k d_{2\pi}(\beta_{h,k},I_{s_h})^2
        \label{eq:euler-range-score}
    \end{equation}
    may still be useful, but outer-angle bounds should not be penalized. Let
    $H_{\mathrm{range}}(h;p)$ denote the profile-appropriate choice among these diagnostics,
    or zero when the profile places no range prior. The corresponding violation fraction
    $f_{\mathrm{range}}(h;\epsilon)$ should use the same literal or circular distance rule.
    No range quantity is a universal validity test.

    A stronger continuity heuristic explicitly accounts for the two exact branches. For
    each sample, form branch-equivalent triples $a_k^{(0)}$ and $a_k^{(1)}$ using
    \eqref{eq:tait-bryan-branch} or \eqref{eq:proper-euler-branch}. Let $\wrap$ act
    componentwise into $(-\pi,\pi]$, let $W\succeq0$ be a fixed angle-difference weight,
    and let $\lambda_{\mathrm{sw}}\ge0$ penalize branch switches. Define
    \begin{equation}
        H_{\mathrm{branch}}(h)
        =\min_{b_0,\ldots,b_N\in\{0,1\}}
        \sum_{k=1}^{N}
        \norm{\wrap\!\left(a_k^{(b_k)}-a_{k-1}^{(b_{k-1})}\right)}_{W}^2
        +\lambda_{\mathrm{sw}}
        \sum_{k=1}^{N}\mathbf 1\set{b_k\ne b_{k-1}}.
        \label{eq:euler-branch-dp}
    \end{equation}
    This is a two-state dynamic program and costs $O(N)$. It distinguishes a naturally
    continuous branch from a decoder that requires frequent $\pi$ shifts and branch
    switches. It still encodes a smooth-angle prior, not a theorem that the lowest-cost
    sequence is physically correct.

    The singularity margin is sequence-family dependent:
    \begin{equation}
        \sigma_s(\beta)=
        \begin{cases}
            \abs{\cos\beta}, & s\in\TB,\\
            \abs{\sin\beta}, & s\in\PE.
        \end{cases}
        \label{eq:euler-singularity-margin}
    \end{equation}
    Report a low quantile of $\sigma_s$ and the fraction below a configured threshold. A
    candidate that spends substantial time near singularity may still be correct, but its
    angle channels are a poor coordinate chart and angle-domain smoothness becomes
    unreliable.

    Finally, a source-specific physical envelope may be useful. For a conventional ground
    vehicle, aircraft in non-aerobatic operation, or stabilized sensor, one angle may be
    expected to remain much narrower than another. Such an envelope belongs in the source
    profile $p$ and must not be generalized to spacecraft, aerobatic vehicles, tumbling
    objects, or already-unwrapped angle logs.

    \subsection{How to interpret large raw values}

    \begingroup
    \small
    \par\noindent\begin{tabularx}{\linewidth}{@{}p{0.28\linewidth}YY@{}}
\toprule
Observed pattern & Justified inference & Unjustified inference \\
\midrule
One channel is almost always within $[-90,90]$ degrees while two span a wider range
& It is a plausible Tait--Bryan middle channel under a principal-output profile.
& It must be pitch, or the Euler sequence is uniquely known. \\
A proposed Tait--Bryan middle channel exceeds $90$ degrees in magnitude
& It conflicts with principal-value output, subject to tolerance.
& The rotation is invalid; an equivalent nonprincipal branch exists. \\
A proposed proper-Euler middle channel lies outside $[0,180]$ degrees
& It conflicts with principal-value output.
& Proper Euler angles are impossible for the data. \\
Outer angles exceed $180$ degrees or accumulate beyond $360$ degrees
& The source may use $[0,360)$ values or continuous unwrapping.
& The angle cannot be yaw or roll. \\
All three channels stay inside $[-90,90]$ degrees
& The motion envelope is narrow; range alone supplies little sequence information.
& Every channel is equally likely to be the middle angle. \\
Values that look degree-sized under one unit and extreme under radians
& Unit magnitude is a useful source-profile cue.
& Units are identified without timing, rates, or documented bounds. \\
\bottomrule
    \end{tabularx}
    \endgroup

    The recommended policy is therefore lexicographic. Apply documented schema constraints
    first and use trajectory risk whenever an independent rate witness exists. Only among
    surviving classes, or in an explicitly heuristic attitude-only regime, use a declared
    Euler profile key such as
    \begin{equation}
        K_p^{\mathrm{Euler}}(h)
        =\bigl(
        H_{\mathrm{range}},
        H_{\mathrm{branch}},
        f_{\mathrm{sing}},
        H_{\mathrm{physical}}
        \bigr),
        \label{eq:euler-heuristic-key}
    \end{equation}
    where profile $p$ supplies any documented branch policy, tolerances, or physical motion
    envelope. The fraction $f_{\mathrm{sing}}$ counts samples below a declared singularity-
    margin threshold, and $H_{\mathrm{physical}}$ is a source-specific envelope loss. Exact
    schema compatibility is intentionally absent from this key because it was already
    handled as a hard constraint. Do not add these terms to the geometric risk with
    arbitrary weights and then describe the sum as purely data-driven. If several members
    of one exact quotient class receive different representation keys, the key may select a
    representative, but the class itself receives only one evidential vote.

    \begin{sectionrecap}
        Representation heuristics are allowed to abstain and are kept separate from both hard
        schema facts and domain priors. Candidate turn periods $T_u$ support a wrap-gain test,
        but no observed boundary crossing means no wrap-based unit identification. A learned classifier should target explicitly named representation attributes rather than a
        full convention class whose rate/timing semantics are absent from these features. For Euler data, the standard principal middle interval used here is
        $[-\pi/2,\pi/2]$ for Tait--Bryan sequences and $[0,\pi]$ for proper-Euler sequences;
        these are output-branch conventions, not physical attitude limits. Exact branch
        identities and intrinsic/extrinsic reversal are removed by quotienting before heuristic
        or statistical ranking.

        \noindent\textbf{Key equations.}
        The unit conversion is $s_u=2\pi/T_u$; wrap correction uses
        $r_{u,k}=d_k-m_{u,k}T_u$ and
        $H_{\mathrm{wrap}}(u)=-G_u+\lambda_{\mathrm{wrap}}n_u$.
        A calibrated representation classifier may use
        $q_\theta(v\mid\phi_{\rm repr})$ from \eqref{eq:heuristic-softmax} with an explicit
        abstention rule. Tait--Bryan branches satisfy
        \(E_s(\alpha,\beta,\gamma)=E_s(\alpha+\pi,\pi-\beta,\gamma+\pi)\);
        proper-Euler branches use $-\beta$ instead. Literal principal-output checking uses
        $H_{\mathrm{principal}}$, while branch-aware continuity uses the two-state dynamic
        program $H_{\mathrm{branch}}$.

        \noindent\textbf{Symbols introduced.}
        \par\noindent\begin{tabularx}{\linewidth}{@{}p{0.25\linewidth}Y@{}}
\toprule
Symbol & Meaning \\
\midrule
$T_u,s_u$ & one-turn raw period and conversion from raw unit $u$ to radians \\
$G_u,n_u,H_{\mathrm{wrap}}$ & wrap gain, inferred-wrap count, and penalized wrap heuristic \\
$\phi_{\rm repr},q_\theta$ & representation-only feature vector and optional calibrated attribute classifier \\
$E_s(\alpha,\beta,\gamma)$ & Euler decoder for sequence and composition convention $s$ \\
$\TB,\PE$ & Tait--Bryan and proper-Euler sequence families \\
$I_s,I_{\mathrm{out}}$ & selected principal middle and outer intervals \\
$d_{\R},d_{2\pi}$ & literal and whole-turn-invariant interval distances \\
$H_{\mathrm{range}}$ & profile-selected range diagnostic \\
$H_{\mathrm{branch}}$ & minimum two-branch continuity cost \\
$f_{\mathrm{sing}}$ & fraction of samples within the declared singularity-warning margin \\
$K_p^{\mathrm{Euler}}$ & declared lexicographic Euler representative/class-priority key \\
\bottomrule
        \end{tabularx}
    \end{sectionrecap}


    \section{Primary-documentation survey and domain-prior construction}\label{app:survey}

    Trajectory consistency should remain the primary evidence whenever it is available, but
    it is reasonable to ask a second question after an evidence-preserving contender set has
    been formed: \emph{which convention is more typical for the kind of system that probably
    produced this file?}  This section turns that question into a documented, reproducible
    prior rather than an informal preference. The survey is based on official API,
    protocol, and standards documentation, not on a random sample of deployed systems.
    Consequently its scores are \emph{documentation-support scores}, not estimates of
    industry-wide market share or true convention frequencies.

    \subsection{Why the prior must be hierarchical}

    A broad label such as ``science,'' ``NASA,'' ``robotics,'' or ``game engine'' is often
    too coarse. For example, the NASA Planetary Data System explicitly recognizes both a
    SPICE quaternion type with the scalar component first and a spacecraft-telemetry type
    with the scalar component fourth \cite{nasaPDSQuaternion}. Eigen similarly accepts its
    four-scalar constructor in $(w,x,y,z)$ order while its native coefficient vector and
    mapped memory layout are $(x,y,z,w)$ \cite{eigenQuaternion}. Thus even a single
    organization or library can contain more than one legitimate ordering, depending on the
    artifact being decoded.

    Let the available provenance context be
    \begin{equation}
        z=(d,r,p,\iota),
    \end{equation}
    where $d$ is a broad domain, $r$ is an artifact role such as IMU telemetry, spacecraft
    pointing, rigid-body state, or editor/UI angles, $p$ is an ecosystem or protocol, and
    $\iota$ is a particular API, message, file product, or schema. The recommended strength
    ordering is
    \begin{equation}
        \boxed{\begin{aligned}
                   \text{documented schema/API}
                   &\succ \text{ecosystem/protocol}
                   \succ \text{artifact role} \\
                   &\succ \text{broad domain}
                   \succ \text{uninformative fallback}.
        \end{aligned}}
        \label{eq:prior-hierarchy}
    \end{equation}
    An exact documented API field is metadata and may be a hard constraint. A broad-domain
    association is only a weak tie breaker.

    Each surveyed source profile $p$ defines a set
    \begin{equation}
        \mathcal H_p\subseteq\mathcal H
        \label{eq:profile-compatible-set}
    \end{equation}
    of hypotheses compatible with everything that source explicitly documents. An
    unspecified field does not count as a mismatch: if Unity documents a quaternion field
    order but the cited API does not document the rate-expression frame, both rate frames
    remain in $\mathcal H_p$. This set-valued representation is preferable to multiplying
    independent marginal probabilities because real convention choices are coupled. For
    example, scalar order, quaternion multiplication, active/passive action, frame direction,
    and axis handedness should not be recombined into artificial ``Frankenstein'' profiles
    that no surveyed ecosystem uses.

    \subsection{Survey ledger}

    Table~\ref{tab:domain-survey-ledger} records the convention information that is explicit
    in the surveyed primary documentation. A dash means that the cited interface does not
    supply enough information for this survey dimension; it must not be silently inferred.
    The table distinguishes array or wire order from constructor argument order, and it
    separates Euler display units from angular-rate units because several engines use
    opposite choices for those two quantities.

    \begingroup
    \scriptsize
    \setlength{\tabcolsep}{2.6pt}
    \begin{longtable}{@{}p{0.125\textwidth}p{0.205\textwidth}p{0.245\textwidth}p{0.245\textwidth}@{}}
        \caption{Primary-documentation survey of representative attitude and rate interfaces.}\label{tab:domain-survey-ledger}\\
        \toprule
        Profile & Quaternion convention & Euler convention & Angular-rate convention \\
        \midrule
        \endfirsthead
        \toprule
        Profile & Quaternion convention & Euler convention & Angular-rate convention \\
        \midrule
        \endhead
        \bottomrule
        \endlastfoot
        SciPy \texttt{Rotation}
        & Array default $(x,y,z,w)$; scalar-first is optional; active rotations.
        & Sequence is explicit; uppercase intrinsic and lowercase extrinsic; radians by default.
        & -- \\
        Eigen \texttt{Quaternion}
        & Four-scalar constructor $(w,x,y,z)$; native coefficient vector and mapped memory $(x,y,z,w)$.
        & No universal Euler default is used in this survey.
        & -- \\
        MATLAB Navigation/Robotics
        & \texttt{eul2quat} returns $(w,x,y,z)$.
        & Intrinsic $ZYX$ default; radians.
        & -- \\
        MATLAB Aerospace Toolbox
        & Scalar-first $(w,x,y,z)$; documented as right-handed passive transformation.
        & Intrinsic passive $ZYX$ default; radians.
        & -- \\
        JPL/NAIF SPICE
        & SPICE/Hamilton style $(w,x,y,z)$; interfaces require SPICE style.
        & Axis triple is explicit/configurable; moving-axis construction; radians.
        & CK angular velocity is expressed in the base/reference frame, right-handed, rad/s. \\
        NASA PDS quaternion type
        & \texttt{SPICE}: scalar first; \texttt{Spacecraft Telemetry}: scalar fourth.
        & --
        & -- \\
        ROS core / IMU
        & ROS quaternion message serializes $(x,y,z,w)$; IMU orientation is sensor frame w.r.t. world.
        & REP~103 prefers fixed-axis roll--pitch--yaw about $X,Y,Z$; angles use radians.
        & Gyroscope is sensor/body expressed, right-handed, rad/s. \\
        MAVLink / PX4 flight stack
        & MAVLink attitude quaternions and PX4 attitude messages use $(w,x,y,z)$; PX4 documents Hamilton convention.
        & MAVLink \texttt{ATTITUDE} is intrinsic $ZYX$ in radians.
        & Vehicle/target rates are body-frame, rad/s. \\
        Unity 3D
        & \texttt{Quaternion} component/constructor order $(x,y,z,w)$.
        & Transform Euler values are degrees; extrinsic $ZXY$.
        & 3D rigid-body angular velocity is rad/s; cited API does not establish expression frame. \\
        Unreal Engine
        & \texttt{TQuat} constructor/register order $(x,y,z,w)$; composition $A B$ applies $B$ then $A$.
        & \texttt{TRotator} stores degrees; intrinsic yaw--pitch--roll ($Z,Y,X$) with documented sign conventions.
        & Physics exposes world-space rad/s explicitly and also degree/s accessors. \\
        Godot
        & Hamilton quaternion; four-scalar constructor $(x,y,z,w)$.
        & Radians; intrinsic $YXZ$ default.
        & Rigid-body angular velocity is rad/s; cited API does not establish expression frame. \\
    \end{longtable}
    \endgroup

    The underlying documentation is explicit on several important contrasts. SciPy defaults
    to scalar-last quaternion arrays \cite{scipyQuaternion}; Eigen's storage/constructor split
    is documented directly by Eigen \cite{eigenQuaternion}; MATLAB's navigation and
    Aerospace APIs use scalar-first outputs and a default intrinsic $ZYX$ sequence in radians
    \cite{matlabEul2Quat,matlabAeroAngle2Quat}. SPICE requires scalar-first Hamilton-style
    quaternions and its Euler routine accepts an explicit moving-axis sequence in radians
    \cite{naifQuaternion,naifEuler}; its C-kernel angular velocity is base-frame expressed
    in rad/s \cite{naifCK}. ROS uses SI/radian conventions, $(x,y,z,w)$ quaternion message
    fields, and sensor-frame gyro rates \cite{rosrep103,rosrep145,rosQuaternion}. MAVLink
    uses $(w,x,y,z)$ attitudes, intrinsic $ZYX$ Euler attitude in radians, and body rates in
    rad/s \cite{mavlinkCommon}. The three surveyed game engines all expose scalar-last
    quaternion constructors/fields, but their default Euler charts differ: Unity uses
    extrinsic $ZXY$ in degrees, Unreal uses intrinsic $ZYX$ in degrees, and Godot uses
    intrinsic $YXZ$ in radians \cite{unityQuaternion,unityEuler,unrealQuaternion,unrealRotator,godotQuaternion,godotBasis}.

    \subsection{Ranking quaternion conventions}

    For a coarse domain score, count \emph{independent ecosystem profiles}, not documentation
    pages. Multiple pages describing one protocol do not provide multiple votes. Let
    $\mathcal P_{d,r}$ be the profiles relevant to domain $d$ and optional artifact role $r$.
    For one convention dimension $j$, let $I_{p,j}=1$ when profile $p$ is informative about
    that dimension, and let $m_{p,j}(v)\in[0,1]$ be the support that an informative profile gives value $v$; set it to zero when $I_{p,j}=0$.
    A profile with one documented/default value gives it unit support; a genuinely
    set-valued profile may divide support among the values it explicitly admits. The
    transparent survey score is
    \begin{equation}
        S_{d,r,j}(v)
        =100\,
        \frac{\sum_{p\in\mathcal P_{d,r}}w_p I_{p,j}m_{p,j}(v)}
        {\sum_{p\in\mathcal P_{d,r}}w_p I_{p,j}},
        \qquad 0\le S\le100,
        \label{eq:survey-score}
    \end{equation}
    where $w_p=1$ is the default when there is no defensible usage-frequency information.
    The score must always be reported with its informative-profile coverage $N_{d,r,j}$.
    A score of $100$ from one source is much weaker as a broad-domain claim than $100$ from
    three independent ecosystems. Equation~\eqref{eq:survey-score} is a ranking score, not
    a calibrated probability.

    The resulting coarse quaternion-order ranking is summarized below. These numbers are
    best read as ``what the surveyed defaults/wire formats support,'' not ``what percentage
    of all software in the field uses this convention.''

    \begin{center}
        \small
        \begin{tabularx}{\linewidth}{@{}p{0.27\linewidth}p{0.19\linewidth}p{0.19\linewidth}Y@{}}
\toprule
Context & $(x,y,z,w)$ & $(w,x,y,z)$ & Interpretation \\
\midrule
Major game engines & $100$ (3/3) & $0$ (0/3) & Strong scalar-last cue for Unity/Unreal/Godot-like serialized or constructor data. \\
Scientific numeric arrays & $67$ (2/3) & $33$ (1/3) & SciPy and Eigen coefficient arrays are scalar-last; MATLAB numeric quaternion output is scalar-first. Artifact role matters. \\
ROS vs. flight telemetry & $50$ (1/2) & $50$ (1/2) & ROS message fields and MAVLink/PX4 disagree; ``robotics'' alone should not break the tie. \\
NASA/spacecraft, broad label & mixed & mixed & PDS explicitly recognizes both SPICE scalar-first and spacecraft-telemetry scalar-last; require product provenance. \\
Exact NAIF/SPICE API & $0$ & $100$ & Treat as hard metadata when SPICE provenance is known. \\
Exact MAVLink/PX4 attitude & $0$ & $100$ & Treat as hard metadata when the message type is known. \\
Exact ROS quaternion message & $100$ & $0$ & Treat as hard metadata when the message type is known. \\
\bottomrule
        \end{tabularx}
    \end{center}

    This table also shows why ``Hamilton versus JPL'' cannot be reduced to scalar placement.
    SPICE calls its own scalar-first convention Hamilton style and contrasts it with an
    engineering scalar-last style whose vector signs and rotation-matrix association also
    differ \cite{naifQuaternion}. MATLAB Aerospace, PX4, SciPy, Godot, Unity, and Unreal
    have their own action/sign/frame semantics. The hypothesis engine should therefore
    rank a \emph{joint quaternion profile} whenever those semantics are uncertain, not infer
    multiplication or active/passive meaning from the location of $w$ alone.

    \subsection{Ranking Euler conventions by quotient class}

    Euler conventions should be ranked after exact intrinsic/extrinsic reversal duplicates
    have been quotiented. In particular,
    \begin{equation}
        [\mathrm{XYZ}_{\mathrm{ext}}]
    =
    [\mathrm{ZYX}_{\mathrm{int}}]
    \label{eq:euler-rpy-survey-class}
    \end{equation}
    when the angle channels are reversed consistently, and similarly
    \begin{equation}
        [\mathrm{ZXY}_{\mathrm{ext}}]
    =
    [\mathrm{YXZ}_{\mathrm{int}}].
    \label{eq:euler-yxz-survey-class}
    \end{equation}
    This matters in the survey. ROS prefers fixed-axis roll, pitch, and yaw about
    $X,Y,Z$. MAVLink and MATLAB describe intrinsic $ZYX$ yaw--pitch--roll. After the
    angle-channel semantics are aligned, these descriptions land in the same Euler sequence
    quotient class. Unity's extrinsic
    $ZXY$ and Godot's intrinsic $YXZ$ defaults likewise occupy the same structural sequence
    class.

    Three domain-level observations are therefore more useful than a raw list of 24 names:
    \begin{enumerate}[nosep]
    \item \textbf{Angle units.}  Scientific computing, ROS/IMU, MAVLink, and SPICE all
    provide strong documentation support for radians. Among the three surveyed
    major game engines, Unity and Unreal use degrees for Euler-facing APIs while
    Godot's programmatic Euler API uses radians: the broad game-engine score is
    therefore degrees $67$ (2/3), radians $33$ (1/3). Exact engine provenance
    should override that weak domain score.
    \item \textbf{Tait--Bryan versus proper Euler.}  Every surveyed default that fixes a
    three-axis sequence is Tait--Bryan. SciPy and SPICE also support proper-Euler
    families when explicitly requested, so ``Tait--Bryan'' is a useful default
    prior, not a hard rule.
    \item \textbf{Sequence quotient class.}  In the game-engine sample,
    $[\mathrm{YXZ}_{\mathrm{int}}]=[\mathrm{ZXY}_{\mathrm{ext}}]$ receives support
    from Unity and Godot, or $67$ (2/3), while
    $[\mathrm{ZYX}_{\mathrm{int}}]=[\mathrm{XYZ}_{\mathrm{ext}}]$ receives Unreal's
    $33$ (1/3). In the surveyed flight/robotics RPY interfaces, MATLAB,
    MAVLink, and ROS all support the latter $ZYX/XYZ$ quotient class, despite
    differing terminology.
    \end{enumerate}

    A sequence rank should therefore be stored as a quotient-class preference first and a
    representative naming preference second. Otherwise the survey can accidentally give
    two votes to mathematically identical intrinsic/extrinsic descriptions.

    \subsection{Ranking angular-rate conventions}

    Rate priors should be conditioned on \emph{what generated the rate}, not only on the
    software domain. The survey gives a particularly clean contrast:
    \begin{center}
        \small
        \begin{tabularx}{\linewidth}{@{}p{0.30\linewidth}p{0.19\linewidth}p{0.19\linewidth}Y@{}}
\toprule
Artifact role & Rate unit & Expression frame & Survey implication \\
\midrule
ROS IMU / gyro & rad/s & sensor/body & Strong body-rate profile when the stream is sensor telemetry \cite{rosrep145}. \\
MAVLink/PX4 vehicle rate & rad/s & body/FRD & Strong body-rate profile for autopilot telemetry and attitude targets \cite{mavlinkCommon,px4VehicleOdometry}. \\
NAIF SPICE CK pointing & rad/s & base/reference & Strong world/base-rate profile for CK angular-velocity data \cite{naifCK}. \\
Unity 3D Rigidbody & rad/s & unspecified by cited API & Unit cue only; do not infer body/world frame \cite{unityAngularVelocity}. \\
Unreal physics body & rad/s API and degree/s API & world for explicit world-rad/s getter & API name is highly informative; broad engine unit prior must retain both possibilities \cite{unrealAngularVelocity,unrealAngularDegrees}. \\
Godot RigidBody3D & rad/s & unspecified by cited API & Unit cue only; do not infer body/world frame \cite{godotAngularVelocity}. \\
\bottomrule
        \end{tabularx}
    \end{center}

    Thus ``IMU/flight-controller telemetry'' has a defensible body-frame/rad/s prior, while
    ``SPICE CK pointing'' has the opposite base/world-frame prior. A generic
    ``physics-engine'' label supports radians for 3-D angular velocity but is not sufficient
    to assign the expression frame from this survey. The different engine APIs also show
    why Euler units and rate units must be classified separately: Unity displays Euler
    angles in degrees while its 3-D rigid-body angular velocity is rad/s; Unreal's Rotator
    is degree based while its physics API deliberately exposes both rad/s and degree/s
    forms \cite{unityEuler,unityAngularVelocity,unrealRotator,unrealAngularVelocity,unrealAngularDegrees}.

    \subsection{Realm scorecards and what they actually rank}\label{sec:realm-scorecards}

    The survey is easiest to use when its conclusions are summarized at the same level at
    which a user normally knows provenance. Table~\ref{tab:realm-scorecard} therefore gives
    broad-domain scorecards, while preserving the informative-profile count behind each
    number. ``Mixed'' is an intentional result: it means the broad label should not be used
    to break that dimension's tie.

    \begingroup
    \small
    \setlength{\tabcolsep}{3.4pt}
    \begin{longtable}{@{}p{0.20\textwidth}p{0.24\textwidth}p{0.27\textwidth}p{0.20\textwidth}@{}}
        \caption{Broad realm scorecard derived from the surveyed official interfaces.}
        \label{tab:realm-scorecard}\\
        \toprule
        Realm/context & Quaternion order & Euler prior & Rate prior \\
        \midrule
        \endfirsthead
        \toprule
        Realm/context & Quaternion order & Euler prior & Rate prior \\
        \midrule
        \endhead
        \bottomrule
        \endlastfoot
        Scientific/numerical APIs
        & Scalar-last $67$ (2/3), scalar-first $33$ (1/3) for the surveyed numeric array/output roles.
        & Radians $100$ (2/2 informative defaults).  $ZYX$ is the only fixed default among the surveyed interfaces that supply one, so sequence coverage is weak ($1/1$).
        & No broad rate-frame prior from this sample. \\
        Robotics/autonomy: ROS versus MAVLink/PX4
        & Scalar-last $50$ (1/2), scalar-first $50$ (1/2): broad ``robotics'' order is non-informative.
        & Radians $100$ (2/2); the fixed-$XYZ$/intrinsic-$ZYX$ roll--pitch--yaw descriptions occupy the same structural quotient class, $100$ (2/2).
        & Body/sensor expressed and rad/s $100$ (2/2) for the surveyed IMU/autopilot roles. \\
        NASA/spacecraft, broad
        & Mixed by specification: PDS explicitly admits scalar-first SPICE and scalar-fourth telemetry. No broad scalar-order tie breaker is recommended.
        & SPICE uses radians but leaves the axis sequence explicit; no broad sequence default is inferred.
        & No broad NASA rate-frame prior. For the specific CK pointing role, use base/reference-frame rad/s as metadata/profile information. \\
        Major game engines
        & Scalar-last $100$ (3/3) for Unity, Unreal, and Godot quaternion constructor/component interfaces.
        & Degrees $67$ (2/3), radians $33$ (1/3). Sequence quotient: $[YXZ_{int}/ZXY_{ext}]$ $67$ (2/3), $[ZYX_{int}/XYZ_{ext}]$ $33$ (1/3).
        & Rad/s is documented by all three surveyed rigid-body APIs, but Unreal also exposes degree/s forms; broad expression frame remains unresolved. \\
    \end{longtable}
    \endgroup

    The scorecard has two distinct uses. If two \emph{different quotient classes} remain in
    the evidence contender set, the domain survey may weakly rank those classes. If two raw
    hypotheses are members of the \emph{same exact quotient class}, trajectory evidence cannot
    rank them at all; the survey may only choose the ecosystem-typical representative. In
    symbols, the intended flow is
    \begin{equation}
        \boxed{
            \mathcal T_{\mathrm{ev}}
            \xrightarrow{\;\pi_{\mathrm{survey}}\;}
            C_{\mathrm{op}}
            \xrightarrow{\;\pi_{\mathrm{rep}}\;}
            h_{\mathrm{op}}.}
        \label{eq:survey-two-level-use}
    \end{equation}
    The first arrow is a prior over unresolved physical convention classes; the second arrow
    is a naming/storage policy inside an observational equivalence class. Neither arrow
    changes the recorded trajectory residuals.

    A broad-domain score should also decay in influence when its coverage is sparse or its
    top values are close. A simple reporting diagnostic is the normalized margin
    \begin{equation}
        M_{d,r,j}=\frac{S_{(1)}-S_{(2)}}{100},
        \qquad 0\le M_{d,r,j}\le1,
        \label{eq:survey-margin}
    \end{equation}
    where $S_{(1)}$ and $S_{(2)}$ are the largest two support scores for dimension $j$.
    The margin is not a probability. It is a warning flag: $M=0$ means the survey itself
    has no preference, whereas a large margin with large coverage is a stronger reason to
    use the prior as a tie breaker. Exact API metadata still outranks any such score.

    \subsection{From survey scores to a domain-conditioned prior}

    The marginal scores above are useful diagnostics, but the solver should use the joint
    profile sets $\mathcal H_p$. Given a domain/artifact context $(d,r)$ and no measured
    usage frequencies, use an equal mixture over independent source profiles,
    \begin{equation}
        \mu_p(d,r)=\frac{1}{\abs{\mathcal P_{d,r}}},
    \end{equation}
    unless a project has defensible source-frequency data. Profiles incompatible with all
    classes surviving hard metadata are flagged rather than assigned a zero-denominator
    weight. For quotient class $C\in\mathcal Q$, define the number of quotient classes
    compatible with profile $p$ as
    \begin{equation}
        Z_p=\abs{\set{C'\in\mathcal Q:C'\cap\mathcal H_p\ne\varnothing}}.
    \end{equation}
    Then a documentation-derived class prior is
    \begin{equation}
        \pi_{\mathrm{survey}}(C\mid d,r)
        =
        \sum_{p\in\mathcal P_{d,r}}
        \mu_p(d,r)
        \frac{\mathbf 1\{C\cap\mathcal H_p\ne\varnothing\}}{Z_p}.
        \label{eq:survey-class-prior}
    \end{equation}
    The division by $Z_p$ prevents a partially specified profile from creating extra mass
    merely because it leaves many hypothesis dimensions unspecified. It also avoids
    rewarding a quotient class for containing more redundant enum members.

    When broad-domain information is weak, temper the survey prior toward uniform over the
    evidence contender set $\mathcal T_{\mathrm{ev}}$. First define its surviving survey
    mass
    \begin{equation}
        Z_{\rm ev}(d,r)
        =\sum_{C'\in\mathcal T_{\mathrm{ev}}}\pi_{\mathrm{survey}}(C'\mid d,r).
    \end{equation}
    Then use
    \begin{equation}
        \begin{aligned}
            &\pi_{\beta_{\rm dom}}(C\mid d,r,\mathcal T_{\mathrm{ev}})\\
            &\quad=\begin{cases}
                       \dfrac{1-\beta_{\rm dom}}{\abs{\mathcal T_{\mathrm{ev}}}}
                       +\beta_{\rm dom}\dfrac{\pi_{\mathrm{survey}}(C\mid d,r)}{Z_{\rm ev}(d,r)},
                       & Z_{\rm ev}(d,r)>0,\\[8pt]
                       \dfrac{1}{\abs{\mathcal T_{\mathrm{ev}}}},
                       & Z_{\rm ev}(d,r)=0.
            \end{cases}
        \end{aligned}
        \label{eq:tempered-domain-prior}
    \end{equation}
    for $C\in\mathcal T_{\mathrm{ev}}$ and $0\le\beta_{\rm dom}\le1$.
    Here $\beta_{\rm dom}$ is \emph{prior strength}, not a parameter learned from the trajectory
    score. Suggested interpretation is qualitative: exact API provenance should normally
    be promoted to hard metadata rather than handled with a large $\beta_{\rm dom}$; a known
    protocol or ecosystem may justify a moderate $\beta_{\rm dom}$; a broad domain alone should use
    small $\beta_{\rm dom}$ or simply serve as a lexicographic tie breaker. The second branch
    makes abstention explicit: if the survey has no support inside the contender set, it
    contributes no ranking information and the solver falls back to canonical policy.

    The final class ranking remains evidence first:
    \begin{equation}
        \boxed{\begin{aligned}
                   \text{hard metadata}
                   &\longrightarrow \text{trajectory/statistical contender set} \\
                   &\longrightarrow \text{survey/domain class prior} \\
                   &\longrightarrow \text{within-class representative policy}.
        \end{aligned}}
        \label{eq:survey-ranking-hierarchy}
    \end{equation}
    Within the selected quotient class, the same survey can choose a representative
    semantic description without changing class-level evidence. Let
    \begin{equation}
        A_h(C,d,r)
        =\sum_{p\in\mathcal P_{d,r}}
        \mu_p(d,r)\,\mathbf 1\{h\in\mathcal H_p\},
        \qquad h\in C,
    \end{equation}
    and, when $\sum_{h'\in C}A_{h'}>0$, normalize
    \begin{equation}
        \pi_{\mathrm{rep}}(h\mid C,d,r)
        =\frac{A_h(C,d,r)}{\sum_{h'\in C}A_{h'}(C,d,r)}.
        \label{eq:survey-representative-prior}
    \end{equation}
    If the denominator is zero, the survey abstains and the deterministic canonical policy
    chooses the representative. This two-level construction is important: class priors rank empirically unresolved
    \emph{physical convention classes}; representative priors decide which naming/storage
    member of an exact class is most natural for the suspected ecosystem. The latter must
    never acquire extra evidential weight just because the software enumerated more
    representatives.

    \subsection{Practical scoring output}

    For each run, the hypothesis engine should be able to emit a compact prior report such as
    \begin{quote}
        \small
        Context: \texttt{game-engine / 3d-rigidbody / engine-unknown}.\\
        Quaternion order survey: $(x,y,z,w)$ $100$ with coverage $3/3$.\\
        Euler unit survey: degrees $67$, radians $33$, coverage $3/3$.\\
        Euler sequence class survey: $[YXZ_{int}/ZXY_{ext}]$ $67$,
        $[ZYX_{int}/XYZ_{ext}]$ $33$, coverage $3/3$.\\
        Rate unit survey: radians receives the strongest 3-D rigid-body support, but Unreal also
        has a degree/s API; rate frame remains \texttt{unresolved-by-survey}.\\
        Selection use: tie breaker only; trajectory score unchanged.
    \end{quote}
    For an exact profile such as \texttt{MAVLink ATTITUDE\_QUATERNION} or a SPICE CK product,
    most of these entries should instead be promoted from ``prior'' to ``documented
    metadata.'' This prevents a probabilistic-looking survey score from weakening a known
    schema fact.

    \begin{sectionrecap}
        The survey supplies a hierarchical, documentation-derived prior after trajectory
        evidence, not a replacement for trajectory evidence. Broad domains are frequently
        mixed; exact ecosystem and artifact role are much more informative. Game engines in
        the surveyed set unanimously use scalar-last quaternion constructors/fields, but their
        Euler defaults split across two quotient sequence classes and two angle-unit choices.
        ROS/MAVLink/PX4-like IMU and flight telemetry strongly support body rates in rad/s,
        whereas SPICE CK supports base-frame rates in rad/s. NASA/PDS itself documents both
        scalar-first SPICE and scalar-last spacecraft-telemetry quaternion types.
        Table~\ref{tab:realm-scorecard} summarizes the broad-domain rankings and their coverage.

        \noindent\textbf{Key equations.}
        The raw marginal support score is
        \(S_{d,r,j}(v)=100\sum_p w_p I_{p,j}m_{p,j}(v)/\sum_p w_p I_{p,j}\).
        Joint source profiles define $\mathcal H_p$; class-level survey mass is assigned by
        \eqref{eq:survey-class-prior}, optionally tempered by \eqref{eq:tempered-domain-prior},
        and representative selection uses \eqref{eq:survey-representative-prior}. Euler survey
        counts operate on structural sequence classes such as
        $[XYZ_{ext}]=[ZYX_{int}]$, not duplicated raw labels.

        \noindent\textbf{Symbols introduced.}
        \par\noindent\begin{tabularx}{\linewidth}{@{}p{0.25\linewidth}Y@{}}
\toprule
Symbol & Meaning \\
\midrule
$d,r,p,\iota$ & broad domain, artifact role, ecosystem/profile, and exact API/schema context \\
$\mathcal H_p$ & hypothesis set compatible with documented source profile $p$ \\
$S_{d,r,j}(v)$ & 0--100 documentation-support score for one convention value \\
$N_{d,r,j}$ & number of independent informative profiles behind a marginal score \\
$M_{d,r,j}$ & normalized gap between the two leading marginal survey scores \\
$\mu_p$ & mixture weight assigned to one independent source profile \\
$Z_p$ & number of quotient classes compatible with a partially specified profile \\
$\pi_{\mathrm{survey}}$ & joint documentation-derived quotient-class prior \\
$\beta_{\rm dom}$ & strength used to temper a broad survey prior toward uniform \\
$Z_{\rm ev}$ & survey-prior mass carried by the current evidence contender set \\
$A_h$ & survey support assigned to representative $h$ within one quotient class \\
$\pi_{\mathrm{rep}}$ & normalized within-class source-profile preference over representatives \\
\bottomrule
        \end{tabularx}
    \end{sectionrecap}


    \section{Detailed selection mechanisms and canonicalization}\label{app:selection}

    \subsection{Connection to the main result object}

    Section~\ref{sec:result} defines the returned object
    \begin{equation}
        \mathcal R=
        (\mathcal T_{\mathrm{ev}},\,C_{\mathrm{op}},\,h_{\mathrm{op}},\,
        \ell_{\mathrm{reason}}).
        \label{eq:final-result-object}
    \end{equation}
    This appendix does not redefine that output. It only specifies how numerical ties,
    statistical uncertainty, source-profile priors, practical overrides, and deterministic
    canonicalization act on the four fields. In particular, $\mathcal T_{\mathrm{ev}}$ is
    inferential; $C_{\mathrm{op}}$ and $h_{\mathrm{op}}$ are operational selections whose
    reason must remain explicit.

    The quotient formulation answers what the trajectory can identify, but an operational
    system often still needs one decoder to instantiate. That second decision should be
    reproducible without being confused with statistical evidence. The recommended policy
    has two outputs:
    \begin{align}
        \text{inferential output}
        &=\text{one or more quotient classes},\\
        \text{operational output}
        &=\text{one labeled representative}.
    \end{align}
    The representative is a gauge choice or a prior-driven decision, not a hidden extension
    of the geometric score.

    \subsection{What established selection mechanisms contribute}

    Several familiar mechanisms address different parts of this problem. They should not be
    collapsed into one undocumented scalar penalty.

    \begingroup
    \small
    \setlength{\tabcolsep}{4pt}
    \setlength{\LTleft}{0pt}
    \setlength{\LTright}{0pt}
    \begin{longtable}{@{}>{\raggedright\arraybackslash}p{0.18\textwidth}>{\raggedright\arraybackslash}p{0.39\textwidth}>{\raggedright\arraybackslash}p{0.35\textwidth}@{}}
        \toprule
        Mechanism & Legitimate role & Effect on an exact observational tie \\
        \midrule
        \endfirsthead
        \toprule
        Mechanism & Legitimate role & Effect on an exact observational tie \\
        \midrule
        \endhead
        \midrule
        \multicolumn{3}{r}{\emph{Continued on next page}}\\
        \endfoot
        \bottomrule
        \endlastfoot
        Quotient or set-valued $\argmin$
        & Preserve every data-indistinguishable solution and report the identifiable object;
        a simultaneous ambiguity or model-confidence set quantifies unresolved near ties
        \cite{hansen2011}.
        & Does not break an exact tie; this is the scientific default. \\
        Held-out risk or cross-validation
        & Compare distinct quotient classes after all nuisance parameters have been fitted on
        training blocks; guard against winning through in-sample adaptation \cite{stone1974}.
        & Cannot break a symmetry whose predictions agree on every possible holdout. \\
        Hard metadata or gauge anchor
        & Enforce a documented schema, known frame, known unit, timestamp, or other external
        fact. Geometric inverse problems likewise fix a datum to select one gauge representative
        \cite{triggs2000}.
        & Removes candidates using external information. \\
        Bayesian or loss-based prior
        & Combine class-level prior belief with a likelihood or a calibrated loss update
        \cite{kass1995,bissiri2016}.
        & Equal data evidence leaves posterior odds equal to prior odds. \\
        AIC or BIC
        & Penalize genuine likelihood models for fitted dimension
        \cite{akaike1974,schwarz1978}.
        & Equal maximized likelihood and equal parameter count remain tied. \\
        Minimum description length
        & Compare an explicitly coded model-plus-data description, including continuous
        parameters and residuals \cite{rissanen1978}.
        & A shorter code for ``identity'' than for ``axis flip'' is a declared coding prior unless
        that code is part of the physical data-generating model. \\
        Canonicalization
        & Choose one stable representative by an explicit minimum-intervention or
        interoperability rule.
        & Breaks the software tie only; it adds no physical evidence. \\
    \end{longtable}
    \endgroup

    For distinct quotient classes, held-out predictive risk is the safest default when the
    residual model is intentionally robust rather than probabilistic. Fit period, bias,
    latency, scale, and any other continuous nuisance parameters on training blocks, freeze
    them, and evaluate the class on disjoint validation blocks. Information criteria are
    appropriate only after specifying an actual likelihood and a defensible effective sample
    size. Minimum description length is appropriate only after specifying the code; otherwise
    ``shorter decoding transformation'' is simply another name for a canonical preference.

    \begin{proposition}[Common selectors under exact ties]
Let $h$ and $h'$ have equal maximized likelihood
$\widehat{\mathcal L}_h=\widehat{\mathcal L}_{h'}$ and the same number $k$ of fitted
parameters. Then
\begin{equation}
    \operatorname{AIC}(h)=\operatorname{AIC}(h'),
    \qquad
    \operatorname{BIC}(h)=\operatorname{BIC}(h').
\end{equation}
If an MDL criterion is written as
\begin{equation}
    \operatorname{DL}(h;\mathcal D)
    =\ell_{\mathrm{model}}(h)+\ell_{\mathrm{data}}(\mathcal D\mid h),
\end{equation}
then equal data codelengths remain tied unless the declared model code assigns
$\ell_{\mathrm{model}}(h)\ne\ell_{\mathrm{model}}(h')$. For a loss-based update with
external information $z$,
\begin{equation}
    \Pi_\eta(h\mid\mathcal D,z)
    \propto
    \pi(h\mid z)\exp\!\left(-\eta J_{q(h)}\right),
    \qquad \eta>0,
    \label{eq:gibbs-posterior}
\end{equation}
if $q(h)=q(h')$, then $J_{q(h)}=J_{q(h')}$ and
\begin{equation}
    \frac{\Pi_\eta(h\mid\mathcal D,z)}
    {\Pi_\eta(h'\mid\mathcal D,z)}
    =
    \frac{\pi(h\mid z)}{\pi(h'\mid z)}.
\end{equation}
    \end{proposition}

    \begin{proof}
        AIC and BIC have the forms
        $-2\log\widehat{\mathcal L}+2k$ and
        $-2\log\widehat{\mathcal L}+k\log n$, respectively, so their equalities are immediate.
        The MDL statement follows by cancellation of equal data codelengths. The exponential loss terms in \eqref{eq:gibbs-posterior} cancel when the two hypotheses
        map to the same class risk. The scale
        $\eta$ must be calibrated before \eqref{eq:gibbs-posterior} is interpreted
        probabilistically; the robust score \eqref{eq:robust-risk} does not supply that
        calibration automatically.
    \end{proof}

    The MDL interpretation makes the minimum-intervention rule precise: a code may assign
    short descriptions to identity channel order, native storage, no added inverse, and no
    extra nuisance parameters. Such a code can be operationally useful, but another
    ecosystem can rationally use a different code. MDL therefore justifies a declared
    parsimony policy; it does not turn that code into trajectory evidence
    \cite{rissanen1978}.

    For exact gauge classes, placing a prior directly on every enumerated hypothesis can
    make the result depend on how many redundant encodings the implementation happened to
    list. A cleaner hierarchy places a prior on quotient classes and a separate canonical
    policy within each class:
    \begin{equation}
        \Pi_\eta(C\mid\mathcal D,z)
        \propto
        \pi_{\mathrm{class}}(C\mid z)\exp(-\eta J_C),
        \qquad C\in\mathcal Q,
        \label{eq:class-posterior}
    \end{equation}
    where $J_C$ is the common geometric risk of members of $C$. The within-class
    representative is then selected by a deterministic key, not by duplicating prior mass.

    \subsection{Source profiles instead of a universal convention prior}

    Appendix~\ref{app:survey} develops the primary-documentation survey and the two-level class/representative prior used here. There is no universal quaternion storage or axis profile. ROS REP~103 standardizes SI
    units, right-handed frames, body axes $x$ forward, $y$ left, $z$ up, and ENU geographic
    frames; the standard quaternion message serializes fields as $(x,y,z,w)$
    \cite{rosrep103,rosQuaternion}. ROS REP~145 further specifies a gyroscope rate as the
    rotational velocity of the sensor frame with respect to an inertial frame, expressed in
    the sensor frame, right-handed and in radians per second; a fused quaternion represents
    the sensor frame with respect to the world frame \cite{rosrep145}. MAVLink's common
    attitude message instead specifies the right-handed aeronautical frame $x$ front, $y$
    right, $z$ down and quaternion order $(w,x,y,z)$; its attitude-target message explicitly
    uses a local-NED-to-body-FRD quaternion and body rates in radians per second
    \cite{mavlinkCommon}. SciPy's rotation API defaults to scalar-last active quaternions,
    while Eigen accepts constructor arguments in $(w,x,y,z)$ order but stores coefficients
    internally as $(x,y,z,w)$ \cite{scipyQuaternion,eigenQuaternion}. These are all
    legitimate conventions, and their differences are exactly why provenance must precede a
    global popularity ranking.

    Therefore external information $z$ should first select a \emph{source profile} $p$.
    Examples include \texttt{ros}, \texttt{mavlink}, \texttt{scipy}, a documented vendor
    schema, or a project-specific interface control document. A profile may specify
    \begin{equation}
        p=(\text{units},\text{chirality},\text{axis labels},\text{quaternion order},
        \text{attitude direction},\text{rate frame},\text{timing semantics}).
    \end{equation}
    A schema field that is explicitly documented is a hard constraint. A producer name,
    file extension, variable name, or software-library fingerprint is weaker evidence and
    should enter only as a soft class prior. Provenance-specific rules dominate a generic
    fallback because convention popularity is ecosystem-dependent.

    When provenance is absent, the recommended generic profile is
    \emph{minimum intervention}:
    \begin{enumerate}[nosep]
    \item use right-handed coordinate bases and proper basis changes unless a reflected
    basis is explicitly documented;
    \item prefer SI units, including angular velocity in radians per second
    \cite{bipmSI};
    \item prefer the identity rate-channel map before permutations or sign reversals;
    \item prefer the recorded attitude direction without an added transpose or inverse;
    \item prefer the ingestion schema's native component order without an added
    quaternion permutation;
    \item use body- or sensor-resolved angular velocity as a weak final semantic default,
    reflecting documented IMU and vehicle-rate interfaces in ROS REP~145 and
    MAVLink \cite{rosrep145,mavlinkCommon}, not a universal law;
    \item prefer no extra latency, time warp, or continuous calibration parameter unless
    validation demonstrates that the added flexibility is needed; and
    \item finish with a stable lexicographic code independent of search enumeration.
    \end{enumerate}
    If the input fields are anonymous, such as \texttt{q0}, \texttt{q1}, \texttt{q2},
    \texttt{q3}, the generic profile must not pretend that scalar-first or scalar-last is a
    physical fact. The parser's declared native order may be used as a reproducibility
    choice, but the report must label it as canonical rather than identified.

    \subsection{A lexicographic canonical key}

    A weighted sum can trade a serious profile contradiction against several cosmetic
    preferences. A lexicographic key makes the hierarchy explicit. For profile $p$ and
    metadata $z$, define
    \begin{equation}
        \begin{split}
            K_p(h;z)=(&c_{\mathrm{schema}},c_{\mathrm{profile}},c_{\mathrm{basis}},
            c_{\mathrm{unit}},c_{\mathrm{channel}},c_{\mathrm{direction}},
            c_{\mathrm{storage}},c_{\mathrm{frame}},c_{\mathrm{nuisance}},
            c_{\mathrm{code}}).
            \label{eq:canonical-key}
        \end{split}
    \end{equation}
    The fields are ordered from strongest to weakest:
    \begin{itemize}[nosep]
    \item $c_{\mathrm{schema}}$ is zero only when all hard metadata constraints hold;
    nonzero values make a candidate infeasible rather than merely less preferred;
    \item $c_{\mathrm{profile}}$ counts or weights soft mismatches with the selected
    source profile;
    \item $c_{\mathrm{basis}}$ penalizes an explicitly reflected coordinate basis, but
    does not confuse a sensor-channel sign map with a basis reflection;
    \item $c_{\mathrm{unit}}$ penalizes departure from the profile unit, with SI as the
    generic fallback;
    \item $c_{\mathrm{channel}}$ is itself the tuple
    $(\mathbf 1\{C_h\ne I\},n_{\mathrm{sign}},n_{\mathrm{moved}},n_{\mathrm{inv}})$,
    where the last two entries count moved channels and permutation inversions;
    \item $c_{\mathrm{direction}}$ penalizes an added attitude transpose or inverse;
    \item $c_{\mathrm{storage}}$ penalizes component reordering relative to the profile;
    \item $c_{\mathrm{frame}}=0$ for the profile rate frame and, in the generic profile,
    uses body before world only as a weak fallback;
    \item $c_{\mathrm{nuisance}}$ orders models by added latency, clock, bias, scale, or
    calibration freedom after data fit has tied; and
    \item $c_{\mathrm{code}}$ is a stable serialized hypothesis identifier.
    \end{itemize}
    Within a structural class $C$, select
    \begin{equation}
        h_{\mathrm{can}}(C)
        =\lexargmin_{h\in C}K_p(h;z).
        \label{eq:canonical-representative}
    \end{equation}
    Equation \eqref{eq:canonical-representative} is invariant to hypothesis enumeration and
    is auditable field by field. It should not be folded into the geometric class risk $J_C$.

    \subsection{Numerical ties, statistical ties, and the final hierarchy}

    Structural equivalence is established analytically and should not be inferred from a
    floating-point tolerance. For implementation tests only, let $J_h^{\rm test}$ denote
    the score obtained by directly evaluating raw member $h$ instead of reusing the stored
    class score. If every member of a known class is evaluated, define
    \begin{equation}
        \overline J_C^{\rm test}=\frac{1}{\abs C}\sum_{h\in C}J_h^{\rm test},
        \qquad
        \delta_C^{\mathrm{sym}}
        =\max_{h\in C}\abs{J_h^{\rm test}-\overline J_C^{\rm test}}.
    \end{equation}
    The spread $\delta_C^{\mathrm{sym}}$ is a software diagnostic. It must remain below a
    configured symmetry tolerance; otherwise the implementation has violated a proven
    invariance. Production scoring should preferably evaluate one representative per class
    and evaluate the other members only in tests.

    For distinct classes, define a numerical resolution
    \begin{equation}
        \varepsilon_{\mathrm{num}}
        \ge
        \max\set{
            \varepsilon_{\mathrm{fp}},
            \varepsilon_{M\rightarrow2M},
            \max_C\delta_C^{\mathrm{sym}},
            \varepsilon_{\mathrm{repeat}}
        },
    \end{equation}
    where the terms bound floating-point variation, integration refinement, symmetry-test
    spread, and repeat-run/backend variation. The numerical tie set is
    \begin{equation}
        \mathcal T_{\mathrm{num}}
        =\set{C:J_C-J_{\min}\le\varepsilon_{\mathrm{num}}}.
    \end{equation}
    This is not yet a statistical statement. With independent validation blocks or a
    justified block bootstrap, retain a class whenever a simultaneous confidence interval
    for its excess population risk contains zero. Under the bounded-loss result above, use
    \begin{equation}
        \mathcal T_{\mathrm{stat}}(n,\delta)
        =\mathcal A_n(\delta),
    \end{equation}
    with $\mathcal A_n(\delta)$ defined in \eqref{eq:ambiguity-set}. A full model-confidence-
    set procedure may instead use simultaneous bootstrap tests and sequential elimination
    \cite{hansen2011}. If no uncertainty analysis is available, the report may say
    ``numerically tied,'' but not ``statistically indistinguishable.''

    A third, deliberately nonstatistical set may encode downstream practical indifference.
    For a declared application-scale threshold $\varepsilon_{\mathrm{pr}}\ge0$, define
    \begin{equation}
        \mathcal T_{\mathrm{pr}}
        =\set{C:J_C-J_{\min}\le\varepsilon_{\mathrm{pr}}}.
        \label{eq:practical-tie-set}
    \end{equation}
    This threshold is useful when a tiny score advantage does not justify a less interoperable
    or more complicated decoder. It must be chosen from downstream consequences, not tuned
    until a preferred convention wins. A class may be statistically worse yet practically
    interchangeable; selecting it is a policy decision and must be labeled as such.

    The default operational policy needs one unambiguous contender set. Define the
    \emph{evidence-preserving set}
    \begin{equation}
        \mathcal T_{\mathrm{ev}}
        =
        \begin{cases}
            \mathcal T_{\mathrm{stat}}(n,\delta),
            & \text{when a justified simultaneous uncertainty analysis is available},\\
            \mathcal T_{\mathrm{num}},
            & \text{otherwise}.
        \end{cases}
        \label{eq:evidence-preserving-set}
    \end{equation}
    This set is part of the inferential output. It is never replaced silently by a
    source-preference rule. To extend the within-class key to quotient classes, define
    \begin{equation}
        \overline K_p(C;z)=K_p\!\left(h_{\mathrm{can}}(C);z\right).
    \end{equation}
    If $\mathcal T_{\mathrm{ev}}$ is not a singleton, a declared source-profile prior may
    first reduce it to the highest-priority classes. The induced lexicographic class key then
    breaks any remaining prior tie and supplies the complete fallback when no prior exists:
    \begin{align}
        \mathcal P_{p,z}
        &=
        \begin{cases}
            \displaystyle\operatorname*{arg\,max}_{C\in\mathcal T_{\mathrm{ev}}}
            \pi_{\mathrm{class}}(C\mid z),
            & \text{if a declared class prior is available},\\[6pt]
            \mathcal T_{\mathrm{ev}},
            & \text{otherwise},
        \end{cases}\\
        C_{\mathrm{op}}
        &=\lexargmin_{C\in\mathcal P_{p,z}}\overline K_p(C;z),
        \qquad
        h_{\mathrm{op}}=h_{\mathrm{can}}(C_{\mathrm{op}}).
        \label{eq:default-operational-selection}
    \end{align}
    For exact equal-score classes, the first stage reduces to prior odds and the second to an
    explicit canonical fallback. For statistically unresolved near ties, both stages are an
    operational decision rather than proof that the selected class is physically correct.

    A practical policy may instead choose some $C_{\mathrm{pr}}\in\mathcal T_{\mathrm{pr}}$
    for interoperability or downstream utility. If
    $C_{\mathrm{pr}}\notin\mathcal T_{\mathrm{ev}}$, the report must preserve
    $\mathcal T_{\mathrm{ev}}$, state the sacrificed score
    \begin{equation}
        \Delta J_{\mathrm{override}}=J_{C_{\mathrm{pr}}}-J_{\min},
    \end{equation}
    and label the result as a \texttt{practical-policy-override}. This keeps an engineering
    preference from being retroactively presented as statistical evidence.

    The recommended decision hierarchy is therefore:
    \begin{enumerate}[nosep]
    \item reject hard-schema violations;
    \item quotient exact structural symmetries;
    \item rank quotient classes by held-out, hypothesis-independent trajectory risk;
    \item when candidate forward models have different fitted flexibility, compare them
    by held-out risk or a justified common likelihood/description criterion rather
    than raw in-sample residual alone;
    \item retain and report numerical, confidence-calibrated, and practical near-tie sets
    separately;
    \item form $\mathcal T_{\mathrm{ev}}$ by \eqref{eq:evidence-preserving-set};
    \item choose the default operational class only within $\mathcal T_{\mathrm{ev}}$,
    using a declared source-profile class prior or the induced lexicographic class
    key in \eqref{eq:default-operational-selection};
    \item treat any selection from $\mathcal T_{\mathrm{pr}}\setminus
    \mathcal T_{\mathrm{ev}}$ as an explicit practical-policy override;
    \item choose the within-class representative by the section
    \eqref{eq:canonical-section}, implemented through
    \eqref{eq:canonical-representative}; and
    \item use the stable named-field code only as the final deterministic fallback.
    \end{enumerate}
    The output should carry a selection label. Recommended values are shown explicitly so
    that long machine-readable names do not become prose:
    \begin{quote}
        \small\ttfamily
        trajectory-unique; exact-gauge-class; hard-metadata; domain-prior;\\
        heuristic-only; practical-policy-override; canonical-policy.
    \end{quote}
    A single reported decoder without that label hides the most important distinction in the
    problem.


    \begin{sectionrecap}
        Inference and deployment are separate outputs. Data and uncertainty analysis determine
        an evidence-preserving class set; metadata, a declared class prior, and a deterministic
        lexicographic key may then select one operational class and one representative. Any
        choice outside the evidence-preserving set is a disclosed policy override.

        \noindent\textbf{Key equations.}
        Within class $C$, $h_{\mathrm{can}}(C)=\lexargmin_{h\in C}K_p(h;z)$. The inferential
        contender set is $\mathcal T_{\mathrm{ev}}=\mathcal T_{\mathrm{stat}}$ when calibrated
        uncertainty is available and $\mathcal T_{\mathrm{num}}$ otherwise. The default output is
        \(C_{\mathrm{op}}=\lexargmin_{C\in\mathcal P_{p,z}}\overline K_p(C;z)\) and
        $h_{\mathrm{op}}=h_{\mathrm{can}}(C_{\mathrm{op}})$.

        \noindent\textbf{Symbols introduced.}
        \par\noindent\begin{tabularx}{\linewidth}{@{}p{0.25\linewidth}Y@{}}
\toprule
Symbol & Meaning \\
\midrule
$K_p(h;z)$ & ordered canonical key under source profile $p$ and metadata $z$ \\
$\pi_{\mathrm{class}}(C\mid z)$ & prior probability assigned to a quotient class \\
$\varepsilon_{\mathrm{num}},\varepsilon_{\mathrm{pr}}$ & numerical and application-scale tie thresholds \\
$\mathcal T_{\mathrm{num}},\mathcal T_{\mathrm{stat}},\mathcal T_{\mathrm{pr}}$ & numerical, statistical, and practical contender sets \\
$\mathcal T_{\mathrm{ev}}$ & evidence-preserving set used before policy selection \\
$C_{\mathrm{op}},h_{\mathrm{op}}$ & selected operational class and canonical representative \\
\bottomrule
        \end{tabularx}
    \end{sectionrecap}


    \section{Computational acceleration and reproducible benchmarking}\label{app:acceleration}

    A broad convention search can be expensive because discrete decoders, integration
    substeps, window lengths, and continuous timing parameters multiply. Speedups should be
    divided into \emph{exact} transformations that leave the final score unchanged and
    \emph{screening} transformations that may discard candidates and therefore require a
    recall benchmark.

    Let $H=\abs{\mathcal H}$ be the enumerated hypothesis count, $Q=\abs{\mathcal Q}$ the
    number of structural quotient classes, $N$ the interval count, $M$ the number of
    integration substeps, and $\mathcal L$ the window set. Straightforward integration costs
    $O(HMN)$. If every length-$\ell$ window is recomposed from its $\ell$ interval
    increments, the additional cost is
    \begin{equation}
        O\!\left(HN\sum_{\ell\in\mathcal L}\ell\right).
    \end{equation}
    The first goal is to replace $H$ by $Q$ and the window sum by $\abs{\mathcal L}$.

    \subsection{Exact accelerations}

    \paragraph{Quotient before scoring.}
    Construct canonical representatives of the exact classes in
    Table~\ref{tab:exact-equivalences}. Do not score both members of a
    direction--frame--sign dual, both intrinsic/extrinsic reversal duplicates, or multiple
    Euler branches that decode to the same $R_k$. Derive the omitted member in the report.
    This reduces work without changing any inferential output.

    \paragraph{Factor and cache decoders.}
    Many hypotheses share the same attitude decoder, observed increment sequence, rate
    channel map, unit conversion, or timing profile. Cache
    \begin{equation}
        R_{d,k},\quad
        \Delta R_{d,f,k}^{\mathrm{obs}},\quad
        C_c y_k,\quad
        \phi_{d,f,k}=\Log\Delta R_{d,f,k}^{\mathrm{obs}}
    \end{equation}
    by their actual dependencies rather than recomputing them for every long tuple. Hard
    validity checks---quaternion norm, matrix orthogonality, impossible schema domains, and
    nonmonotone timestamps---should run before group integration.

    \paragraph{Use prefix products for windows.}
    For body increments, define $S_0=I$ and $S_{k+1}=S_kP_k$; then
    \begin{equation}
        P_{B,k}^{(\ell)}=S_k^{-1}S_{k+\ell}.
    \end{equation}
    For world increments, define $S_{k+1}=P_kS_k$; then
    \begin{equation}
        P_{W,k}^{(\ell)}=S_{k+\ell}S_k^{-1}.
    \end{equation}
    After one $O(N)$ prefix pass, every configured window product costs $O(1)$, reducing
    window construction to $O(QN\abs{\mathcal L})$.

    \paragraph{Exploit exact lower bounds.}
    The triangle inequality gives the per-window bound
    \begin{equation}
        d_g(P,O)\ge\abs{d_g(I,P)-d_g(I,O)}.
        \label{eq:angle-only-lower-bound}
    \end{equation}
    It can reject clearly poor candidates before computing a full relative rotation axis.
    For the upper-trimmed mean, a second safe bound is available.

    \begin{proposition}[Partial lower bound for an upper-trimmed mean]
Let $z_1,\ldots,z_n\ge0$, retain the $m$ smallest values, and suppose $r$ values have
been evaluated. Let $z_{(1:r)}\le\cdots\le z_{(r:r)}$ order the evaluated values and set
\begin{equation}
    q_r=\max\set{0,m-(n-r)}.
\end{equation}
Then every completion of the unevaluated values satisfies
\begin{equation}
    T_{\alpha_{\rm trim}}(z_1,\ldots,z_n)
    \ge
    \frac{1}{m}\sum_{i=1}^{q_r}z_{(i:r)}.
    \label{eq:trimmed-mean-lower-bound}
\end{equation}
    \end{proposition}
    \begin{proof}
        At most $n-r$ retained values can come from the unseen set, so at least $q_r$ evaluated
        values must appear among the final $m$ retained values. Their smallest possible sum is
        the sum of the $q_r$ smallest evaluated values. Unseen values are nonnegative.
    \end{proof}
    Weighted per-scale versions of \eqref{eq:trimmed-mean-lower-bound} can be accumulated.
    Once a candidate's lower bound exceeds the incumbent verified score, evaluation can stop
    without changing the winner.

    \paragraph{Optimize identifiable nuisance coordinates.}
    Search $\lambda=s_t\kappa$ or $\lambda=\Delta\kappa$, not gauge-redundant pairs. Use the
    first-order period fit and a lag cross-correlation only to seed a bounded global search;
    cache objective evaluations shared by candidates with the same decoded rate direction.

    \subsection{Coarse-to-fine screening}

    Approximate screening is often the largest speedup but must preserve the true quotient
    class with high probability. A useful sequence is:
    \begin{enumerate}[nosep]
    \item hard schema and numerical validity filters;
    \item a small subset $\mathcal S$ of intervals away from the cut locus, scored by the
    first-order tangent proxy
    \begin{equation}
        \widetilde J_C
        =\median_{k\in\mathcal S}
        \norm{\phi_{C,k}-\Delta t_k\bar\omega_{C,k}}^2;
        \label{eq:coarse-tangent-score}
    \end{equation}
    \item one-step or low-$M$ group scoring of a wider contender set;
    \item full-$M$, multiscale, robust scoring of the top $K$ classes plus every class
    inside a configured screening margin; and
    \item held-out validation and statistical ambiguity analysis only for finalists.
    \end{enumerate}
    The proxy in \eqref{eq:coarse-tangent-score} is unreliable for large increments,
    noncommuting within-interval motion, or logarithms near $\pi$. It should never be the
    final score. The benchmark acceptance metric is top-$K$ \emph{quotient-class recall}
    across the deployment envelope, not average runtime on easy trajectories.

    Candidates are embarrassingly parallel, but vectorization is not automatically faster.
    Large candidate-by-sample temporary arrays can become memory-bandwidth bound. Batch size,
    quaternion versus matrix kernels, and CPU versus GPU execution should be measured rather
    than assumed.

    \subsection{Benchmark design}

    A reproducible benchmark should cross the following factors:
    \begin{itemize}[nosep]
    \item availability regime: $\mathsf A+\mathsf G$ with known time, unknown scale,
    unknown latency, order only, $\mathsf A$ only, and $\mathsf G$ only;
    \item trajectory: static, coaxial, planar, full-rank noncommuting, periodic,
    near-Euler-singular, and aliasing-prone;
    \item size: $N$, $H$, $Q$, $M$, and $\mathcal L$;
    \item corruption: attitude noise, gyro noise, bias, scale error, outliers, dropped
    samples, and timestamp jitter; and
    \item implementation: exhaustive, quotient-aware, cached/prefix, batched,
    coarse-to-fine, and early-abandon variants.
    \end{itemize}
    Report wall time, candidate-intervals per second, peak memory, exact class recovery,
    top-$K$ recall after each screen, score deviation from exhaustive evaluation, fitted
    nuisance error, and invariance-test failures. Record hardware, thread count, software
    versions, seeds, and whether compilation or cache warm-up is included.

    Table~\ref{tab:microbenchmark} is an illustrative kernel benchmark generated by the
    accompanying script. It uses synthetic noiseless full-rank body-rate motion, $384$
    enumerated candidates, $192$ exact dual classes, a one-step quaternion score, five
    median repeats, NumPy~2.3.5, SciPy~1.17.0, Python~3.13.5, in a container exposing five virtual AMD EPYC CPU cores. The coarse stage scores every sixteenth interval, retains $K=16$ classes, and in
    this clean test retains the true class at all sizes. These numbers demonstrate the cost
    structure; they are not a statistical recovery claim.

    \begin{table}[H]
        \centering
        \small
        \caption{Illustrative reference-kernel runtime. Speedup is relative to all $384$
            candidates.}\label{tab:microbenchmark}
        \begin{tabular}{@{}rrrrrr@{}}
            \toprule
            $N$    & Exhaustive (s) & Quotient (s) & Batched quotient (s) & Coarse-to-fine (s) & Coarse speedup \\
            \midrule
            1,000  & 0.130          & 0.066        & 0.054                & 0.025              & $5.15\times$   \\
            5,000  & 0.497          & 0.236        & 0.255                & 0.048              & $10.30\times$  \\
            20,000 & 1.908          & 0.962        & 0.888                & 0.152              & $12.58\times$  \\
            \bottomrule
        \end{tabular}
    \end{table}

    The benchmark should be extended with Monte Carlo noise and outliers before choosing
    $K$, stride, or pruning thresholds. A screen that is ten times faster but loses the true
    class in one difficult regime is not acceptable unless that failure rate is explicitly
    part of the product requirement.

    \begin{sectionrecap}
        Exact acceleration comes first: quotient structural duplicates, cache by dependency,
        construct windows with prefix products, optimize only identifiable nuisance coordinates,
        and stop candidates using proven lower bounds. Coarse-to-fine screening can add an order
        of magnitude in the reference kernel, but its governing metric is quotient-class recall
        under difficult trajectories, not speed alone.

        \noindent\textbf{Key equations.}
        Prefix products give $P_{B,k}^{(\ell)}=S_k^{-1}S_{k+\ell}$ and
        $P_{W,k}^{(\ell)}=S_{k+\ell}S_k^{-1}$. Safe pruning uses
        $d_g(P,O)\ge|d_g(I,P)-d_g(I,O)|$ and the partial trimmed-mean bound in
        \eqref{eq:trimmed-mean-lower-bound}. Approximate screening uses
        $\widetilde J_C$ only to choose finalists.

        \noindent\textbf{Symbols introduced.}
        \par\noindent\begin{tabularx}{\linewidth}{@{}p{0.23\linewidth}Y@{}}
\toprule
Symbol & Meaning \\
\midrule
$H,Q$ & enumerated hypothesis count and structural quotient-class count \\
$N,M,\mathcal L$ & intervals, integration substeps, and window-length set \\
$S_k$ & prefix product used for $O(1)$ window construction \\
$q_r$ & minimum number of evaluated values forced into a trimmed mean \\
$\widetilde J_C$ & first-order class screening score, never the final risk \\
$K$ & number of classes retained for full evaluation \\
\bottomrule
        \end{tabularx}
    \end{sectionrecap}


    \section{Reporting schema and implementation invariants}\label{app:report}

    A machine-readable result should preserve the same hierarchy used in the mathematics.
    At minimum, report:
    \begin{enumerate}[nosep]
    \item the data-availability label from Table~\ref{tab:availability-modes};
    \item the best structural quotient class or classes, listing all exact members;
    \item the evidence-preserving contender set and the criterion used to form it;
    \item the operational representative, if one is required;
    \item a per-field provenance label distinguishing hard metadata, trajectory evidence,
    heuristic classification, and canonical policy; and
    \item diagnostics for rejected near-ties, timing modes, wrap/unit ambiguity, and
    excitation/alias warnings; and
    \item when cohort fusion is used, the immutable local packet, quotient/calibration
    versions, reference packet IDs, relationship model, lineage audit, prior and
    posterior sharedness, and local/borrowed/fused class distributions.
    \end{enumerate}


    A defensible report should contain at least the following:
    \begin{itemize}[nosep]
    \item the winning structural equivalence class, not only one representative;
    \item the class-level robust risk $J_C$, untrimmed risk, per-scale risks, and class gap;
    \item exact dual partners under direction--frame--sign inversion;
    \item any rate-scale/sample-period gauge partners;
    \item fitted nuisance parameters, active bounds, and alternate local minima;
    \item total motion, Gramian eigenvalues, and effective rank;
    \item per-interval alias bounds and fraction near the $\pi$ cut locus;
    \item Euler singularity diagnostics where applicable;
    \item residual covariance eigenvalues and temporal correlation diagnostics;
    \item integration-convergence comparison between $M$ and $2M$;
    \item a held-out or block-bootstrap stability result;
    \item hard constraints and priors shown separately from trajectory evidence;
    \item the selected source profile, provenance cues, and whether each cue was hard or soft;
    \item the availability regime and which claims are trajectory-confirmed versus attitude-only or rate-only heuristics;
    \item the clock model, identifiable scale products, latency modes, and any time-warp regularization;
    \item Euler principal-range, branch-complexity, and singularity diagnostics where applicable;
    \item the structural-equivalence catalog entry used to group every exact pair;
    \item benchmarked top-$K$ recall, runtime, and peak memory for every enabled screening stage;
    \item numerical, statistical, and practical tie sets with the rule used for each;
    \item the canonical key $K_p$, canonical representative, and selection-basis label;
    \item confirmation that the result is invariant to hypothesis enumeration order; and
    \item the downstream sensitivity to every class retained in the reported contender set;
    \item the prior-separated class evidence $\lambda_i$, local prior $\pi_i^{(0)}$, and
    calibration status for every exported evidence packet;
    \item atomic measurement/derivation lineage and the result of every dependency check;
    \item the cohort model (M1, M2, or M3), reference packet set, $\rho_0$ sensitivity,
    posterior sharedness, and no-borrowing replay; and
    \item the pairwise local-versus-borrowed log-odds decomposition for the selected and
    runner-up classes.
    \end{itemize}

    The following identities should be enforced as property-based tests:
    \begin{align}
        d_g(A,B)&=d_g(QA,QB)=d_g(AQ,BQ),\\
        d_g(A,B)&=d_g(A^{-1},B^{-1}),\\
        \Exp(Qv)&=Q\Exp(v)Q^{\T} \quad (Q\in\SO),\\
        P_W(-\omega)&=P_B(\omega)^{-1},\\
        J_h^{\rm test}&=J_{h^\dagger}^{\rm test},\\
        J_C(\kappa,\Delta)&=J_C(c\kappa,\Delta/c),\\
        h_{\mathrm{can}}(C;\text{order}_1)&=h_{\mathrm{can}}(C;\text{order}_2),\\
        q\!\left(s_{p,z}(C)\right)&=C,\\
        \operatorname{deprior}\!\left(q_i,\pi_i^{(0)}\right)&\equiv\lambda_i+\text{constant},\\
        q_T^{\mathrm{fused}}(\rho_0=0)&=q_T^{\mathrm{baseline}},\\
        \mathcal F_T(\text{reference order}_1)&=\mathcal F_T(\text{reference order}_2).
    \end{align}
    The final cohort tests verify prior removal, the no-borrowing limit, and invariance to
    reference enumeration. Strict-mode tests should additionally reject overlapping atomic
    lineage and should reproduce each member's result using a leave-one-out cohort prior. A
    failure indicates prior double-counting, hidden packet dependence, or a fusion result that
    depends on container order rather than evidence.

    \begin{sectionrecap}
        A report must expose the quotient class, every retained tie set, fitted nuisance modes,
        excitation and alias diagnostics, the policy used to choose a representative, and the
        runtime/recall behavior of screening. Property tests should enforce analytical
        symmetries independently of candidate enumeration order.

        \noindent\textbf{Key invariants.}
        The implementation should verify, to numerical precision,
        \begin{align*}
            J_h^{\rm test}&=J_{h^\dagger}^{\rm test},\\
            J_C(\kappa,\Delta)&=J_C(c\kappa,\Delta/c),\\
            h_{\mathrm{can}}(C;\mathrm{order}_1)
            &=h_{\mathrm{can}}(C;\mathrm{order}_2).
        \end{align*}

        \noindent\textbf{Reported quantities.}
        \par\noindent\begin{tabularx}{\linewidth}{@{}p{0.28\linewidth}Y@{}}
\toprule
Quantity & Interpretation \\
\midrule
$J_C$ and per-scale risks & class-level geometric evidence, both trimmed and untrimmed \\
$\delta_C^{\mathrm{sym}}$ & implementation spread among analytically equivalent members \\
$\widehat{\mathcal C}_{\min},\widehat\gamma$ and tie sets & empirical minima, next-class separation, and unresolved alternatives \\
$\widehat\lambda,\widehat\delta_t$ & identifiable time/rate scale product and latency modes \\
$K_p,C_{\mathrm{op}},h_{\mathrm{op}}$ & declared policy key, operational class, and representative \\
$\lambda_i,\mathsf{lin}_i,\mathsf{cal}_i$ & prior-separated packet evidence, atomic lineage, and calibration identity \\
$\rho_0,\rho^{\mathrm{post}},q_T^{\mathrm{fused}}$ & declared prior sharedness, posterior compatibility, and fused target result \\
Top-$K$ recall and peak memory & screening reliability and computational cost \\
\bottomrule
        \end{tabularx}
    \end{sectionrecap}


    \section{Stress evaluation}\label{app:stress}

    Synthetic evaluation should separate implementation correctness from statistical
    performance. Deterministic tests should include:
    \begin{itemize}[nosep]
    \item exact direction--frame--sign duality to machine precision;
    \item exact rate-scale/sample-period invariance when timestamps are absent;
    \item quaternion $\OO(4)$ distance invariance in orientation-only mode;
    \item pure-axis motion producing the predicted body/world ambiguity;
    \item noncommuting time-varying rates verifying second-order midpoint convergence;
    \item intervals near and beyond the $\pi$ alias boundary;
    \item proper-Euler and Tait--Bryan singular neighborhoods;
    \item random quaternion sign flips, which must not affect decoded rotations;
    \item independent attitude noise producing the predicted lag-one residual structure;
    \item outlier fractions below, at, and above the configured trim fraction;
    \item candidate-dependent period fits with multiple alias-related minima;
    \item randomized hypothesis enumeration producing the same canonical representative;
    \item exact ties under ROS, MAVLink, SciPy, and generic profiles selecting the documented
    profile representative without changing the quotient-class score;
    \item practical-band tests showing that policy selection never changes the reported
    trajectory ordering; and
    \item deliberate profile contradictions verifying that hard metadata rejects rather
    than merely down-ranks a candidate;
    \item attitude-only data showing exact quaternion-isometry ties despite known timestamps;
    \item rate-only data verifying that every decoder self-integrates and remains unranked without an external target;
    \item intrinsic/extrinsic Euler reversal pairs and branch-equivalent triples matching exactly;
    \item principal-output versus unwrapped Euler logs exercising the range and branch heuristics;
    \item affine clock/rate-scale transformations preserving predictions to machine precision;
    \item constant, aperiodic, and periodic rate histories for latency observability and multimodality;
    \item quotient, prefix-product, early-abandon, and coarse-to-fine kernels matching exhaustive finalist scores; and
    \item screening Monte Carlo studies reporting top-$K$ class recall rather than only speed;
    \item A/B-strong, C-weak cohort studies under correct sharing, partial exchangeability,
    and deliberate target--cohort conflict;
    \item duplicated, overlapping, and derived packets verifying that lineage checks prevent
    confidence inflation;
    \item calibration-shift studies across domains, excitation regimes, and packet lengths;
    and
    \item leave-one-out cohort reports verifying that no target packet contributes to its
    own borrowed prior.
    \end{itemize}

    Monte Carlo studies should report recovery of equivalence classes, not arbitrary
    representatives. Useful curves include class-selection frequency versus trajectory
    length, risk gap, excitation spectrum, noise level, outlier fraction, sample rate, and
    latency. Calibration claims require held-out simulations generated from assumptions at
    least as broad as those expected in deployment.

    \begin{sectionrecap}
        Deterministic tests establish implementation invariants; Monte Carlo tests establish
        recovery behavior under a stated deployment envelope. Both must score equivalence
        classes rather than arbitrary representatives. Screening is acceptable only when its
        class recall and score agreement are measured across weak excitation, aliasing,
        latency, noise, and outlier regimes.
    \end{sectionrecap}


    \section{Limitations and extensions}\label{app:limits}

    The base model assumes one constant convention, synchronized clocks, knot-sampled rates,
    linear intersample rate behavior, no gyro bias, no scale-factor matrix beyond discrete
    unit choices, no axis misalignment beyond signed permutations, and no time-varying
    calibration. Bias, full calibration matrices, latency, clock drift, and piecewise
    convention changes are natural extensions, but each introduces continuous parameters
    and new gauge freedoms.

    If a full $3\times3$ rate calibration matrix is introduced without constraints, it can
    absorb permutations, signs, scales, and some frame errors; identifiability then requires
    regularization or calibration references. Likewise, an unconstrained time warp can fit
    rate magnitude errors. Model flexibility should be added only with explicit anchors and
    held-out validation.

    Attitude-only and rate-only modes require especially careful labels. Attitude-only smoothness and Euler-range rules are priors or source-profile heuristics; rate-only self-integration is not a convention test. Likewise, benchmark speedups obtained on clean synthetic data do not establish screening recall under noise, aliasing, latency, or weak excitation.

    Trajectory consistency also cannot establish semantic labels that never enter the
    kinematics. For example, it cannot decide which physical axis an otherwise self-consistent
    coordinate triad should be called ``north'' without an external world reference.

    \begin{sectionrecap}
        Greater model flexibility introduces new gauges as well as new fit capability. Full rate
        calibration matrices and unconstrained time warps can absorb the very convention errors
        being tested, so they require anchors, regularization, and held-out validation.
        Attitude-only and rate-only outputs must be labeled as heuristic or externally anchored,
        not as ordinary trajectory-confirmed recovery.
    \end{sectionrecap}

    \begin{thebibliography}{99}

        \bibitem{shuster1993}
        M. D. Shuster,
        ``A Survey of Attitude Representations,''
        \emph{The Journal of the Astronautical Sciences},
        vol. 41, no. 4, pp. 439--517, 1993.

        \bibitem{diebel2006}
        J. Diebel,
        ``Representing Attitude: Euler Angles, Unit Quaternions, and Rotation Vectors,''
        Stanford University, Technical Report, 2006.

        \bibitem{huynh2009}
        D. Q. Huynh,
        ``Metrics for 3D Rotations: Comparison and Analysis,''
        \emph{Journal of Mathematical Imaging and Vision},
        vol. 35, no. 2, pp. 155--164, 2009.

        \bibitem{iserles2000}
        A. Iserles, H. Z. Munthe-Kaas, S. P. N\o rsett, and A. Zanna,
        ``Lie-Group Methods,''
        \emph{Acta Numerica}, vol. 9, pp. 215--365, 2000.

        \bibitem{blanes2009}
        S. Blanes, F. Casas, J. A. Oteo, and J. Ros,
        ``The Magnus Expansion and Some of Its Applications,''
        \emph{Physics Reports}, vol. 470, nos. 5--6, pp. 151--238, 2009.

        \bibitem{barfoot2017}
        T. D. Barfoot,
        \emph{State Estimation for Robotics}.
        Cambridge University Press, 2017.

        \bibitem{markley2014}
        F. L. Markley and J. L. Crassidis,
        \emph{Fundamentals of Spacecraft Attitude Determination and Control}.
        Springer, 2014.

        \bibitem{lee2018}
        T. Lee,
        ``Bayesian Attitude Estimation with the Matrix Fisher Distribution on $\SO$,''
        \emph{IEEE Transactions on Automatic Control},
        vol. 63, no. 10, pp. 3377--3392, 2018.

        \bibitem{hoeffding1963}
        W. Hoeffding,
        ``Probability Inequalities for Sums of Bounded Random Variables,''
        \emph{Journal of the American Statistical Association},
        vol. 58, no. 301, pp. 13--30, 1963.

        \bibitem{akaike1974}
        H. Akaike,
        ``A New Look at the Statistical Model Identification,''
        \emph{IEEE Transactions on Automatic Control},
        vol. 19, no. 6, pp. 716--723, 1974,
        doi: 10.1109/TAC.1974.1100705.

        \bibitem{schwarz1978}
        G. Schwarz,
        ``Estimating the Dimension of a Model,''
        \emph{The Annals of Statistics}, vol. 6, no. 2, pp. 461--464, 1978,
        doi: 10.1214/aos/1176344136.

        \bibitem{stone1974}
        M. Stone,
        ``Cross-Validatory Choice and Assessment of Statistical Predictions,''
        \emph{Journal of the Royal Statistical Society: Series B},
        vol. 36, no. 2, pp. 111--133, 1974,
        doi: 10.1111/j.2517-6161.1974.tb00994.x.

        \bibitem{rissanen1978}
        J. Rissanen,
        ``Modeling by Shortest Data Description,''
        \emph{Automatica}, vol. 14, no. 5, pp. 465--471, 1978,
        doi: 10.1016/0005-1098(78)90005-5.

        \bibitem{kass1995}
        R. E. Kass and A. E. Raftery,
        ``Bayes Factors,''
        \emph{Journal of the American Statistical Association},
        vol. 90, no. 430, pp. 773--795, 1995,
        doi: 10.1080/01621459.1995.10476572.

        \bibitem{bissiri2016}
        P. G. Bissiri, C. C. Holmes, and S. G. Walker,
        ``A General Framework for Updating Belief Distributions,''
        \emph{Journal of the Royal Statistical Society: Series B},
        vol. 78, no. 5, pp. 1103--1130, 2016,
        doi: 10.1111/rssb.12158.

        \bibitem{hansen2011}
        P. R. Hansen, A. Lunde, and J. M. Nason,
        ``The Model Confidence Set,''
        \emph{Econometrica}, vol. 79, no. 2, pp. 453--497, 2011,
        doi: 10.3982/ECTA5771.

        \bibitem{triggs2000}
        B. Triggs, P. F. McLauchlan, R. I. Hartley, and A. W. Fitzgibbon,
        ``Bundle Adjustment---A Modern Synthesis,'' in \emph{Vision Algorithms: Theory and
        Practice}, Lecture Notes in Computer Science, vol. 1883, pp. 298--372, 2000,
        doi: 10.1007/3-540-44480-7\_21.

        \bibitem{rosrep103}
        T. Foote and M. Purvis,
        ``REP 103: Standard Units of Measure and Coordinate Conventions,''
        Open Robotics, active informational standard, accessed August 2026.

        \bibitem{rosrep145}
        P. Bovbel,
        ``REP 145: Conventions for IMU Sensor Drivers,''
        Open Robotics, draft informational standard, accessed August 2026.

        \bibitem{rosQuaternion}
        Open Robotics,
        ``geometry\_msgs/Quaternion Message Definition,''
        ROS documentation, accessed August 2026.

        \bibitem{mavlinkCommon}
        MAVLink Project,
        ``MAVLink Common Message Set,''
        protocol documentation, accessed August 2026.

        \bibitem{scipyQuaternion}
        SciPy Developers,
        ``scipy.spatial.transform.Rotation.from\_quat,''
        SciPy reference documentation, accessed August 2026.

        \bibitem{eigenQuaternion}
        Eigen Developers,
        ``Eigen::Quaternion Class Template Reference,''
        Eigen 3 documentation, accessed August 2026.


        \bibitem{shusterMarkley2006}
        M. D. Shuster and F. L. Markley,
        ``General Formula for Extracting the Euler Angles,''
        \emph{Journal of Guidance, Control, and Dynamics},
        vol. 29, no. 1, pp. 215--221, 2006,
        doi: 10.2514/1.16622.

        \bibitem{bernardes2022}
        E. Bernardes and S. Viollet,
        ``Quaternion to Euler Angles Conversion: A Direct, General and Computationally
        Efficient Method,''
        \emph{PLOS ONE}, vol. 17, no. 11, e0276302, 2022,
        doi: 10.1371/journal.pone.0276302.

        \bibitem{scipyEulerRanges}
        SciPy Developers,
        ``\texttt{scipy.spatial.transform.Rotation.as\_euler},''
        SciPy Reference Manual, version 1.17, 2026.

        \bibitem{liMourikis2014}
        M. Li and A. I. Mourikis,
        ``Online Temporal Calibration for Camera--IMU Systems: Theory and Algorithms,''
        \emph{The International Journal of Robotics Research},
        vol. 33, no. 7, pp. 947--964, 2014,
        doi: 10.1177/0278364913515286.

        \bibitem{bipmSI}
        Bureau International des Poids et Mesures,
        \emph{The International System of Units (SI)},
        9th ed., version 4.01, revised June 2026,
        doi: 10.59161/AUEZ1291.


        \bibitem{matlabEul2Quat}
        The MathWorks,
        ``\texttt{eul2quat}: Convert Euler Angles to Quaternion,''
        MATLAB Navigation/Robotics documentation, accessed August 2026.

        \bibitem{matlabAeroAngle2Quat}
        The MathWorks,
        ``\texttt{angle2quat}: Convert Rotation Angles to Quaternion,''
        Aerospace Toolbox documentation, accessed August 2026.

        \bibitem{naifQuaternion}
        NASA/JPL Navigation and Ancillary Information Facility,
        ``Quaternion to Matrix / Quaternion Styles,''
        SPICE Toolkit documentation, accessed August 2026.

        \bibitem{naifEuler}
        NASA/JPL Navigation and Ancillary Information Facility,
        ``\texttt{eul2m\_c}: Euler Angles to Matrix,''
        CSPICE documentation, accessed August 2026.

        \bibitem{naifCK}
        NASA/JPL Navigation and Ancillary Information Facility,
        ``C-Kernel Required Reading: Angular Velocity Vectors,''
        SPICE Toolkit documentation, accessed August 2026.

        \bibitem{nasaPDSQuaternion}
        NASA Planetary Data System,
        ``Quaternion \texttt{type} Attribute,''
        PDS Data Dictionary, accessed August 2026.

        \bibitem{px4VehicleOdometry}
        PX4 Project,
        ``\texttt{VehicleOdometry} Message Definition,''
        \texttt{px4\_msgs} documentation, accessed August 2026.

        \bibitem{unityQuaternion}
        Unity Technologies,
        ``\texttt{Quaternion} Scripting API,''
        Unity documentation, accessed August 2026.

        \bibitem{unityEuler}
        Unity Technologies,
        ``Rotation and Orientation in Unity,''
        Unity Manual, accessed August 2026.

        \bibitem{unityAngularVelocity}
        Unity Technologies,
        ``\texttt{Rigidbody.angularVelocity},''
        Unity Scripting API, accessed August 2026.

        \bibitem{unrealQuaternion}
        Epic Games,
        ``\texttt{TQuat} API Reference,''
        Unreal Engine documentation, accessed August 2026.

        \bibitem{unrealRotator}
        Epic Games,
        ``\texttt{TRotator} API Reference,''
        Unreal Engine documentation, accessed August 2026.

        \bibitem{unrealAngularVelocity}
        Epic Games,
        ``\texttt{FBodyInstance::GetUnrealWorldAngularVelocityInRadians},''
        Unreal Engine documentation, accessed August 2026.

        \bibitem{unrealAngularDegrees}
        Epic Games,
        ``\texttt{UPrimitiveComponent::GetPhysicsAngularVelocityInDegrees},''
        Unreal Engine documentation, accessed August 2026.

        \bibitem{godotQuaternion}
        Godot Engine Project,
        ``\texttt{Quaternion} Class Reference,''
        Godot Engine documentation, accessed August 2026.

        \bibitem{godotBasis}
        Godot Engine Project,
        ``\texttt{Basis} Class Reference,''
        Godot Engine documentation, accessed August 2026.

        \bibitem{godotAngularVelocity}
        Godot Engine Project,
        ``\texttt{RigidBody3D.angular\_velocity},''
        Godot Engine documentation, accessed August 2026.

        \bibitem{numpyUnwrap}
        NumPy Developers,
        ``\texttt{numpy.unwrap},''
        NumPy Reference Manual, version 2.5, 2026.



        \bibitem{ibrahim2000}
        J. G. Ibrahim and M.-H. Chen,
        ``Power Prior Distributions for Regression Models,''
        \emph{Statistical Science}, vol. 15, no. 1, pp. 46--60, 2000,
        doi: 10.1214/ss/1009212673.

        \bibitem{hobbs2011}
        B. P. Hobbs, B. P. Carlin, S. J. Mandrekar, and D. J. Sargent,
        ``Hierarchical Commensurate and Power Prior Models for Adaptive Incorporation of
        Historical Information in Clinical Trials,''
        \emph{Biometrics}, vol. 67, no. 3, pp. 1047--1056, 2011,
        doi: 10.1111/j.1541-0420.2011.01564.x.

        \bibitem{schmidli2014}
        H. Schmidli, S. Gsteiger, S. Roychoudhury, A. O'Hagan, D. Spiegelhalter, and
        B. Neuenschwander,
        ``Robust Meta-Analytic-Predictive Priors in Clinical Trials with Historical Control
        Information,''
        \emph{Biometrics}, vol. 70, no. 4, pp. 1023--1032, 2014,
        doi: 10.1111/biom.12242.

        \bibitem{yang2023}
        P. Yang, Y. Zhao, L. Nie, J. Vallejo, and Y. Yuan,
        ``SAM: Self-Adapting Mixture Prior to Dynamically Borrow Information from Historical
        Data in Clinical Trials,''
        \emph{Biometrics}, vol. 79, no. 4, pp. 2857--2868, 2023,
        doi: 10.1111/biom.13927.

        \bibitem{gneiting2007}
        T. Gneiting and A. E. Raftery,
        ``Strictly Proper Scoring Rules, Prediction, and Estimation,''
        \emph{Journal of the American Statistical Association}, vol. 102, no. 477,
        pp. 359--378, 2007,
        doi: 10.1198/016214506000001437.

    \end{thebibliography}

\end{document}
